% ============================================================================
% Bitwise Digits: Theory and Design of a Hybrid Logic-Based
% Neural Network for MNIST
% ============================================================================
\documentclass[11pt,a4paper,twoside,openright]{book}

% --- Shared preamble ---
% ============================================================================
% Shared preamble for the theory book
% ============================================================================

% --- Core packages ---
\usepackage{fontspec}
\usepackage[english]{babel}
\usepackage[top=2.5cm, bottom=2.5cm, inner=3cm, outer=2.5cm]{geometry}

% --- Typography ---
\usepackage{microtype}
\setmainfont{Latin Modern Roman}
\setsansfont{Latin Modern Sans}
\setmonofont{Latin Modern Mono}

% --- Math ---
\usepackage{amsmath, amssymb, amsthm}
\usepackage{mathtools}
\usepackage{bm}

% --- Graphics and diagrams ---
\usepackage{tikz}
\usetikzlibrary{
  arrows.meta, positioning, calc, decorations.pathreplacing,
  shapes.geometric, shapes.misc, fit, backgrounds, patterns,
  matrix, chains, circuits.logic.US, automata
}
\usepackage{pgfplots}
\pgfplotsset{compat=1.18}

% --- Tables and lists ---
\usepackage{booktabs}
\usepackage{tabularx}
\usepackage{enumitem}

% --- Colors ---
\usepackage{xcolor}
\definecolor{bitblue}{HTML}{2563EB}
\definecolor{bitgreen}{HTML}{16A34A}
\definecolor{bitorange}{HTML}{EA580C}
\definecolor{bitpurple}{HTML}{9333EA}
\definecolor{bitgray}{HTML}{6B7280}
\definecolor{lightbg}{HTML}{F8FAFC}
\definecolor{darkbg}{HTML}{1E293B}

% --- Code and algorithms ---
\usepackage{algorithm}
\usepackage{algpseudocode}

% --- Headers and footers ---
\usepackage{fancyhdr}
\pagestyle{fancy}
\fancyhf{}
\fancyhead[LE]{\small\itshape\leftmark}
\fancyhead[RO]{\small\itshape\rightmark}
\fancyfoot[LE,RO]{\thepage}
\renewcommand{\headrulewidth}{0.4pt}

% --- Hyperlinks and references ---
\usepackage[
  colorlinks=true,
  linkcolor=bitblue,
  citecolor=bitgreen,
  urlcolor=bitorange
]{hyperref}
\usepackage[capitalise,noabbrev]{cleveref}

% --- Bibliography ---
\usepackage[
  backend=biber,
  style=numeric-comp,
  sorting=nyt,
  maxbibnames=99
]{biblatex}
\addbibresource{references.bib}

% --- Theorem environments ---
\theoremstyle{definition}
\newtheorem{definition}{Definition}[chapter]
\newtheorem{example}{Example}[chapter]
\theoremstyle{remark}
\newtheorem{remark}{Remark}[chapter]
\theoremstyle{plain}
\newtheorem{theorem}{Theorem}[chapter]
\newtheorem{lemma}[theorem]{Lemma}

% --- Custom commands ---
\newcommand{\bitop}[1]{\texttt{\textcolor{bitblue}{#1}}}
\newcommand{\concept}[1]{\textbf{\textcolor{bitpurple}{#1}}}
\newcommand{\bvec}[1]{\bm{#1}}  % bold vector
\newcommand{\bmat}[1]{\bm{#1}}  % bold matrix
\newcommand{\popcount}{\operatorname{popcount}}
\newcommand{\sign}{\operatorname{sign}}
\newcommand{\TruthTable}{\operatorname{TT}}

% --- Custom environments ---
\usepackage{tcolorbox}
\tcbuselibrary{skins, breakable}

\newtcolorbox{keyinsight}[1][]{
  enhanced, breakable,
  colback=bitblue!5, colframe=bitblue!70,
  fonttitle=\bfseries,
  title={Key Insight},
  #1
}

\newtcolorbox{engineernote}[1][]{
  enhanced, breakable,
  colback=bitgreen!5, colframe=bitgreen!70,
  fonttitle=\bfseries,
  title={Engineer's Note},
  #1
}

\newtcolorbox{analogy}[1][]{
  enhanced, breakable,
  colback=bitorange!5, colframe=bitorange!70,
  fonttitle=\bfseries,
  title={Analogy},
  #1
}

\newtcolorbox{warning}[1][]{
  enhanced, breakable,
  colback=red!5, colframe=red!70,
  fonttitle=\bfseries,
  title={Warning},
  #1
}

% --- Chapter style ---
\usepackage{titlesec}
\titleformat{\chapter}[display]
  {\normalfont\huge\bfseries}
  {\chaptertitlename\ \thechapter}{20pt}{\Huge}
\titlespacing*{\chapter}{0pt}{-30pt}{40pt}


% --- Document metadata ---
\title{%
  \vspace{-2cm}
  {\Huge\bfseries Bitwise Digits}\\[1em]
  {\Large Theory and Design of a Hybrid Logic-Based\\
   Neural Network for MNIST}\\[2em]
  {\large A Practical Guide for Engineers}
}
\author{Bernhard Gerlach}
\date{March 2026}

% ============================================================================
\begin{document}
% ============================================================================

\maketitle
\thispagestyle{empty}

% --- Front matter ---
\frontmatter

\chapter*{Preface}

This book describes the theory and design of a neural network that uses no
floating-point arithmetic at inference time. Instead of the
multiply-accumulate operations that dominate modern AI, our architecture
relies on Boolean logic, integer counting, and memory lookups---operations
that every CPU can perform cheaply and efficiently.

The architecture combines three paradigms:
\begin{itemize}
  \item \textbf{Tsetlin Machines} learn interpretable if--then rules using
    simple counting automata. They form the feature extraction layer.
  \item \textbf{Logic Gate Networks} learn which logic gates (AND, OR,
    XOR, \ldots) best combine those features. They form the composition
    layer.
  \item \textbf{Lookup Table compilation} converts any remaining learned
    functions into precomputed tables. It forms the deployment layer.
\end{itemize}

The target task is handwritten digit recognition on the MNIST dataset---a
deliberately simple benchmark chosen to validate the approach, not to chase
records. The implementation is in Rust.

\paragraph{Who this book is for.}
Engineers who are comfortable with university-level mathematics (linear
algebra, basic probability, logic) but who are not specialists in machine
learning or theoretical computer science. Every concept is introduced from
intuition and diagrams before it is formalised.

\paragraph{How to read this book.}
Part~I builds the foundations: Boolean algebra, supervised learning, and
data preparation. Parts~II--IV cover the three components of the hybrid
architecture in depth. Part~V assembles everything into a complete system
and analyses its properties. Each part is designed to be self-contained
enough that you can read the Tsetlin chapters (Part~II) and the Logic Gate
chapters (Part~III) in either order, provided you read both before Part~V.

\paragraph{Conventions.}
Throughout this book, coloured boxes highlight important ideas:
\begin{keyinsight}
Blue boxes contain key theoretical insights---the ideas that, once
understood, make everything else click.
\end{keyinsight}
\begin{engineernote}
Green boxes contain practical engineering observations---implementation
tips, performance implications, and hardware considerations.
\end{engineernote}
\begin{analogy}
Orange boxes contain analogies that map unfamiliar concepts to everyday
experience. They sacrifice precision for intuition.
\end{analogy}

\vfill
\noindent\textit{March 2026}

\tableofcontents

% --- Main matter ---
\mainmatter

% Part I: Foundations
\part{Foundations}
\label{part:foundations}

% ============================================================================
\chapter{The Problem and the Promise}
\label{ch:problem}
% ============================================================================

\section{What This Book Is About}

This book describes a single, concrete neural network architecture and
explains every piece of it from the ground up. By the end, you will
understand not just \emph{what} the architecture does, but \emph{why}
each design choice was made and \emph{how} to implement it.

The architecture is a hybrid of three components:
\begin{enumerate}
  \item A \textbf{Tsetlin Machine}~\cite{granmo2018tsetlin} that learns interpretable Boolean
    rules from raw pixel data.
  \item A \textbf{Logic Gate Network}~\cite{petersen2022difflogic} that learns which logic gates
    best combine those rules into a classification decision.
  \item A \textbf{Lookup Table compiler}~\cite{liu2024kan} that converts any remaining
    learned functions into precomputed memory tables.
\end{enumerate}

The task is handwritten digit recognition on the MNIST dataset. The
implementation language is Rust. The ambition is simple: demonstrate
that a neural network built entirely from Boolean operations and memory
lookups can recognise digits with competitive accuracy---and do so
faster and with less energy than a conventional floating-point model~\cite{horowitz2014energy}.


\section{The MNIST Digit Recognition Task}
\label{sec:mnist-task}

The MNIST dataset~\cite{lecun1998mnist} is the ``Hello World'' of
machine learning. It consists of 70,000 images of handwritten digits
(0 through 9), each a $28 \times 28$ pixel grayscale image.

\begin{figure}[htbp]
\centering
\begin{tikzpicture}[scale=0.15]
  % Grid
  \fill[bitblue!5] (0,0) rectangle (28,28);
  \draw[bitgray!20, very thin, step=1] (0,0) grid (28,28);

  % Stylised "7"
  \foreach \x/\y in {
    8/24, 9/24, 10/24, 11/24, 12/24, 13/24, 14/24, 15/24, 16/24, 17/24, 18/24, 19/24,
    18/23, 19/23,
    17/22, 18/22,
    16/21, 17/21,
    15/20, 16/20,
    14/19, 15/19,
    14/18, 15/18,
    13/17, 14/17,
    13/16, 14/16,
    12/15, 13/15,
    12/14, 13/14,
    11/13, 12/13,
    11/12, 12/12,
    11/11, 12/11
  }{
    \fill[bitblue!60] (\x,\y) rectangle ++(1,1);
  }

  % Annotations
  \draw[<->, thick, bitgray] (-1, 0) -- (-1, 28)
    node[midway, left, font=\small\sffamily] {28\,px};
  \draw[<->, thick, bitgray] (0, -1) -- (28, -1)
    node[midway, below, font=\small\sffamily] {28\,px};
\end{tikzpicture}
\hspace{2cm}
\begin{tikzpicture}[scale=1]
  % Flattened vector
  \node[font=\sffamily\bfseries] at (0, 3.5) {Flattened};
  \draw[thick] (0, 0) -- (0, 3);
  \foreach \y/\v in {2.8/0, 2.5/0, 2.2/0.72, 1.9/0.95, 1.6/0.88,
    1.3/\ldots, 1.0/0, 0.7/0.31, 0.4/0} {
    \node[font=\tiny\ttfamily, right] at (0.1, \y) {\v};
  }
  \node[font=\scriptsize\sffamily, bitgray] at (0.5, -0.3)
    {784 values};

  % Arrow
  \draw[-{Stealth}, thick] (1.5, 1.5) -- (3, 1.5);
  \node[font=\scriptsize\sffamily] at (2.25, 1.9) {classify};

  % Output
  \node[draw, rounded corners, fill=bitgreen!20, minimum width=1.5cm,
    minimum height=0.8cm, font=\sffamily\bfseries] at (4, 1.5) {7};
\end{tikzpicture}
\caption{The MNIST task. A $28 \times 28$ grayscale image is flattened into
  a vector of 784 values (each in $[0, 255]$ or normalised to $[0, 1]$),
  and the network must output the correct digit label.}
\label{fig:mnist-task}
\end{figure}

\begin{itemize}
  \item \textbf{Training set:} 60,000 labelled images used to teach the
    model.
  \item \textbf{Test set:} 10,000 labelled images, never seen during
    training, used to measure accuracy.
  \item \textbf{Baseline:} A standard two-layer floating-point neural
    network achieves approximately 98\% accuracy on the test set.
\end{itemize}

MNIST is deliberately simple. We choose it not to set records but to
have a well-understood benchmark where we can focus entirely on the
architecture rather than on data engineering or scale.


\section{Why Bitwise? The Energy and Hardware Argument}
\label{sec:why-bitwise}

Every neural network in production today is built on a single
operation: \concept{multiply-accumulate} (MAC). Multiply a number by
a weight, add to a running total. Repeat billions of times.

This operation is expensive. \Cref{tab:energy-cost} shows the energy
cost of different operations in a modern chip.

\begin{table}[htbp]
\centering
\caption{Energy cost per operation in 45\,nm technology~\cite{horowitz2014energy}.}
\label{tab:energy-cost}
\begin{tabular}{@{}lrr@{}}
  \toprule
  \textbf{Operation} & \textbf{Energy (pJ)} & \textbf{vs.\ Logic Gate} \\
  \midrule
  Bitwise AND/OR/XOR  & 0.1  & $1\times$ \\
  Integer add (32-bit) & 0.9  & $9\times$ \\
  FP multiply (32-bit) & 3.7  & $37\times$ \\
  SRAM read            & 5.0  & $50\times$ \\
  DRAM read            & 640  & $6{,}400\times$ \\
  \bottomrule
\end{tabular}
\end{table}

\begin{keyinsight}
A floating-point multiplication costs 37 times more energy than a
bitwise logic operation~\cite{horowitz2014energy}. But moving data from main memory costs
6,400 times more. Any architecture that keeps its model small enough
to fit in cache---and uses logic operations instead of
multiplication---wins on both fronts.
\end{keyinsight}

There is also a market structure argument. Training and running large
floating-point models requires GPUs, a market dominated by a single
manufacturer. Bitwise operations run on every CPU ever made. An
architecture that needs only AND, OR, XOR, integer addition, and table
lookups can run on a microcontroller, an FPGA, or a 20-year-old laptop.


\section{The Hybrid Thesis}
\label{sec:thesis}

Our thesis is:

\begin{quote}
\itshape
A three-stage pipeline---Tsetlin Machine feature extraction, Logic
Gate Network composition, and Lookup Table deployment---can classify
MNIST digits with competitive accuracy ($>$97\%) while using
\emph{zero} floating-point operations at inference time.
\end{quote}

\begin{figure}[htbp]
\centering
\begin{tikzpicture}[
  stage/.style={draw, rounded corners=5pt, minimum width=3cm,
    minimum height=1.6cm, align=center, font=\small\sffamily,
    line width=0.8pt},
  arr/.style={-{Stealth[length=5pt]}, thick},
]
  \node[stage, fill=lightbg] (in) at (0, 0)
    {\textbf{Input}\\$28 \times 28$ pixels};
  \node[stage, fill=bitgreen!10, draw=bitgreen] (tm) at (4.5, 0)
    {\textbf{Tsetlin Machine}\\Boolean clauses};
  \node[stage, fill=bitorange!10, draw=bitorange] (lgn) at (9, 0)
    {\textbf{Logic Gates}\\Learned circuit};
  \node[stage, fill=bitpurple!10, draw=bitpurple] (lut) at (13.5, 0)
    {\textbf{LUT Scoring}\\Table lookups};

  \draw[arr] (in) -- (tm) node[midway, above, font=\scriptsize\sffamily]
    {booleanise};
  \draw[arr] (tm) -- (lgn) node[midway, above, font=\scriptsize\sffamily]
    {clause votes};
  \draw[arr] (lgn) -- (lut) node[midway, above, font=\scriptsize\sffamily]
    {gate outputs};

  % Training phases
  \node[font=\scriptsize\sffamily, bitgreen] at (4.5, -1.4)
    {Phase 1: automata};
  \node[font=\scriptsize\sffamily, bitorange] at (9, -1.4)
    {Phase 2: backprop};
  \node[font=\scriptsize\sffamily, bitpurple] at (13.5, -1.4)
    {Phase 3: compile};
\end{tikzpicture}
\caption{The three-stage hybrid architecture. Each stage is trained
  separately, and each uses a different learning mechanism. At inference
  time, the entire pipeline runs on Boolean operations, integer addition,
  and table lookups.}
\label{fig:hybrid-pipeline}
\end{figure}

The training happens in three phases:
\begin{description}[leftmargin=1em]
  \item[Phase 1 (Tsetlin):]
    Train Boolean clauses on booleanised pixel data using Tsetlin
    automata~\cite{tsetlin1961}. No floating-point arithmetic---only integer increment/decrement
    and random number generation.
  \item[Phase 2 (Logic Gates):]
    Freeze the clauses. Train a Logic Gate Network~\cite{petersen2024difflogic} on the clause outputs
    using gradient descent with a softmax relaxation over gate types.
    This phase uses floating-point arithmetic during training only.
  \item[Phase 3 (LUT Compilation):]
    Compile any remaining continuous scoring functions into integer
    lookup tables. This is a one-time offline step.
\end{description}

After Phase~3, the inference model contains no floating-point operations.


\section{What the Reader Needs}
\label{sec:prerequisites}

This book assumes:
\begin{itemize}
  \item \textbf{Linear algebra:} Vectors, matrices, dot products. You
    should be comfortable with $\bvec{y} = \bmat{W}\bvec{x} + \bvec{b}$.
  \item \textbf{Basic probability:} Random variables, expected value,
    conditional probability. You do not need measure theory.
  \item \textbf{Logic:} AND, OR, NOT, truth tables. Chapter~2 reviews
    this in detail.
  \item \textbf{Programming intuition:} You should be comfortable
    reading pseudocode. Rust experience is helpful for Chapter~14 but
    not required.
\end{itemize}

You do \emph{not} need prior experience with machine learning,
neural networks, or AI. Every concept is introduced from scratch.


\section{Notation}
\label{sec:notation}

We use the following conventions throughout:

\begin{table}[htbp]
\centering
\caption{Notation conventions used in this book.}
\label{tab:notation}
\begin{tabular}{@{}ll@{}}
  \toprule
  \textbf{Symbol} & \textbf{Meaning} \\
  \midrule
  $x, y, z$ & Scalars (lowercase italic) \\
  $\bvec{x}, \bvec{w}$ & Vectors (bold lowercase) \\
  $\bmat{W}$ & Matrices (bold uppercase) \\
  $\{0, 1\}$ & Binary / Boolean values \\
  $\{-1, +1\}$ & Bipolar binary values \\
  $\bar{x}$ & Boolean complement (NOT $x$) \\
  $\wedge, \vee$ & Logical AND, OR \\
  $\oplus$ & Logical XOR \\
  $|S|$ & Cardinality of set $S$ \\
  $[n]$ & The set $\{1, 2, \ldots, n\}$ \\
  $\popcount(\cdot)$ & Number of 1-bits in a binary word \\
  \bottomrule
\end{tabular}
\end{table}

% ============================================================================
\chapter{Boolean Algebra and Digital Logic}
\label{ch:boolean}
% ============================================================================

This chapter reviews the mathematics of Boolean values and logic gates.
If you work with digital systems, much of this will be familiar---but
we establish notation and highlight the specific concepts that the later
chapters depend on.

% ----------------------------------------------------------------------------
\section{Bits and Boolean Values}
\label{sec:bits}
% ----------------------------------------------------------------------------

A \concept{bit} is the smallest unit of information: it is either 0
or~1. We write $x \in \{0, 1\}$ and call $x$ a \concept{Boolean
variable}.

A \concept{bit string} of length $n$ is an ordered sequence of $n$
bits:
\begin{equation}
  \bvec{x} = (x_1, x_2, \ldots, x_n) \in \{0, 1\}^n
  \label{eq:bitstring}
\end{equation}

For MNIST, each image will become a bit string of length 784 (one bit
per pixel) or 1568 (including negations). Understanding how to
manipulate bit strings efficiently is the foundation of everything that
follows.


% ----------------------------------------------------------------------------
\section{The Fundamental Logic Operations}
\label{sec:logic-ops}
% ----------------------------------------------------------------------------

Three operations form the basis of all Boolean computation:

\begin{definition}[AND, OR, NOT]
\label{def:basic-ops}
For $a, b \in \{0, 1\}$:
\begin{align}
  a \wedge b \; (\text{AND}) &=
    \begin{cases} 1 & \text{if } a = 1 \text{ and } b = 1 \\ 0 & \text{otherwise} \end{cases}
    \label{eq:and} \\[6pt]
  a \vee b \; (\text{OR}) &=
    \begin{cases} 1 & \text{if } a = 1 \text{ or } b = 1 \\ 0 & \text{otherwise} \end{cases}
    \label{eq:or} \\[6pt]
  \bar{a} \; (\text{NOT}) &=
    \begin{cases} 1 & \text{if } a = 0 \\ 0 & \text{if } a = 1 \end{cases}
    \label{eq:not}
\end{align}
\end{definition}

From these, we derive two more operations that appear constantly in this
book:

\begin{definition}[XOR and XNOR]
\label{def:xor}
\begin{align}
  a \oplus b \; (\text{XOR}) &= (a \wedge \bar{b}) \vee (\bar{a} \wedge b)
    = \begin{cases} 1 & \text{if } a \neq b \\ 0 & \text{if } a = b \end{cases}
    \label{eq:xor} \\[6pt]
  \overline{a \oplus b} \; (\text{XNOR}) &=
    \begin{cases} 1 & \text{if } a = b \\ 0 & \text{if } a \neq b \end{cases}
    \label{eq:xnor}
\end{align}
\end{definition}

\begin{engineernote}
XNOR is the ``equality test'' for bits. When we compare a binary weight
with a binary input (as in Binary Neural Networks~\cite{courbariaux2016binarized, rastegari2016xnornet}), XNOR tells us
whether they match. This will become important in \cref{ch:clauses}.
\end{engineernote}


% ----------------------------------------------------------------------------
\section{Truth Tables and the 16 Two-Input Functions}
\label{sec:truth-tables}
% ----------------------------------------------------------------------------

A \concept{truth table} exhaustively lists the output of a Boolean
function for every possible input combination. Two binary inputs
$(a, b)$ have $2^2 = 4$ possible combinations, so a 2-input function
is fully described by 4 output bits.

Since each of those 4 output bits can be 0 or~1, there are
$2^4 = 16$ possible 2-input Boolean functions. We list the six most
important ones:

\begin{table}[htbp]
\centering
\caption{Truth tables for six key 2-input Boolean functions.
  The 4-bit ``code'' column uniquely identifies each function
  by its output for inputs $(0,0), (0,1), (1,0), (1,1)$.}
\label{tab:truth-tables}
\begin{tabular}{@{}cc|cccccc@{}}
  \toprule
  $a$ & $b$ & AND & OR & XOR & NAND & NOR & XNOR \\
  \midrule
  0 & 0 & 0 & 0 & 0 & 1 & 1 & 1 \\
  0 & 1 & 0 & 1 & 1 & 1 & 0 & 0 \\
  1 & 0 & 0 & 1 & 1 & 1 & 0 & 0 \\
  1 & 1 & 1 & 1 & 0 & 0 & 0 & 1 \\
  \midrule
  \multicolumn{2}{c|}{4-bit code} & 0001 & 0111 & 0110 & 1110 & 1000 & 1001 \\
  \bottomrule
\end{tabular}
\end{table}

\begin{keyinsight}
Every 2-input Boolean function can be represented as a 4-bit number
(its truth table outputs read as a binary code). This is how Logic
Gate Networks encode gate types: each gate is identified by one of the
16 possible 4-bit codes~\cite{petersen2022difflogic}. We will use this encoding extensively in
\cref{ch:diff-gates}.
\end{keyinsight}

The complete set of 16 functions is listed in \cref{app:gates}.


% ----------------------------------------------------------------------------
\section{Literals and Clauses}
\label{sec:literals-clauses}
% ----------------------------------------------------------------------------

Two concepts from propositional logic are essential for understanding
Tsetlin Machines~\cite{granmo2018tsetlin}.

\begin{definition}[Literal]
\label{def:literal}
Given a Boolean variable $x_k$, a \concept{literal} is either $x_k$
(the \emph{positive literal}) or $\bar{x}_k$ (the \emph{negative
literal}). If $x_k = 1$, then $x_k$ evaluates to 1 and $\bar{x}_k$
evaluates to 0, and vice versa.
\end{definition}

For $n$ Boolean variables, we have $2n$ available literals.

\begin{definition}[Clause]
\label{def:clause}
A \concept{clause} (specifically, a \emph{conjunctive clause}) is a
logical AND of selected literals:
\begin{equation}
  C = \ell_1 \wedge \ell_2 \wedge \cdots \wedge \ell_m
  \label{eq:clause-def}
\end{equation}
where each $\ell_i$ is a literal drawn from the available set.
A clause outputs 1 if and only if \emph{all} its selected literals
are simultaneously true.
\end{definition}

\begin{example}
Consider three Boolean variables $x_1, x_2, x_3$. The clause
$C = x_1 \wedge \bar{x}_2 \wedge x_3$
outputs 1 only when $x_1 = 1$, $x_2 = 0$, and $x_3 = 1$.
For all other input combinations, $C = 0$.
\end{example}

\begin{analogy}
Think of a clause as a checklist at a factory quality control station.
Each item on the checklist says ``this measurement must be within
spec'' (positive literal) or ``this defect must be absent'' (negative
literal). The product passes inspection (clause outputs~1) only if
every single item on the checklist is satisfied.
\end{analogy}

\subsection{Clauses as Pattern Detectors}

A clause with $m$ selected literals out of $2n$ available is a very
specific pattern detector. It fires (outputs~1) only for inputs that
match \emph{all} $m$ conditions. The more literals a clause includes,
the more specific it becomes---firing for fewer and fewer input
patterns.

This is exactly what we want for digit recognition: a clause that
includes literals for specific pixel positions will fire only when the
image matches a particular pattern of bright and dark pixels.


% ----------------------------------------------------------------------------
\section{Bitwise Parallelism on CPUs}
\label{sec:bitwise-parallel}
% ----------------------------------------------------------------------------

A key engineering advantage of Boolean computation is
\concept{bitwise parallelism}. When a CPU executes an AND instruction
on two 64-bit registers, it performs 64 independent AND operations in
a single clock cycle.

\begin{figure}[htbp]
\centering
\begin{tikzpicture}[
  bit/.style={draw, minimum width=0.4cm, minimum height=0.5cm,
    font=\tiny\ttfamily, inner sep=0pt},
  lbl/.style={font=\small\sffamily, bitgray}
]
  % Register A
  \node[lbl, left] at (-0.3, 1.5) {Reg A:};
  \foreach \b [count=\i from 0] in {1,0,1,1,0,1,0,1,1,1,0,0,1,0,1,0} {
    \node[bit, fill=bitblue!15] at (\i*0.45, 1.5) {\b};
  }
  \node[font=\tiny\sffamily, right] at (7.5, 1.5) {$\cdots$ 64 bits};

  % Register B
  \node[lbl, left] at (-0.3, 0.8) {Reg B:};
  \foreach \b [count=\i from 0] in {1,1,0,1,0,0,1,1,0,1,1,0,0,1,0,1} {
    \node[bit, fill=bitgreen!15] at (\i*0.45, 0.8) {\b};
  }

  % AND operation
  \node[font=\sffamily\bfseries, bitpurple] at (3.4, 0.3) {AND};

  % Result
  \node[lbl, left] at (-0.3, -0.3) {Result:};
  \foreach \b [count=\i from 0] in {1,0,0,1,0,0,0,1,0,1,0,0,0,0,0,0} {
    \node[bit, fill=bitpurple!15] at (\i*0.45, -0.3) {\b};
  }

  % Clock annotation
  \node[draw, rounded corners, fill=lightbg, font=\scriptsize\sffamily]
    at (3.4, -1.1) {1 clock cycle for all 64 bits};
\end{tikzpicture}
\caption{Bitwise parallelism: one AND instruction processes 64 bit
  positions simultaneously. With SIMD extensions (AVX-512), this
  scales to 512~bits per instruction.}
\label{fig:bitwise-parallel}
\end{figure}

\begin{engineernote}
For our MNIST architecture, each image is 784~bits (after
booleanisation). That fits in $\lceil 784 / 64 \rceil = 13$ machine
words. Evaluating whether a clause matches an image requires 13~AND
instructions followed by a check that all bits are set---a total of
roughly 26~instructions. A conventional neuron with 784 floating-point
multiply-adds requires 784~multiplications and 784~additions.

The ratio is approximately $1{,}568 / 26 \approx 60\times$ fewer
instructions per ``neuron equivalent.''
\end{engineernote}

\subsection{The Popcount Instruction}

The \concept{popcount} (population count) instruction counts the
number of 1-bits in a machine word. Most modern CPUs provide this as
a hardware instruction.

\begin{equation}
  \popcount(10110101_2) = 5
  \label{eq:popcount-example}
\end{equation}

Popcount is useful for:
\begin{itemize}
  \item Counting how many literals in a clause matched (clause
    evaluation).
  \item Computing Hamming distance between two bit strings (number
    of positions where they differ).
  \item Replacing summation: $\sum_i a_i$ for binary $a_i$ is just
    $\popcount(\bvec{a})$.
\end{itemize}


% ----------------------------------------------------------------------------
\section{Boolean Functions as Lookup Tables}
\label{sec:bool-lut}
% ----------------------------------------------------------------------------

Any Boolean function of $n$ inputs can be represented as a table with
$2^n$ entries---one entry for each possible input combination. This is
the \concept{lookup table} (LUT) representation.

\begin{example}
The AND function of 2 inputs is the table:

\begin{center}
\begin{tabular}{@{}cc|c@{}}
  \toprule
  \textbf{Input $(a,b)$} & \textbf{Index} & \textbf{Output} \\
  \midrule
  $(0,0)$ & 0 & 0 \\
  $(0,1)$ & 1 & 0 \\
  $(1,0)$ & 2 & 0 \\
  $(1,1)$ & 3 & 1 \\
  \bottomrule
\end{tabular}
\end{center}

To evaluate AND$(a,b)$, compute the index $i = 2a + b$ and return
$\text{table}[i]$. No logic needed---just a memory read.
\end{example}

This observation---that any function can be replaced by a table---is
the foundation of LUT compilation (\cref{ch:lut}). The trade-off is
clear: a table of $2^n$ entries replaces $n$-input computation with a
single memory access, but the table size grows exponentially with
$n$. For small $n$ (2--8 inputs), this is efficient. For $n = 784$,
it is impossible ($2^{784}$ entries).

\begin{keyinsight}
The architecture in this book uses lookup tables at two scales:
\begin{enumerate}
  \item \textbf{Gate level:} Each 2-input logic gate is a 4-entry
    table (tiny, fits in a register).
  \item \textbf{Function level:} Learned scoring functions are sampled
    and stored as $L$-entry tables (typically $L = 64$ to $256$).
\end{enumerate}
The Tsetlin Machine and Logic Gate Network reduce the
dimensionality from 784 to a manageable number of intermediate
values, making LUT compilation tractable.
\end{keyinsight}


% ----------------------------------------------------------------------------
\section{Functional Completeness}
\label{sec:completeness}
% ----------------------------------------------------------------------------

\begin{definition}[Functional completeness]
A set of logic operations is \concept{functionally complete} if every
possible Boolean function can be built from that set alone.
\end{definition}

The set $\{\text{AND}, \text{OR}, \text{NOT}\}$ is functionally
complete. More remarkably, $\{\text{NAND}\}$ alone is functionally
complete---every Boolean function can be built from NAND gates only.

This matters for our architecture because it guarantees that a network
of logic gates (even with a limited gate vocabulary) can in principle
represent \emph{any} Boolean classification function. The Logic Gate
Network has access to all 16 two-input functions, which is more than
sufficient.


% ----------------------------------------------------------------------------
\section{Worked Example: Classifying a Toy Image}
\label{sec:toy-example}
% ----------------------------------------------------------------------------

Let us build intuition with a minimal example. Consider a $2 \times 2$
image with 4 pixels, each either black (0) or white (1). We want to
classify whether the image contains a ``diagonal stripe'' pattern.

The two target patterns are:
\begin{center}
\begin{tikzpicture}[
  px/.style={draw, minimum size=0.8cm, font=\small\ttfamily}
]
  % Pattern 1
  \node[px, fill=bitblue!30] at (0, 0.8) {1};
  \node[px] at (0.8, 0.8) {0};
  \node[px] at (0, 0) {0};
  \node[px, fill=bitblue!30] at (0.8, 0) {1};
  \node[font=\small\sffamily, below] at (0.4, -0.4) {Pattern A};

  % Pattern 2
  \node[px] at (3, 0.8) {0};
  \node[px, fill=bitblue!30] at (3.8, 0.8) {1};
  \node[px, fill=bitblue!30] at (3, 0) {1};
  \node[px] at (3.8, 0) {0};
  \node[font=\small\sffamily, below] at (3.4, -0.4) {Pattern B};
\end{tikzpicture}
\end{center}

\paragraph{Step 1: Define variables.}
Let $x_1, x_2, x_3, x_4$ be the four pixel values (top-left,
top-right, bottom-left, bottom-right).

\paragraph{Step 2: Write clauses.}
\begin{align*}
  C_1 &= x_1 \wedge \bar{x}_2 \wedge \bar{x}_3 \wedge x_4
    && \text{(detects Pattern A)} \\
  C_2 &= \bar{x}_1 \wedge x_2 \wedge x_3 \wedge \bar{x}_4
    && \text{(detects Pattern B)}
\end{align*}

\paragraph{Step 3: Combine with OR.}
\[
  \text{is\_diagonal}(\bvec{x}) = C_1 \vee C_2
\]

This outputs 1 if and only if the input matches either diagonal
pattern. The entire classifier uses only AND, OR, and NOT---no
multiplication, no floating-point arithmetic.

\begin{engineernote}
In this toy example, we hand-crafted the clauses. The power of the
Tsetlin Machine (Part~II) is that it \emph{discovers} such clauses
automatically from training data, without human intervention.
\end{engineernote}

% ============================================================================
\chapter{Supervised Learning from an Engineer's Perspective}
\label{ch:learning}
% ============================================================================

This chapter introduces the concepts of supervised learning---the
framework used to train neural networks from labelled data. If you
have a machine learning background, you can skim this chapter; it
establishes the vocabulary and intuition needed for the training
methods in Parts~II--IV.


% ----------------------------------------------------------------------------
\section{Classification as Function Approximation}
\label{sec:classification}
% ----------------------------------------------------------------------------

A \concept{classifier} is a function that maps an input (an image,
a sensor reading, a document) to one of several predefined categories.
For MNIST:
\begin{equation}
  f : \{0, 1, \ldots, 255\}^{784} \to \{0, 1, 2, \ldots, 9\}
  \label{eq:classifier}
\end{equation}

We do not know the ``true'' function $f^*$ that perfectly classifies
every possible handwritten digit. Instead, we have a dataset of
examples $({\bvec{x}_i, y_i})$ and we want to find a function
$\hat{f}$ that agrees with these examples and \emph{generalises} to
unseen inputs.

\begin{analogy}
Learning a classifier from data is like learning to identify birds from
a field guide. You study examples (``this is a robin, this is a
sparrow''), and eventually you can identify birds you have never seen
before---not because you memorised every bird, but because you learned
the distinguishing features (beak shape, colouring, size).
\end{analogy}


% ----------------------------------------------------------------------------
\section{Training Data, Test Data, and Overfitting}
\label{sec:overfitting}
% ----------------------------------------------------------------------------

The 70,000 MNIST~\cite{lecun1998mnist} images are split into:
\begin{itemize}
  \item \textbf{Training set} (60,000 images): used to adjust the
    model's parameters.
  \item \textbf{Test set} (10,000 images): used only to evaluate
    accuracy. The model never sees these during training.
\end{itemize}

\concept{Overfitting} occurs when a model memorises the training
examples instead of learning general patterns. An overfitting model
achieves high training accuracy but poor test accuracy.

\begin{keyinsight}
The test set is the only honest measure of a model's quality. Any
accuracy number reported in this book refers to the test set unless
explicitly stated otherwise.
\end{keyinsight}


% ----------------------------------------------------------------------------
\section{The Traditional Artificial Neuron}
\label{sec:neuron}
% ----------------------------------------------------------------------------

The building block of conventional neural networks is the
\concept{artificial neuron}. It computes:

\begin{equation}
  y = \sigma\!\left(\sum_{i=1}^{n} w_i \cdot x_i + b\right)
    = \sigma(\bvec{w} \cdot \bvec{x} + b)
  \label{eq:neuron}
\end{equation}

where:
\begin{itemize}
  \item $\bvec{x} = (x_1, \ldots, x_n)$ is the input vector.
  \item $\bvec{w} = (w_1, \ldots, w_n)$ is the \concept{weight}
    vector (learned parameters).
  \item $b$ is the \concept{bias} (a learned offset).
  \item $\sigma$ is the \concept{activation function} (a nonlinear
    function such as ReLU: $\sigma(z) = \max(0, z)$).
\end{itemize}

\begin{figure}[htbp]
\centering
\begin{tikzpicture}[
  input/.style={circle, draw, minimum size=0.6cm, font=\small},
  neuron/.style={circle, draw, minimum size=1cm, font=\small,
    fill=bitblue!10, line width=0.8pt},
  >=Stealth
]
  \foreach \i/\v in {1/x_1, 2/x_2, 3/x_n} {
    \ifnum\i=3
      \node[input] (in\i) at (0, 3-\i*1.2) {$\v$};
    \else
      \node[input] (in\i) at (0, 3-\i*1.2) {$\v$};
    \fi
  }
  \node at (0, -0.4) {$\vdots$};

  \node[neuron] (n) at (3.5, 0.6) {$\Sigma$};
  \node[draw, rounded corners, fill=bitorange!10, minimum width=0.8cm,
    font=\small] (act) at (5.5, 0.6) {$\sigma$};
  \node[right=0.5cm, font=\small] (out) at (act.east) {$y$};

  \foreach \i in {1,2,3} {
    \draw[->, thick] (in\i) -- node[above, font=\scriptsize, sloped] {$w_\i$} (n);
  }
  \draw[->, thick] (n) -- (act);
  \draw[->, thick] (act) -- (out);
\end{tikzpicture}
\caption{A traditional artificial neuron. The weighted sum
  $\sum w_i x_i + b$ is computed first, then passed through a nonlinear
  activation function $\sigma$.}
\label{fig:neuron}
\end{figure}

The weighted sum $\bvec{w} \cdot \bvec{x} + b$ is the
\concept{multiply-accumulate} (MAC) operation that dominates the cost
of conventional networks. For a neuron with 784 inputs (one per MNIST
pixel), this requires 784 floating-point multiplications and 784
floating-point additions per neuron.


% ----------------------------------------------------------------------------
\section{Loss Functions}
\label{sec:loss}
% ----------------------------------------------------------------------------

To train a model, we need a way to measure \emph{how wrong} its
predictions are. This measurement is called the \concept{loss function}
(or cost function).

For 10-class classification (MNIST), the standard loss is
\concept{cross-entropy loss}. The model outputs a probability
distribution over 10 classes, and we penalise it for assigning low
probability to the correct class.

\begin{equation}
  \mathcal{L} = -\log p_{\hat{y}}
  \label{eq:cross-entropy}
\end{equation}

where $p_{\hat{y}}$ is the model's predicted probability for the
correct class $\hat{y}$.

\begin{example}
Suppose the true label is ``3'' and the model outputs:
\[
  \bvec{p} = (0.01, 0.02, 0.05, \mathbf{0.80}, 0.03, 0.01, 0.02, 0.01, 0.03, 0.02)
\]
The loss is $\mathcal{L} = -\log(0.80) \approx 0.22$.

If the model were less confident, outputting $p_3 = 0.10$, the loss
would be $\mathcal{L} = -\log(0.10) \approx 2.30$---much larger,
reflecting the worse prediction.
\end{example}

\begin{engineernote}
Not all paradigms use cross-entropy loss. Tsetlin Machines
(\cref{ch:clauses}) use a completely different training signal based on
reinforcement learning. Logic Gate Networks (\cref{ch:diff-gates}) use
cross-entropy or similar differentiable losses. Understanding the
concept of ``measuring wrongness'' is the key; the specific formula
varies.
\end{engineernote}


% ----------------------------------------------------------------------------
\section{Gradient Descent}
\label{sec:gradient-descent}
% ----------------------------------------------------------------------------

Given a loss function that measures prediction error, \concept{gradient
descent} is the algorithm that adjusts the model's parameters to reduce
that error.

\begin{enumerate}
  \item \textbf{Forward pass:} Run an input through the model, get a
    prediction.
  \item \textbf{Compute loss:} Measure how wrong the prediction is.
  \item \textbf{Compute gradient:} For each parameter $w$, compute
    $\frac{\partial \mathcal{L}}{\partial w}$---how much the loss
    changes if we nudge $w$ slightly.
  \item \textbf{Update:} Move each parameter in the direction that
    reduces the loss:
    \begin{equation}
      w \leftarrow w - \eta \cdot \frac{\partial \mathcal{L}}{\partial w}
      \label{eq:gradient-update}
    \end{equation}
    where $\eta > 0$ is the \concept{learning rate} (step size).
\end{enumerate}

\begin{analogy}
Imagine you are blindfolded on a hillside and want to reach the valley.
You feel the slope under your feet (the gradient) and take a step
downhill. Repeat. You will eventually reach the bottom---maybe not the
absolute lowest point, but a low one. Gradient descent does the same in
the high-dimensional space of model parameters.
\end{analogy}

\begin{figure}[htbp]
\centering
\begin{tikzpicture}[>=Stealth]
  \begin{axis}[
    width=9cm, height=5cm,
    xlabel={Parameter $w$}, ylabel={Loss $\mathcal{L}$},
    xmin=-3, xmax=3, ymin=0, ymax=5,
    axis lines=middle,
    every axis label/.style={font=\small\sffamily},
    ticks=none
  ]
    % Loss curve
    \addplot[thick, bitblue, smooth, domain=-3:3, samples=100]
      {0.5*x^2 + 0.3*sin(deg(2*x)) + 1};

    % Gradient descent steps
    \addplot[only marks, mark=*, mark size=2.5pt, bitorange]
      coordinates {(2, 3.21) (1.2, 1.92) (0.5, 1.41) (0.1, 1.11)};
    \draw[->, thick, bitorange] (2, 3.21) -- (1.2, 1.92);
    \draw[->, thick, bitorange] (1.2, 1.92) -- (0.5, 1.41);
    \draw[->, thick, bitorange] (0.5, 1.41) -- (0.1, 1.11);

    \node[font=\scriptsize\sffamily, bitorange, right] at (2.1, 3.5)
      {start};
  \end{axis}
\end{tikzpicture}
\caption{Gradient descent iteratively moves the parameter $w$ in the
  direction that reduces the loss. Each step follows the negative
  gradient (downhill).}
\label{fig:gradient-descent}
\end{figure}


% ----------------------------------------------------------------------------
\section{Backpropagation}
\label{sec:backprop}
% ----------------------------------------------------------------------------

In a multi-layer network, computing $\frac{\partial \mathcal{L}}
{\partial w}$ for a parameter deep in the network requires the
\concept{chain rule} of calculus.

If the network computes $y = f(g(h(x)))$, then:
\begin{equation}
  \frac{\partial y}{\partial x} =
    \frac{\partial f}{\partial g} \cdot
    \frac{\partial g}{\partial h} \cdot
    \frac{\partial h}{\partial x}
  \label{eq:chain-rule}
\end{equation}

\concept{Backpropagation} applies this chain rule layer by layer, from
the output back to the input, efficiently computing gradients for every
parameter in the network.

\begin{engineernote}
Backpropagation requires that every operation in the forward pass be
\emph{differentiable}---that is, it must have a well-defined
derivative. This is the central challenge for bitwise networks:
discrete operations like AND and OR are \emph{not} differentiable.
Their output jumps from 0 to~1 with no smooth transition, so the
derivative is zero everywhere (useless for learning) or undefined
(at the jump). The straight-through estimator~\cite{bengio2013ste} is one common
workaround for this problem.
\end{engineernote}


% ----------------------------------------------------------------------------
\section{The Differentiability Problem}
\label{sec:diff-problem}
% ----------------------------------------------------------------------------

Consider the AND function: $f(a, b) = a \wedge b$. The function maps
$\{0,1\}^2$ to $\{0,1\}$---it is defined only at four points. Between
those points, there is no continuous surface, so there is no gradient.

\begin{figure}[htbp]
\centering
\begin{tikzpicture}[>=Stealth]
  % Continuous neuron
  \begin{scope}
    \node[font=\sffamily\bfseries] at (2, 3.2) {Continuous (ReLU)};
    \begin{axis}[
      at={(0,0)}, width=4.5cm, height=3.5cm,
      xmin=-2, xmax=2, ymin=-0.5, ymax=2.5,
      axis lines=middle, ticks=none,
      every axis label/.style={font=\tiny\sffamily}
    ]
      \addplot[very thick, bitblue, domain=-2:2, samples=50]
        {max(0, x)};
      \node[font=\tiny\sffamily, bitblue] at (axis cs:1.5, 2)
        {smooth};
      \draw[->, thick, bitgreen] (axis cs:1, 1) -- (axis cs:1.3, 1)
        node[right, font=\tiny\sffamily, bitgreen] {$\nabla$};
    \end{axis}
  \end{scope}

  % Discrete gate
  \begin{scope}[xshift=6cm]
    \node[font=\sffamily\bfseries] at (2, 3.2) {Discrete (AND gate)};
    \begin{axis}[
      at={(0,0)}, width=4.5cm, height=3.5cm,
      xmin=-0.5, xmax=1.5, ymin=-0.5, ymax=1.5,
      axis lines=middle, ticks=none,
      every axis label/.style={font=\tiny\sffamily}
    ]
      \addplot[only marks, mark=*, mark size=3pt, bitorange]
        coordinates {(0,0) (1,0) (0,0) (1,1)};
      \node[font=\tiny\sffamily, bitorange] at (axis cs:0.8, 0.5)
        {no gradient};
    \end{axis}
  \end{scope}
\end{tikzpicture}
\caption{Left: a continuous activation (ReLU) has a gradient that
  tells us which direction to move. Right: a discrete logic gate has
  no gradient---its outputs are isolated points with no smooth
  connection.}
\label{fig:diff-problem}
\end{figure}

This is the fundamental obstacle that every bitwise neural network must
overcome. The three paradigms in this book handle it in three different
ways:


% ----------------------------------------------------------------------------
\section{Three Solutions to the Differentiability Problem}
\label{sec:three-solutions}
% ----------------------------------------------------------------------------

\begin{description}[leftmargin=1em]
  \item[Tsetlin Machines~\cite{granmo2018tsetlin}: avoid gradients entirely.]
    Instead of computing derivatives, Tsetlin Machines use
    \emph{reinforcement learning}---simple reward/penalty signals that
    nudge automata states up or down. No calculus required. The
    learning rule is: ``if the clause helped classify correctly,
    reinforce it; if it caused an error, weaken it.'' This is covered
    in Chapters~\ref{ch:automaton}--\ref{ch:tm-mnist}.

  \item[Logic Gate Networks~\cite{petersen2022difflogic}: use a continuous relaxation.]
    During training, each node does not commit to a single gate type.
    Instead, it maintains a \emph{probability distribution} over all 16
    gate types, and its output is a differentiable weighted average.
    Gradients flow through this weighted average. As training
    progresses, the distribution sharpens until each node converges to
    a single gate. This is covered in Chapters~\ref{ch:diff-gates}
    --\ref{ch:lgn-training}.

  \item[KAN/LUT~\cite{liu2024kan}: train in continuous space, deploy discretely.]
    During training, learnable functions are smooth B-spline curves~\cite{deboor1978splines}
    that support standard backpropagation. After training, these
    curves are \emph{compiled} into discrete lookup tables. The
    discretisation happens only once, at deployment time. This is
    covered in Chapters~\ref{ch:splines}--\ref{ch:lut}.
\end{description}

\begin{keyinsight}
Each paradigm makes a different trade-off:
\begin{itemize}
  \item Tsetlin: no gradients at all (simplest training, slowest
    convergence).
  \item LGN: gradients through a relaxation (standard training, must
    anneal to discrete).
  \item KAN/LUT: full gradients during training, discretisation at
    deployment (highest training cost, smallest deployment cost).
\end{itemize}
The hybrid architecture uses all three, each where it is most
effective.
\end{keyinsight}

% ============================================================================
\chapter{Preparing MNIST for Bitwise Processing}
\label{ch:mnist-prep}
% ============================================================================

Before any of our three paradigms can process MNIST~\cite{lecun1998mnist} images, the raw
pixel data must be converted to a Boolean representation. This chapter
covers the conversion methods and their trade-offs.


% ----------------------------------------------------------------------------
\section{The Raw Data}
\label{sec:raw-data}
% ----------------------------------------------------------------------------

Each MNIST image is a $28 \times 28$ grid of pixels. Each pixel is an
unsigned 8-bit integer in the range $[0, 255]$, where 0~is black
(background) and 255~is white (ink).

In practice, the images are stored as a flat array of 784 bytes:
\begin{equation}
  \bvec{x}_{\text{raw}} = (p_1, p_2, \ldots, p_{784}),
  \quad p_i \in \{0, 1, \ldots, 255\}
  \label{eq:raw-mnist}
\end{equation}

\begin{engineernote}
MNIST pixels are heavily bimodal: most pixels are either very dark
(background, $p \approx 0$) or quite bright (ink, $p > 200$). The
middle values ($50$--$150$) are relatively rare, occurring mainly at
the edges of strokes. This bimodal distribution is why simple
thresholding works well.
\end{engineernote}


% ----------------------------------------------------------------------------
\section{Simple Thresholding}
\label{sec:thresholding}
% ----------------------------------------------------------------------------

The simplest conversion to Boolean is \concept{thresholding}: compare
each pixel to a fixed threshold $\theta$ and output 1 if the pixel is
bright enough, 0 otherwise.

\begin{equation}
  x_k = \begin{cases}
    1 & \text{if } p_k > \theta \\
    0 & \text{if } p_k \leq \theta
  \end{cases}
  \label{eq:threshold}
\end{equation}

A common choice is $\theta = 127$ (the midpoint of the $[0, 255]$
range).

\begin{figure}[htbp]
\centering
\begin{tikzpicture}[
  px/.style={draw, minimum size=0.35cm, inner sep=0pt, font=\tiny\ttfamily}
]
  % Original
  \node[font=\small\sffamily\bfseries] at (2, 3.5) {Grayscale};
  \foreach \x/\y/\v in {
    0/2/0, 1/2/30, 2/2/180, 3/2/255,
    0/1/0, 1/1/200, 2/1/255, 3/1/150,
    0/0/0, 1/0/50, 2/0/220, 3/0/0
  }{
    \pgfmathsetmacro{\intensity}{\v/255*100}
    \node[px, fill=bitblue!\intensity] at (\x*0.5, \y*0.5) {\v};
  }

  % Arrow
  \draw[-{Stealth}, thick] (2.5, 0.75) -- (4, 0.75)
    node[midway, above, font=\scriptsize\sffamily] {$\theta = 127$};

  % Binary
  \node[font=\small\sffamily\bfseries] at (6, 3.5) {Boolean};
  \foreach \x/\y/\v in {
    0/2/0, 1/2/0, 2/2/1, 3/2/1,
    0/1/0, 1/1/1, 2/1/1, 3/1/1,
    0/0/0, 1/0/0, 2/0/1, 3/0/0
  }{
    \ifnum\v=1
      \def\fillcol{bitblue!40}
    \else
      \def\fillcol{white}
    \fi
    \node[px, fill=\fillcol] at (\x*0.5 + 4.5, \y*0.5) {\v};
  }
\end{tikzpicture}
\caption{Thresholding a $4 \times 3$ pixel patch at $\theta = 127$.
  Pixels above the threshold become~1 (foreground); others become~0
  (background).}
\label{fig:thresholding}
\end{figure}

After thresholding, the image is a bit string:
\begin{equation}
  \bvec{x} = (x_1, x_2, \ldots, x_{784}) \in \{0, 1\}^{784}
  \label{eq:bool-mnist}
\end{equation}


% ----------------------------------------------------------------------------
\section{Literal Expansion}
\label{sec:literal-expansion}
% ----------------------------------------------------------------------------

Tsetlin Machines~\cite{granmo2018tsetlin} need access to both a feature and its negation. For
every Boolean feature $x_k$, we also provide $\bar{x}_k = 1 - x_k$.

\begin{equation}
  \bvec{x}_{\text{expanded}} =
    (x_1, \bar{x}_1, x_2, \bar{x}_2, \ldots, x_{784}, \bar{x}_{784})
    \in \{0, 1\}^{1568}
  \label{eq:expanded}
\end{equation}

\begin{keyinsight}
Why include negations? Because a clause that says ``pixel~42 must be
dark'' needs the literal $\bar{x}_{42}$, which is 1 when the pixel is
dark (below threshold). Without negations, clauses could only express
``pixel must be bright'' conditions, halving their expressive power.
\end{keyinsight}

\begin{example}
For a $2 \times 2$ image with pixels $(200, 50, 180, 10)$ and
$\theta = 127$:
\begin{align*}
  \text{Thresholded:}\quad & (1, 0, 1, 0) \\
  \text{Expanded:}\quad & (1, 0, \; 0, 1, \; 1, 0, \; 0, 1)
\end{align*}
The expanded vector has 8 entries (twice the original). Notice that
$x_k$ and $\bar{x}_k$ are always complementary: exactly one of each
pair is~1.
\end{example}


% ----------------------------------------------------------------------------
\section{Thermometer Encoding}
\label{sec:thermometer}
% ----------------------------------------------------------------------------

Simple thresholding discards all information about pixel brightness
beyond ``above or below 127.'' A pixel value of 130 and a pixel value
of 255 both become~1, even though the latter is much brighter.

\concept{Thermometer encoding} uses multiple thresholds per pixel to
retain finer information:

\begin{definition}[Thermometer encoding]
Given $L$ threshold levels
$\theta_1 < \theta_2 < \cdots < \theta_L$, pixel $p_k$ is encoded as:
\begin{equation}
  t_k^{(j)} = \begin{cases}
    1 & \text{if } p_k > \theta_j \\
    0 & \text{otherwise}
  \end{cases}
  \quad \text{for } j = 1, \ldots, L
  \label{eq:thermometer}
\end{equation}
\end{definition}

\begin{example}
With thresholds $(50, 100, 150, 200)$, a pixel value of 175 encodes
as $(1, 1, 1, 0)$. A pixel value of 80 encodes as $(1, 0, 0, 0)$.
The encoding looks like a thermometer: 1s fill up from the bottom as
the value increases.
\end{example}

\begin{figure}[htbp]
\centering
\begin{tikzpicture}[
  bit/.style={draw, minimum width=0.6cm, minimum height=0.5cm,
    font=\scriptsize\ttfamily, inner sep=0pt}
]
  % Pixel value
  \node[font=\small\sffamily] at (-1.5, 0) {$p_k = 175$};

  % Thresholds
  \foreach \t/\i/\v in {50/0/1, 100/1/1, 150/2/1, 200/3/0} {
    \ifnum\v=1
      \def\fillcol{bitblue!30}
    \else
      \def\fillcol{white}
    \fi
    \node[bit, fill=\fillcol] (t\i) at (\i*0.8, 0) {\v};
    \node[font=\tiny\sffamily, bitgray, above] at (\i*0.8, 0.35) {$>\t$};
  }

  \draw[decorate, decoration={brace, amplitude=4pt, mirror}, thick]
    (-0.3, -0.5) -- (2.1, -0.5)
    node[midway, below=4pt, font=\scriptsize\sffamily] {4 bits per pixel};
\end{tikzpicture}
\caption{Thermometer encoding with 4 thresholds. The pixel value 175
  exceeds thresholds 50, 100, and 150 but not 200, yielding the code
  $(1, 1, 1, 0)$.}
\label{fig:thermometer}
\end{figure}

With $L$ thresholds per pixel, plus negations, the total input size is:
\begin{equation}
  \text{Input length} = 784 \times L \times 2
  \label{eq:thermo-size}
\end{equation}

For $L = 4$: $784 \times 4 \times 2 = 6{,}272$ literals. More
information is preserved, but the model must process a larger input.


% ----------------------------------------------------------------------------
\section{The Booleanisation Trade-Off}
\label{sec:bool-tradeoff}
% ----------------------------------------------------------------------------

\begin{table}[htbp]
\centering
\caption{Comparison of booleanisation strategies for MNIST.}
\label{tab:bool-comparison}
\begin{tabular}{@{}lccc@{}}
  \toprule
  \textbf{Method} & \textbf{Bits/pixel} & \textbf{Total literals}
    & \textbf{Info preserved} \\
  \midrule
  Simple threshold ($\theta = 127$) & 1 & 1{,}568 & Low \\
  Thermometer ($L = 4$) & 4 & 6{,}272 & Medium \\
  Thermometer ($L = 8$) & 8 & 12{,}544 & High \\
  Full 8-bit binary    & 8 & 12{,}544 & Full \\
  \bottomrule
\end{tabular}
\end{table}

\begin{engineernote}
For the proof-of-concept implementation described in this book, we
recommend starting with simple thresholding ($\theta = 127$). It
produces 1{,}568 literals, keeps the model small and fast, and achieves
accuracy within 1--2\% of the more sophisticated encodings. Once the
architecture is working, thermometer encoding can be added as an
optimisation.
\end{engineernote}


% ----------------------------------------------------------------------------
\section{Worked Example: One MNIST Digit Through Each Encoding}
\label{sec:encoding-walkthrough}
% ----------------------------------------------------------------------------

Let us trace a single pixel row from an MNIST ``3'' through the
encoding pipeline. Consider pixels at row 14 (the middle of the
image), positions 10--17, with raw values:

\begin{center}
\begin{tabular}{@{}l*{8}{c}@{}}
  \toprule
  \textbf{Position} & 10 & 11 & 12 & 13 & 14 & 15 & 16 & 17 \\
  \midrule
  Raw value & 0 & 42 & 190 & 255 & 255 & 180 & 30 & 0 \\
  \midrule
  Threshold ($\theta = 127$) & 0 & 0 & 1 & 1 & 1 & 1 & 0 & 0 \\
  \midrule
  Thermo ($>50$) & 0 & 0 & 1 & 1 & 1 & 1 & 0 & 0 \\
  Thermo ($>100$) & 0 & 0 & 1 & 1 & 1 & 1 & 0 & 0 \\
  Thermo ($>150$) & 0 & 0 & 1 & 1 & 1 & 1 & 0 & 0 \\
  Thermo ($>200$) & 0 & 0 & 0 & 1 & 1 & 0 & 0 & 0 \\
  \bottomrule
\end{tabular}
\end{center}

The simple threshold captures the stroke location (positions 12--15
are ink). The thermometer encoding additionally reveals that positions
13--14 are the \emph{darkest} part of the stroke (exceeding even the
200 threshold), while positions 12 and~15 are lighter edges.

After literal expansion, position~12 with simple thresholding
contributes two literals: $x_{12} = 1$ and $\bar{x}_{12} = 0$.
With 4-level thermometer encoding, it contributes 8 literals:
$(1, 0, 1, 0, 1, 0, 0, 1)$ for the four thresholds and their
negations.


% Part II: Tsetlin Machines
\part{Tsetlin Machines --- Learning Logic Without Gradients}
\label{part:tsetlin}

% ============================================================================
\chapter{The Tsetlin Automaton}
\label{ch:automaton}
% ============================================================================

The Tsetlin Machine is built from a single, remarkably simple
component: the \concept{Tsetlin Automaton}. This chapter introduces the
automaton, explains its behaviour, and shows why it can learn from
noisy feedback without any calculus.


% ----------------------------------------------------------------------------
\section{Finite-State Learning Automata}
\label{sec:automata-intro}
% ----------------------------------------------------------------------------

A \concept{learning automaton} is a device that repeatedly chooses an
action, receives feedback from its environment, and updates its
internal state to improve future choices. The idea dates to the 1960s
and the work of M.~L.~Tsetlin~\cite{tsetlin1961}.

The simplest version has:
\begin{itemize}
  \item A finite set of \emph{states} (numbered 1 through $2N$ for
    some positive integer $N$).
  \item Two \emph{actions}: action~0 and action~1.
  \item A rule that maps the current state to an action.
  \item An update rule that moves the state up or down based on
    feedback (reward or penalty).
\end{itemize}

\begin{analogy}
Think of a physical slider on a rail, marked with positions 1 through
$2N$. The slider starts somewhere in the middle. A ``reward'' pushes
it toward the nearest end; a ``penalty'' pushes it toward the centre.
Once the slider reaches an end, it tends to stay there. The two ends
represent two different decisions, and the slider's position reflects
the automaton's confidence in that decision.
\end{analogy}


% ----------------------------------------------------------------------------
\section{The Two-Action Tsetlin Automaton}
\label{sec:two-action}
% ----------------------------------------------------------------------------

\begin{definition}[Two-action Tsetlin Automaton with $2N$ states]
\label{def:tsetlin-automaton}
The automaton has states $\{1, 2, \ldots, 2N\}$ and two actions:
\begin{equation}
  \text{action}(s) = \begin{cases}
    0 & \text{if } s \leq N \\
    1 & \text{if } s > N
  \end{cases}
  \label{eq:action}
\end{equation}

States 1 through $N$ map to action~0 (the ``exclude'' action).
States $N{+}1$ through $2N$ map to action~1 (the ``include'' action).
\end{definition}

The state transition rules depend on the feedback signal:

\begin{description}[leftmargin=1em]
  \item[Reward:] Move the state away from the decision boundary
    (toward the nearest end of the rail):
    \begin{equation}
      s \leftarrow \begin{cases}
        \max(1, \; s - 1) & \text{if } s \leq N \\
        \min(2N, \; s + 1) & \text{if } s > N
      \end{cases}
      \label{eq:reward}
    \end{equation}

  \item[Penalty:] Move the state toward the decision boundary
    (toward the centre):
    \begin{equation}
      s \leftarrow \begin{cases}
        \min(N, \; s + 1) & \text{if } s \leq N \\
        \max(N{+}1, \; s - 1) & \text{if } s > N
      \end{cases}
      \label{eq:penalty}
    \end{equation}
\end{description}

Note that the penalty can push the state across the boundary (from
$N{+}1$ to $N$ or vice versa), changing the action. This is how the
automaton changes its mind.

\begin{figure}[htbp]
\centering
\begin{tikzpicture}[
  state/.style={circle, draw, minimum size=0.9cm, font=\small\ttfamily,
    inner sep=0pt, line width=0.6pt},
  action0/.style={state, fill=bitorange!15},
  action1/.style={state, fill=bitblue!15},
  trans/.style={-{Stealth[length=4pt]}, thick},
  >=Stealth
]
  % States for N=4 (8 states)
  \foreach \i in {1,...,4} {
    \node[action0] (s\i) at (\i*1.6, 0) {\i};
  }
  \foreach \i in {5,...,8} {
    \node[action1] (s\i) at (\i*1.6, 0) {\i};
  }

  % Decision boundary
  \draw[dashed, thick, bitgray] (7.2, -1.2) -- (7.2, 1.2)
    node[above, font=\scriptsize\sffamily] {boundary};

  % Action labels
  \node[font=\small\sffamily\bfseries, bitorange] at (4, -1.5)
    {Action 0 (exclude)};
  \node[font=\small\sffamily\bfseries, bitblue] at (10.4, -1.5)
    {Action 1 (include)};

  % Reward arrows (above)
  \draw[trans, bitgreen] (s2.north) to[bend left=30]
    node[above, font=\tiny\sffamily] {reward} (s1.north);
  \draw[trans, bitgreen] (s7.north) to[bend left=30]
    node[above, font=\tiny\sffamily] {reward} (s8.north);

  % Penalty arrows (below)
  \draw[trans, red!70] (s2.south) to[bend right=30]
    node[below, font=\tiny\sffamily] {penalty} (s3.south);
  \draw[trans, red!70] (s7.south) to[bend right=30]
    node[below, font=\tiny\sffamily] {penalty} (s6.south);

  % Boundary crossing
  \draw[trans, red!70] (s5.south) to[bend right=30]
    node[below, font=\tiny\sffamily] {penalty} (s4.south);
\end{tikzpicture}
\caption{A Tsetlin Automaton with $2N = 8$ states. States 1--4 map to
  action~0 (exclude); states 5--8 map to action~1 (include). Rewards
  push toward the ends; penalties push toward the centre boundary.
  A penalty at state~5 crosses the boundary to state~4, flipping the
  action from 1 to~0.}
\label{fig:automaton}
\end{figure}


% ----------------------------------------------------------------------------
\section{Why the Automaton Works: Noise Absorption}
\label{sec:noise-absorption}
% ----------------------------------------------------------------------------

The key property of the Tsetlin Automaton is its ability to resist
noise. Consider an automaton in state~1 (deep in the action~0 region).
Even if it receives a few spurious penalty signals, it must be pushed
through states 2, 3, \ldots, $N$, $N{+}1$ before it changes its
action. The parameter $N$ controls this \concept{noise margin}: larger
$N$ means the automaton needs more consecutive penalties to change its
mind.

\begin{keyinsight}
The number of states $N$ per action controls the trade-off between
\emph{stability} (resistance to noise) and \emph{agility} (speed of
adaptation). $N = 1$ gives a memoryless flip-flop that changes on
every penalty. Large $N$ gives a conservative automaton that changes
slowly but reliably.
\end{keyinsight}

\begin{example}
With $N = 100$ and the automaton in state~1, it takes at least 100
consecutive penalties with no intervening rewards to flip the
action. If the correct action is truly~0, the environment will
produce more rewards than penalties, and the state will drift toward
1 (maximum confidence). If the correct action is truly~1, the
accumulated penalties will eventually push the state past 100 to the
action~1 side, where rewards will then anchor it.
\end{example}


% ----------------------------------------------------------------------------
\section{Convergence Guarantees}
\label{sec:convergence}
% ----------------------------------------------------------------------------

Tsetlin~\cite{tsetlin1961} proved the following result (stated informally):

\begin{theorem}[Tsetlin convergence, informal]
\label{thm:tsetlin-convergence}
If one action produces rewards more frequently than the other, a
two-action Tsetlin Automaton with sufficiently many states will
converge to the better action with probability approaching~1 as
$N \to \infty$.
\end{theorem}

More precisely, let $r_0$ be the reward probability for action~0 and
$r_1$ the reward probability for action~1. If $r_1 > r_0$, then the
probability that the automaton is in an action~1 state converges to~1
as $N \to \infty$.

\begin{engineernote}
In practice, $N$ between 50 and 200 works well for MNIST. The
convergence is fast enough that the automaton settles within a few
training epochs (passes through the data). The default in most
implementations is $N = 100$, giving 200 states per automaton.
\end{engineernote}

The convergence does \emph{not} require independence between
successive feedback signals~\cite{zhang2021convergence}. This is important because in a Tsetlin
Machine, the feedback to each automaton depends on the current input
pattern and the states of other automata---it is highly correlated.
Tsetlin's original proof accommodates this, which is one reason the
framework is robust.


% ----------------------------------------------------------------------------
\section{Multiple Automata Working Together}
\label{sec:team}
% ----------------------------------------------------------------------------

A single automaton makes a single binary decision: include or exclude.
The power of the Tsetlin Machine~\cite{granmo2018tsetlin} comes from assembling many automata
into a \concept{team} that collectively builds a pattern detector (a
clause).

Consider a clause with 10 input features ($x_1, \ldots, x_5$ and
their negations $\bar{x}_1, \ldots, \bar{x}_5$). We assign one
automaton to each of the 10 literals. Each automaton independently
decides whether its literal is ``included'' (action~1) or
``excluded'' (action~0) from the clause.

\begin{figure}[htbp]
\centering
\begin{tikzpicture}[
  auto/.style={draw, rounded corners, minimum width=1.4cm,
    minimum height=0.7cm, font=\scriptsize\sffamily, inner sep=2pt},
  inc/.style={auto, fill=bitblue!15},
  exc/.style={auto, fill=bitgray!10},
  >=Stealth
]
  % Literals
  \foreach \i/\lit/\act in {
    0/$x_1$/include,
    1/$\bar{x}_1$/exclude,
    2/$x_2$/include,
    3/$\bar{x}_2$/include,
    4/$x_3$/exclude
  } {
    \pgfmathsetmacro{\xpos}{\i * 2.5}
    \ifnum\i=1
      \node[exc] (a\i) at (\xpos, 0) {\lit: \act};
    \else\ifnum\i=4
      \node[exc] (a\i) at (\xpos, 0) {\lit: \act};
    \else
      \node[inc] (a\i) at (\xpos, 0) {\lit: \act};
    \fi\fi
  }

  % Clause output
  \node[draw, rounded corners, fill=bitgreen!15, minimum width=2cm,
    font=\small\sffamily\bfseries] (clause) at (5, -1.5)
    {$C = x_1 \wedge \bar{x}_2$};

  % Note
  \node[font=\scriptsize\sffamily, bitgray] at (5, -2.3)
    {(only included literals appear in the clause)};

  % Arrows from included automata to clause
  \draw[->, thick] (a0.south) -- (clause);
  \draw[->, thick] (a2.south) -- (clause);
  \draw[->, thick] (a3.south) -- (clause);
\end{tikzpicture}
\caption{A team of automata forming a clause. Each automaton controls
  one literal. The clause is the AND of all included literals.
  Here, automata for $x_1$, $x_2$, and $\bar{x}_2$ chose ``include,''
  so $C = x_1 \wedge x_2 \wedge \bar{x}_2$. (In practice, including
  both $x_2$ and $\bar{x}_2$ would make the clause always false---the
  training algorithm learns to avoid this.)}
\label{fig:automata-team}
\end{figure}

The clause output is then:
\begin{equation}
  C = \bigwedge_{\substack{k : \text{automaton}_k \\ \text{says include}}} \ell_k
  \label{eq:clause-automata}
\end{equation}

The training algorithm (\cref{ch:clauses}) will adjust the automata
states so that useful literals are included and irrelevant ones are
excluded.


% ----------------------------------------------------------------------------
\section{The Automaton as a Memory Device}
\label{sec:memory}
% ----------------------------------------------------------------------------

It is useful to think of the automaton's state as a form of memory
that accumulates evidence. Each reward strengthens the current
decision; each penalty weakens it. The ``depth'' of memory (controlled
by $N$) determines how much evidence is needed before the automaton
changes its mind.

This is fundamentally different from gradient-based learning:
\begin{itemize}
  \item There is no gradient computation---just integer
    increment/decrement.
  \item There is no learning rate to tune---the step size is always
    $\pm 1$ state.
  \item There is no loss landscape to navigate---the automaton simply
    drifts in the direction that gets more rewards.
\end{itemize}

\begin{keyinsight}
The Tsetlin Automaton replaces all of gradient descent's
machinery---differentiation, learning rates, optimisers---with a
single operation: increment or decrement an integer. This extreme
simplicity is what makes the Tsetlin Machine suitable for
hardware implementation and embedded deployment~\cite{wheeldon2020pervasive}.
\end{keyinsight}


% ----------------------------------------------------------------------------
\section{Stochastic Feedback and Exploration}
\label{sec:stochastic}
% ----------------------------------------------------------------------------

In the Tsetlin Machine, feedback is often made
\concept{stochastic}---that is, the reward or penalty is applied with
some probability rather than deterministically. This serves two
purposes:

\begin{enumerate}
  \item \textbf{Exploration:} Even when a clause is performing well,
    random penalties occasionally push automata toward the boundary,
    allowing the clause to ``try out'' different literal combinations.
  \item \textbf{Diversity:} If all clauses received identical
    deterministic feedback, they would converge to the same pattern.
    Stochastic feedback introduces variation, ensuring that different
    clauses discover different patterns.
\end{enumerate}

The probability of feedback is controlled by a parameter $s$, called
the \concept{specificity} parameter. We will define it precisely in
\cref{ch:clauses}, but the intuition is:

\begin{itemize}
  \item Large $s$ $\Rightarrow$ high probability of reinforcing
    included literals $\Rightarrow$ clauses become more specific
    (include more literals, match fewer inputs).
  \item Small $s$ $\Rightarrow$ lower reinforcement probability
    $\Rightarrow$ clauses stay general (fewer literals, match more
    inputs).
\end{itemize}

\begin{engineernote}
The specificity parameter $s$ is the single most important
hyperparameter of the Tsetlin Machine. Typical values for MNIST are
$s = 3.0$ to $10.0$. A good starting point is $s = 5.0$.
\end{engineernote}


% ----------------------------------------------------------------------------
\section{Summary}
\label{sec:ch5-summary}
% ----------------------------------------------------------------------------

The Tsetlin Automaton is a finite-state machine with a simple update
rule: rewards push the state toward confidence, penalties push it
toward uncertainty. Its key properties are:

\begin{enumerate}
  \item \textbf{Noise tolerance:} The $N$-state memory buffer filters
    out random fluctuations.
  \item \textbf{Guaranteed convergence:} With enough states, the
    automaton provably finds the better action.
  \item \textbf{No calculus:} Learning requires only integer
    increment and decrement.
  \item \textbf{Composability:} Teams of automata can collectively
    learn complex pattern detectors (clauses).
\end{enumerate}

In the next chapter, we will see how the Tsetlin Machine orchestrates
hundreds of these automata, organised into clauses, and trains them
using a clever feedback scheme called Type~I and Type~II feedback.

% ============================================================================
\chapter{Tsetlin Machine Clause Learning}
\label{ch:clauses}
% ============================================================================

This is the core chapter of the Tsetlin Machine. It explains how
clauses are organised into a classifier, how feedback drives learning,
and how the entire system learns to recognise patterns. Read this
chapter carefully---the mechanisms here are the foundation for everything
that follows.


% ----------------------------------------------------------------------------
\section{Architecture Overview}
\label{sec:tm-architecture}
% ----------------------------------------------------------------------------

A Tsetlin Machine~\cite{granmo2018tsetlin} for $K$-class classification (in our case $K = 10$
for digits 0--9) is organised as follows:

\begin{enumerate}
  \item Each class $k \in \{0, 1, \ldots, K{-}1\}$ has $m$ clauses.
    Half are designated \emph{positive} ($C_1^+, C_2^+, \ldots,
    C_{m/2}^+$) and half \emph{negative} ($C_1^-, C_2^-, \ldots,
    C_{m/2}^-$).
  \item Each clause is a conjunctive formula (an AND of selected
    literals), controlled by a team of Tsetlin Automata as described
    in \cref{ch:automaton}.
  \item The \concept{clause vote} for class $k$ is:
    \begin{equation}
      v_k = \sum_{j=1}^{m/2} C_j^+ - \sum_{j=1}^{m/2} C_j^-
      \label{eq:clause-vote}
    \end{equation}
    Positive clauses vote \emph{for} the class; negative clauses
    vote \emph{against}.
  \item The predicted class is the one with the highest vote:
    \begin{equation}
      \hat{y} = \arg\max_k \; v_k
      \label{eq:predict}
    \end{equation}
\end{enumerate}

\begin{figure}[htbp]
\centering
\begin{tikzpicture}[
  clause/.style={draw, rounded corners, minimum width=1.5cm,
    minimum height=0.6cm, font=\scriptsize\sffamily, inner sep=2pt},
  posclause/.style={clause, fill=bitblue!15},
  negclause/.style={clause, fill=bitorange!15},
  >=Stealth
]
  % Input
  \node[font=\small\sffamily\bfseries] at (-2, 0) {Input $\bvec{x}$};

  % Class 0 clauses
  \node[font=\small\sffamily\bfseries] at (3, 2.5) {Class 0};
  \node[posclause] (c0p1) at (1.5, 1.8) {$C_1^+$};
  \node[posclause] (c0p2) at (3.0, 1.8) {$C_2^+$};
  \node[negclause] (c0n1) at (4.5, 1.8) {$C_1^-$};
  \node[negclause] (c0n2) at (6.0, 1.8) {$C_2^-$};

  % Vote
  \node[draw, circle, minimum size=0.8cm, font=\small\sffamily,
    fill=bitgreen!10] (v0) at (8, 1.8) {$v_0$};

  \draw[->, thick] (c0p1) -- (v0) node[midway, above, font=\tiny] {$+$};
  \draw[->, thick] (c0p2) -- (v0);
  \draw[->, thick] (c0n1) -- (v0) node[midway, below, font=\tiny] {$-$};
  \draw[->, thick] (c0n2) -- (v0);

  % Class 1 (abbreviated)
  \node[font=\small\sffamily\bfseries] at (3, 0.3) {Class 1};
  \node[font=\scriptsize\sffamily, bitgray] at (3.75, -0.3) {$\cdots$ (same structure)};
  \node[draw, circle, minimum size=0.8cm, font=\small\sffamily,
    fill=bitgreen!10] (v1) at (8, -0.3) {$v_1$};

  % Dots
  \node[font=\sffamily] at (3, -1.2) {$\vdots$};
  \node[font=\sffamily] at (8, -1.2) {$\vdots$};

  % Class 9
  \node[font=\small\sffamily\bfseries] at (3, -2.1) {Class 9};
  \node[font=\scriptsize\sffamily, bitgray] at (3.75, -2.7) {$\cdots$};
  \node[draw, circle, minimum size=0.8cm, font=\small\sffamily,
    fill=bitgreen!10] (v9) at (8, -2.7) {$v_9$};

  % Argmax
  \node[draw, rounded corners, fill=bitpurple!10, minimum width=1.5cm,
    font=\small\sffamily\bfseries] (argmax) at (10.5, -0.3)
    {$\arg\max$};
  \draw[->, thick] (v0) -- (argmax);
  \draw[->, thick] (v1) -- (argmax);
  \draw[->, thick] (v9) -- (argmax);

  \node[font=\small\sffamily\bfseries, right] at (11.5, -0.3) {$\hat{y}$};
\end{tikzpicture}
\caption{Tsetlin Machine architecture for 10-class classification.
  Each class has $m$ clauses (half positive, half negative). The
  clause vote is the sum of positive clause outputs minus negative
  clause outputs. The predicted class is the argmax of all votes.}
\label{fig:tm-architecture}
\end{figure}


% ----------------------------------------------------------------------------
\section{Clause Evaluation}
\label{sec:clause-eval}
% ----------------------------------------------------------------------------

Given an input $\bvec{x} \in \{0, 1\}^n$ (after booleanisation and
literal expansion, so $n = 2 \times 784 = 1568$ for MNIST), a clause
with included literal set $I \subseteq \{1, \ldots, n\}$ evaluates as:

\begin{equation}
  C(\bvec{x}) = \bigwedge_{k \in I} \ell_k(\bvec{x})
  \label{eq:clause-eval}
\end{equation}

If $I = \emptyset$ (no literals included), the clause outputs 1 for
all inputs---it is vacuously true. The training algorithm prevents
this degenerate case through Type~I feedback (\cref{sec:type1}).

\begin{engineernote}
Clause evaluation maps directly to bitwise AND on packed bit vectors.
For MNIST's 1568 literals, we need $\lceil 1568 / 64 \rceil = 25$
64-bit words. The clause's ``include mask'' is also 25 words. The
evaluation is:
\begin{enumerate}
  \item Bitwise AND the include mask with the literal vector.
  \item Compare the result to the include mask.
  \item If they are equal (all included literals are 1), the clause
    outputs 1.
\end{enumerate}
This takes roughly 50 instructions regardless of how many literals
are included.
\end{engineernote}

In practice, the evaluation is often written as:
\begin{equation}
  C(\bvec{x}) = \begin{cases}
    1 & \text{if } (\bvec{m} \wedge \bvec{x}) = \bvec{m} \\
    0 & \text{otherwise}
  \end{cases}
  \label{eq:clause-bitmask}
\end{equation}
where $\bvec{m}$ is the clause's bitmask with 1s at positions of
included literals.


% ----------------------------------------------------------------------------
\section{The Voting Mechanism and Thresholding}
\label{sec:voting}
% ----------------------------------------------------------------------------

The raw clause vote $v_k$ from \cref{eq:clause-vote} is typically
\concept{clipped} to a threshold $T$:

\begin{equation}
  v_k^{\text{clipped}} = \max(-T, \; \min(T, \; v_k))
  \label{eq:clipped-vote}
\end{equation}

The threshold $T$ limits how strongly any single class can dominate.
It also plays a role in feedback: training feedback is only generated
when $|v_k| < T$, so clauses stop being reinforced once the class
vote is already strong enough. This prevents over-learning.

\begin{keyinsight}
The threshold $T$ controls the number of clauses that actively
participate in learning. If $v_k$ already exceeds $T$, additional
positive clauses are not reinforced---this frees up ``training
capacity'' for other patterns. Think of it as a budget: each class
can spend at most $T$ votes of confidence, encouraging diverse
clause patterns rather than redundant ones.
\end{keyinsight}


% ----------------------------------------------------------------------------
\section{The Feedback Mechanism}
\label{sec:feedback}
% ----------------------------------------------------------------------------

Training a Tsetlin Machine means adjusting the states of all the
automata so that clauses learn to detect useful patterns. The feedback
mechanism is the heart of the algorithm. It comes in two types.

Given a training example $(\bvec{x}, y)$ where $y$ is the true class:

\begin{description}[leftmargin=1em]
  \item[Type~I feedback] is given to clauses of the \textbf{correct
    class} $y$. Its purpose is to make these clauses \emph{more
    specific}---to include more literals that match the current input
    and exclude literals that do not.

  \item[Type~II feedback] is given to clauses of an
    \textbf{incorrect class} $k \neq y$ (specifically, the class
    with the highest or a randomly chosen high vote). Its purpose is
    to make these clauses \emph{less matching}---to push them away
    from recognising the current input.
\end{description}


% ----------------------------------------------------------------------------
\section{Type~I Feedback: Reinforcing the Correct Class}
\label{sec:type1}
% ----------------------------------------------------------------------------

Type~I feedback is given to clauses belonging to the true class $y$.
It operates differently on positive and negative clauses, and
differently depending on whether the clause matched the current input.

\subsection{Type~I on Positive Clauses ($C^+$)}

The goal is to strengthen positive clauses that match patterns from
class $y$.

\paragraph{Case 1: The clause output is 1 (clause matches the input).}

For each literal $\ell_k$ in the clause:
\begin{itemize}
  \item If $\ell_k = 1$ (the literal is true in this input) and the
    automaton says ``include'': \textbf{reward} the automaton with
    probability $\frac{s-1}{s}$. This reinforces the inclusion of
    matching literals.
  \item If $\ell_k = 0$ (the literal is false in this input) and the
    automaton says ``include'': this means the clause includes a
    literal that is currently 0, but the clause still output 1. This
    cannot actually happen (the clause is an AND of included
    literals, so all included literals must be 1 for the clause to
    output 1). So this case is impossible.
  \item If $\ell_k = 1$ and the automaton says ``exclude'':
    \textbf{penalty} the automaton with probability $\frac{s-1}{s}$.
    This pushes excluded-but-true literals toward inclusion,
    making the clause more specific.
  \item If $\ell_k = 0$ and the automaton says ``exclude'':
    \textbf{reward} the automaton with probability $\frac{1}{s}$.
    This gently reinforces the exclusion of currently-false literals.
\end{itemize}

\paragraph{Case 2: The clause output is 0 (clause does not match the input).}

For each literal:
\begin{itemize}
  \item If the automaton says ``include'': \textbf{penalty} with
    probability $\frac{1}{s}$. This gently pushes included literals
    toward exclusion, making the clause less specific (so it might
    match more inputs in the future).
  \item If the automaton says ``exclude'': \textbf{reward} with
    probability $\frac{1}{s}$. This gently reinforces exclusion.
\end{itemize}

\begin{engineernote}
Notice the asymmetry: when the clause matches ($C = 1$), feedback is
applied with high probability $\frac{s-1}{s}$. When the clause does
not match ($C = 0$), feedback uses the low probability $\frac{1}{s}$.
This means matching clauses are strongly reinforced while non-matching
clauses drift slowly. The specificity parameter $s$ controls this
ratio.
\end{engineernote}

\subsection{Type~I on Negative Clauses ($C^-$)}

Negative clauses vote \emph{against} the class. For the true class
$y$, we want negative clauses to \emph{not match}. So Type~I feedback
on negative clauses is the mirror image: it tries to push them away
from the current input pattern.

In practice, Type~I feedback on negative clauses for the true class
is either omitted or applies only the ``clause output is 0'' case
(gently reinforcing exclusion). The key training signal for negative
clauses comes from Type~II feedback on incorrect classes.


% ----------------------------------------------------------------------------
\section{Type~II Feedback: Weakening Incorrect Classes}
\label{sec:type2}
% ----------------------------------------------------------------------------

Type~II feedback targets clauses of an \textbf{incorrect} class
$k \neq y$ that have high votes for the current input. The goal is to
break these clauses so they stop matching.

\paragraph{When does Type~II apply?}
For a positive clause $C^+$ of class $k \neq y$ that outputs~1 (it
wrongly matches the input):

For each literal $\ell_k$ that is \textbf{included} in the clause:
\begin{itemize}
  \item If $\ell_k = 0$ (the literal is false in this input): this
    cannot happen (clause output is 1, so all included literals are
    true).
  \item If $\ell_k = 1$ (the literal is true): \textbf{penalty} the
    automaton, pushing it from ``include'' toward ``exclude.''
\end{itemize}

For each literal that is \textbf{excluded}:
\begin{itemize}
  \item If $\ell_k = 0$ (the literal is false): \textbf{reward} the
    automaton (reinforcing exclusion of this false literal---if it
    were included, it would break the clause, which is what we want,
    but we also want the automaton to ``remember'' this exclusion).
\end{itemize}

\begin{keyinsight}
Type~II feedback has a specific and elegant effect: it tries to
include literals that are 0 in the current input. If it succeeds in
including even one such literal, the clause will output 0 for this
input (because the AND of a 0 is 0), breaking the false match. This
is a \emph{targeted sabotage}: instead of destroying the entire
clause, Type~II just adds one ``poison pill'' literal that prevents
the clause from matching this particular wrong input.
\end{keyinsight}

\begin{figure}[htbp]
\centering
\begin{tikzpicture}[
  >=Stealth,
  box/.style={draw, rounded corners, minimum width=3cm, minimum height=1cm,
    font=\small\sffamily, align=center}
]
  % Type I
  \node[box, fill=bitblue!10] (t1) at (0, 0)
    {\textbf{Type~I Feedback}\\Correct class $y$};
  \node[font=\scriptsize\sffamily, text width=4cm, align=center]
    at (0, -1.5)
    {Strengthen matching clauses\\Include true literals\\Exclude false literals};

  % Type II
  \node[box, fill=bitorange!10] (t2) at (7, 0)
    {\textbf{Type~II Feedback}\\Incorrect class $k \neq y$};
  \node[font=\scriptsize\sffamily, text width=4cm, align=center]
    at (7, -1.5)
    {Break wrongly matching clauses\\Include a false literal\\to sabotage the clause};

  % Arrows
  \draw[->, thick, bitblue] (-2, 1.5) -- (t1)
    node[midway, left, font=\scriptsize\sffamily] {train example};
  \draw[->, thick, bitorange] (9, 1.5) -- (t2)
    node[midway, right, font=\scriptsize\sffamily] {train example};
\end{tikzpicture}
\caption{The two types of feedback. Type~I reinforces clauses of the
  correct class, making them more specific. Type~II sabotages clauses
  of incorrect classes that wrongly matched the input.}
\label{fig:feedback-types}
\end{figure}


% ----------------------------------------------------------------------------
\section{The Complete Feedback Tables}
\label{sec:feedback-tables}
% ----------------------------------------------------------------------------

For reference, we present the complete feedback rules in tabular form.
Let $C$ be the clause output, $\ell$ the literal value in the current
input, and $a$ the automaton action (0 = exclude, 1 = include).

\subsection{Type~I Feedback (Positive Clause, True Class)}

\begin{table}[htbp]
\centering
\caption{Type~I feedback for a positive clause of the true class.
  ``$\nearrow$'' means reward (push away from boundary);
  ``$\searrow$'' means penalty (push toward boundary);
  ``---'' means no feedback. Probabilities are shown in parentheses.}
\label{tab:type1-pos}
\begin{tabular}{@{}ccc|l@{}}
  \toprule
  $C$ & $\ell$ & $a$ & Feedback \\
  \midrule
  1 & 1 & 1 (include) & $\nearrow$ Reward with prob.\ $\frac{s-1}{s}$ \\
  1 & 1 & 0 (exclude) & $\searrow$ Penalty with prob.\ $\frac{s-1}{s}$ \\
  1 & 0 & 1 (include) & impossible ($C = 1$ requires all included $\ell = 1$) \\
  1 & 0 & 0 (exclude) & $\nearrow$ Reward with prob.\ $\frac{1}{s}$ \\
  \midrule
  0 & 1 & 1 (include) & $\searrow$ Penalty with prob.\ $\frac{1}{s}$ \\
  0 & 1 & 0 (exclude) & $\nearrow$ Reward with prob.\ $\frac{1}{s}$ \\
  0 & 0 & 1 (include) & $\searrow$ Penalty with prob.\ $\frac{1}{s}$ \\
  0 & 0 & 0 (exclude) & $\nearrow$ Reward with prob.\ $\frac{1}{s}$ \\
  \bottomrule
\end{tabular}
\end{table}

\subsection{Type~II Feedback (Positive Clause, Wrong Class)}

\begin{table}[htbp]
\centering
\caption{Type~II feedback for a positive clause of an incorrect class.
  Only clauses with $C = 1$ receive Type~II feedback.}
\label{tab:type2-pos}
\begin{tabular}{@{}ccc|l@{}}
  \toprule
  $C$ & $\ell$ & $a$ & Feedback \\
  \midrule
  1 & 1 & 1 (include) & --- (no action) \\
  1 & 1 & 0 (exclude) & $\searrow$ Penalty (push toward include) \\
  1 & 0 & 1 (include) & impossible \\
  1 & 0 & 0 (exclude) & $\searrow$ Penalty (push toward include) \\
  \midrule
  0 & \multicolumn{2}{c|}{any} & --- (no feedback) \\
  \bottomrule
\end{tabular}
\end{table}

\begin{engineernote}
The tables above describe the most common formulation
from~\cite{granmo2018tsetlin}. Variants exist: some
implementations apply Type~II feedback deterministically (probability
1), others use $\frac{1}{s}$. The variant matters less than the
principle: Type~I builds, Type~II sabotages.
\end{engineernote}


% ----------------------------------------------------------------------------
\section{Putting It All Together: One Training Step}
\label{sec:training-step}
% ----------------------------------------------------------------------------

\begin{algorithm}[htbp]
\caption{One training step of the Tsetlin Machine.}
\label{alg:tm-train}
\begin{algorithmic}[1]
  \State \textbf{Input:} Training example $(\bvec{x}, y)$, threshold $T$,
    specificity $s$
  \State Booleanise input: $\bvec{x}_{\text{bool}} \leftarrow
    \text{threshold\_and\_expand}(\bvec{x})$
  \For{each class $k = 0, 1, \ldots, K{-}1$}
    \State Evaluate all clauses: compute $C_j^+(\bvec{x}_{\text{bool}})$
      and $C_j^-(\bvec{x}_{\text{bool}})$
    \State Compute vote: $v_k \leftarrow \sum C_j^+ - \sum C_j^-$
    \State Clip: $v_k \leftarrow \text{clamp}(v_k, -T, T)$
  \EndFor
  \Statex
  \For{each class $k = 0, 1, \ldots, K{-}1$}
    \If{$k = y$} \Comment{True class}
      \State Apply Type~I feedback to clauses of class $k$
      \Statex \quad\quad with probability $\frac{T - v_k}{2T}$
    \Else \Comment{Wrong class}
      \State Apply Type~II feedback to clauses of class $k$
      \Statex \quad\quad with probability $\frac{T + v_k}{2T}$
    \EndIf
  \EndFor
\end{algorithmic}
\end{algorithm}

The probabilities on lines 10 and 13 deserve attention:

\begin{itemize}
  \item For the true class $y$: feedback probability is
    $\frac{T - v_y}{2T}$. When $v_y$ is low (the correct class is
    losing), this probability is high---urgent reinforcement. When
    $v_y \approx T$ (the correct class is winning strongly), the
    probability drops to near zero---no need to train further.

  \item For wrong classes $k \neq y$: feedback probability is
    $\frac{T + v_k}{2T}$. When $v_k$ is high (the wrong class is
    winning), the probability is high---urgent correction. When
    $v_k$ is low or negative, the probability drops.
\end{itemize}

\begin{keyinsight}
The voting margin directly controls the training intensity. Clauses
that are already doing well receive little feedback; clauses that are
struggling receive intense feedback. This is an automatic form of
\emph{hard example mining}---the system focuses its learning effort
where it is most needed.
\end{keyinsight}


% ----------------------------------------------------------------------------
\section{Worked Example: Learning a Simple Pattern}
\label{sec:worked-example}
% ----------------------------------------------------------------------------

Let us trace the training of a single positive clause for class~``1''
on a tiny 4-pixel image. We use $N = 3$ states per action (6 states
total per automaton) and $s = 3.0$.

\subsection{Setup}

The input has 4 pixels, giving 8 literals after expansion:
$x_1, \bar{x}_1, x_2, \bar{x}_2, x_3, \bar{x}_3, x_4, \bar{x}_4$.

We assign one automaton to each literal, all starting in state~4
(just above the boundary, action = include). So initially the clause
includes all 8 literals: $C = x_1 \wedge \bar{x}_1 \wedge \cdots$,
which always outputs 0 (since $x_k \wedge \bar{x}_k = 0$).

\begin{center}
\begin{tabular}{@{}l*{8}{c}@{}}
  \toprule
  \textbf{Literal} & $x_1$ & $\bar{x}_1$ & $x_2$ & $\bar{x}_2$ &
    $x_3$ & $\bar{x}_3$ & $x_4$ & $\bar{x}_4$ \\
  \midrule
  Initial state & 4 & 4 & 4 & 4 & 4 & 4 & 4 & 4 \\
  Action & inc & inc & inc & inc & inc & inc & inc & inc \\
  \bottomrule
\end{tabular}
\end{center}

\subsection{Training Example 1}

Input: image of ``1'' with pixels $(1, 0, 1, 0)$, so:
\[
  (x_1, \bar{x}_1, x_2, \bar{x}_2, x_3, \bar{x}_3, x_4, \bar{x}_4)
  = (1, 0, 0, 1, 1, 0, 0, 1)
\]

The clause output is 0 (since many included literals are 0). Since
$C = 0$ and this is a positive clause for the true class, Type~I
feedback applies with the ``clause output = 0'' rules.

Each automaton receives a penalty with probability $\frac{1}{s} =
\frac{1}{3} \approx 0.33$. Suppose random draws penalise automata
for $\bar{x}_1$, $x_2$, $\bar{x}_3$, and $x_4$. These automata
move from state~4 to state~3 (crossing the boundary from include to
exclude):

\begin{center}
\begin{tabular}{@{}l*{8}{c}@{}}
  \toprule
  \textbf{Literal} & $x_1$ & $\bar{x}_1$ & $x_2$ & $\bar{x}_2$ &
    $x_3$ & $\bar{x}_3$ & $x_4$ & $\bar{x}_4$ \\
  \midrule
  State & 4 & \textbf{3} & \textbf{3} & 4 & 4 & \textbf{3} & \textbf{3} & 4 \\
  Action & inc & \textbf{exc} & \textbf{exc} & inc & inc & \textbf{exc} & \textbf{exc} & inc \\
  \bottomrule
\end{tabular}
\end{center}

Now the clause is $C = x_1 \wedge \bar{x}_2 \wedge x_3 \wedge
\bar{x}_4$. Let us check: for the input $(1, 0, 1, 0)$, the literal
values at included positions are $x_1 = 1$, $\bar{x}_2 = 1$,
$x_3 = 1$, $\bar{x}_4 = 1$. All are 1, so $C = 1$. The clause
now matches this ``1'' pattern.

\subsection{Training Example 2}

Input: image of ``0'' with pixels $(1, 1, 1, 1)$, so all positive
literals are 1 and all negated literals are 0. The clause evaluates:
$x_1 = 1$, $\bar{x}_2 = 0$---stop, $C = 0$. Good: the clause does
not match this ``0'' image.

But now suppose another ``0'' image arrives with pixels $(1, 0, 1, 1)$:
$x_1 = 1$, $\bar{x}_2 = 1$, $x_3 = 1$, $\bar{x}_4 = 0$---stop,
$C = 0$. Still correct.

And a ``1'' image with pixels $(1, 0, 0, 0)$:
$x_1 = 1$, $\bar{x}_2 = 1$, $x_3 = 0$---stop, $C = 0$. This ``1''
pattern does not match because $x_3 = 0$. This clause has learned a
\emph{specific} pattern for ``1'' that requires both $x_1$ and $x_3$
to be bright.

\subsection{Continuing Training}

Over many examples, Type~I feedback will:
\begin{itemize}
  \item Strongly reinforce included literals that are consistently
    true for class ``1'' images (pushing automata toward state~6).
  \item Gently push excluded literals further into exclusion (toward
    state~1).
\end{itemize}

If a Type~II signal arrives (because this clause wrongly matches an
image from another class), it will try to include a literal that is
0 in that image, preventing the false match.


% ----------------------------------------------------------------------------
\section{The Role of Multiple Clauses}
\label{sec:multiple-clauses}
% ----------------------------------------------------------------------------

A single clause can only detect one specific pattern. The digit ``3''
can be written in many ways---some with a flat top, some with a
rounded top, some thick, some thin. No single conjunctive clause can
match all variants.

This is why each class has $m$ clauses (typically $m = 100$ to
$2000$). Each clause specialises in a different sub-pattern:

\begin{itemize}
  \item Clause~1 might detect ``3'' with a flat top bar.
  \item Clause~2 might detect ``3'' with a rounded top.
  \item Clause~3 might detect a common stroke in the middle.
  \item And so on.
\end{itemize}

The vote aggregation (sum of clause outputs) acts as a soft OR: if
\emph{any} clause matches, the vote for class ``3'' increases.

\begin{analogy}
Imagine a committee of experts, each looking at a different aspect of
the digit. Expert~1 checks if the top stroke looks right. Expert~2
checks the middle junction. Expert~3 checks the bottom curve. They
vote, and the digit is classified based on the majority opinion. The
Tsetlin Machine learns both the experts' criteria (which pixels to
check) and which experts to trust (positive vs.\ negative polarity).
\end{analogy}


% ----------------------------------------------------------------------------
\section{Negative Clauses: Learning What a Digit Is Not}
\label{sec:negative-clauses}
% ----------------------------------------------------------------------------

Negative clauses vote \emph{against} their class. A negative clause
for class ``3'' might detect patterns that look like ``8'' (which is
similar to ``3'' but has a closed top loop). When this pattern is
present, the negative clause fires and reduces the vote for ``3,''
making it less likely to be confused with ``8.''

\begin{example}
Suppose class ``7'' has:
\begin{itemize}
  \item A positive clause that fires when there is a horizontal top
    stroke and a diagonal down stroke.
  \item A negative clause that fires when there is also a horizontal
    middle stroke (suggesting ``1'' written with a crossbar, or
    some styles of ``4'').
\end{itemize}
The positive clause says ``this might be a 7'' and the negative clause
says ``but this feature argues against 7.'' The vote is the net
effect.
\end{example}


% ----------------------------------------------------------------------------
\section{Hyperparameters and Their Effects}
\label{sec:hyperparams}
% ----------------------------------------------------------------------------

The Tsetlin Machine has remarkably few hyperparameters:

\begin{table}[htbp]
\centering
\caption{Tsetlin Machine hyperparameters.}
\label{tab:tm-hyperparams}
\begin{tabular}{@{}llp{6cm}@{}}
  \toprule
  \textbf{Symbol} & \textbf{Typical range} & \textbf{Effect} \\
  \midrule
  $m$ & 100--2000 & Number of clauses per class. More clauses
    $\Rightarrow$ more capacity to learn diverse patterns, but
    slower training and inference. \\
  $T$ & 10--100 & Vote threshold. Controls how many clauses
    actively train. Lower $T$ $\Rightarrow$ fewer active clauses,
    faster convergence but potentially lower accuracy. \\
  $s$ & 2.0--10.0 & Specificity. Higher $s$ $\Rightarrow$ more
    specific clauses (more included literals). Lower $s$ $\Rightarrow$
    more general clauses. \\
  $N$ & 50--200 & States per action in each automaton. Higher $N$
    $\Rightarrow$ more noise resistance, slower adaptation. \\
  \bottomrule
\end{tabular}
\end{table}

\begin{engineernote}
A good starting configuration for MNIST with simple thresholding:
$m = 500$ clauses per class (250 positive, 250 negative), $T = 50$,
$s = 5.0$, $N = 100$. This achieves approximately 96--97\% test
accuracy and trains in under a minute on a modern CPU.
\end{engineernote}


% ----------------------------------------------------------------------------
\section{Why Stochastic Feedback Creates Diversity}
\label{sec:diversity}
% ----------------------------------------------------------------------------

A crucial question: if all clauses for class ``3'' receive the same
training examples, why do they learn different patterns?

The answer is the stochastic feedback. Each automaton independently
receives or does not receive a reward/penalty based on random number
generation. This means:

\begin{enumerate}
  \item Two clauses seeing the same input receive different feedback
    on different literals.
  \item Over time, random variations cause clauses to diverge: one
    clause might include $x_{42}$ while another excludes it, leading
    them to specialise on different sub-patterns.
  \item The vote threshold $T$ amplifies this effect: once enough
    clauses are matching, others stop receiving feedback and are free
    to explore alternative patterns.
\end{enumerate}

\begin{keyinsight}
Diversity in the Tsetlin Machine is not engineered---it is an emergent
property of stochastic feedback combined with vote thresholding. Each
clause follows its own random walk through literal space, and the
aggregate system self-organises into a diverse committee of pattern
detectors.
\end{keyinsight}


% ----------------------------------------------------------------------------
\section{Comparison with Gradient-Based Learning}
\label{sec:tm-vs-gradient}
% ----------------------------------------------------------------------------

It is instructive to compare the Tsetlin Machine training loop with
standard gradient descent:

\begin{table}[htbp]
\centering
\caption{Tsetlin Machine vs.\ gradient descent.}
\label{tab:tm-vs-gd}
\begin{tabular}{@{}lll@{}}
  \toprule
  \textbf{Aspect} & \textbf{Gradient descent} & \textbf{Tsetlin Machine} \\
  \midrule
  Parameter type & Floating-point weights & Integer automaton states \\
  Update rule & $w \leftarrow w - \eta \nabla$ & $s \leftarrow s \pm 1$ \\
  Requires & Differentiable ops & Any reward/penalty signal \\
  Arithmetic & FP multiply-add & Integer increment/decrement \\
  Hyperparams & Learning rate, optimiser, & $s$, $T$, $m$, $N$ \\
   & momentum, weight decay & \\
  Memory & 32 bits per weight & $\lceil \log_2(2N) \rceil$ bits per automaton \\
  Interpretability & Opaque weight matrices & Readable Boolean clauses \\
  \bottomrule
\end{tabular}
\end{table}

The Tsetlin Machine's main advantages are simplicity, interpretability,
and integer-only arithmetic~\cite{wheeldon2020pervasive}. Its main disadvantage is slower convergence
on complex tasks~\cite{zhang2021convergence}, which is partly why we combine it with a Logic Gate
Network~\cite{petersen2022difflogic} in the hybrid architecture.


% ----------------------------------------------------------------------------
\section{Summary}
\label{sec:ch6-summary}
% ----------------------------------------------------------------------------

The Tsetlin Machine training algorithm can be summarised in four
principles:

\begin{enumerate}
  \item \textbf{Type~I reinforces:} For the correct class, strengthen
    clauses that match the input by rewarding included-and-true
    literals and penalising excluded-but-true ones.
  \item \textbf{Type~II sabotages:} For incorrect classes, break
    wrongly-matching clauses by including literals that are false
    in the current input.
  \item \textbf{Stochastic diversity:} Random feedback probabilities
    ensure different clauses learn different patterns.
  \item \textbf{Vote thresholding:} The threshold $T$ balances
    training effort across clauses, preventing any single pattern
    from monopolising the vote.
\end{enumerate}

The next chapter applies these principles to MNIST and walks through
the training process on real digit images.

% ============================================================================
\chapter{Tsetlin Machines on MNIST}
\label{ch:tm-mnist}
% ============================================================================

This chapter applies the Tsetlin Machine~\cite{granmo2018tsetlin} to MNIST~\cite{lecun1998mnist} digit recognition.
We walk through the complete pipeline from raw pixels to classification,
examine what the learned clauses look like, and discuss the performance
and accuracy one can expect.


% ----------------------------------------------------------------------------
\section{The End-to-End Pipeline}
\label{sec:tm-pipeline}
% ----------------------------------------------------------------------------

The complete inference pipeline for a single image:

\begin{enumerate}
  \item \textbf{Load image:} Read the $28 \times 28$ pixel values
    $p_1, \ldots, p_{784} \in [0, 255]$.
  \item \textbf{Booleanise:} Apply threshold at $\theta = 127$:
    $x_k = \mathbf{1}[p_k > 127]$.
  \item \textbf{Expand literals:} Create
    $\bvec{x}_{\text{exp}} = (x_1, \bar{x}_1, \ldots, x_{784},
    \bar{x}_{784}) \in \{0,1\}^{1568}$.
  \item \textbf{Evaluate clauses:} For each class $k$ and each
    clause $C_j^+, C_j^-$, compute $C(\bvec{x}_{\text{exp}})$ using
    bitwise AND.
  \item \textbf{Sum votes:} $v_k = \sum C_j^+ - \sum C_j^-$, clipped
    to $[-T, T]$.
  \item \textbf{Classify:} $\hat{y} = \arg\max_k \, v_k$.
\end{enumerate}

\begin{figure}[htbp]
\centering
\begin{tikzpicture}[
  box/.style={draw, rounded corners, minimum width=2.2cm,
    minimum height=1cm, font=\small\sffamily, align=center,
    line width=0.6pt},
  arr/.style={-{Stealth[length=4pt]}, thick},
]
  \node[box, fill=bitgray!10] (img) at (0, 0) {$28 \times 28$\\pixels};
  \node[box, fill=bitorange!10] (bool) at (3.5, 0) {Booleanise\\$\theta = 127$};
  \node[box, fill=bitorange!10] (exp) at (7, 0) {Expand\\literals};
  \node[box, fill=bitblue!10] (eval) at (10.5, 0) {Evaluate\\clauses};
  \node[box, fill=bitgreen!10] (vote) at (14, 0) {Vote\\$\arg\max$};

  \draw[arr] (img) -- (bool);
  \draw[arr] (bool) -- (exp) node[midway, above, font=\tiny\sffamily] {784 bits};
  \draw[arr] (exp) -- (eval) node[midway, above, font=\tiny\sffamily] {1568 bits};
  \draw[arr] (eval) -- (vote);

  \node[font=\small\sffamily\bfseries, right] at (15.5, 0) {digit};
\end{tikzpicture}
\caption{The Tsetlin Machine inference pipeline for MNIST.}
\label{fig:tm-pipeline}
\end{figure}


% ----------------------------------------------------------------------------
\section{Computational Cost Analysis}
\label{sec:tm-cost}
% ----------------------------------------------------------------------------

Let us count the operations for a single image classification with
$m = 500$ clauses per class and 10 classes.

\paragraph{Booleanisation and expansion.}
784 comparisons (threshold) plus 784 NOT operations (negation).
Total: roughly 1,568 operations, or 25 64-bit word operations using
bitwise parallel NOT.

\paragraph{Clause evaluation.}
Each clause requires $\lceil 1568 / 64 \rceil = 25$ AND operations
plus a comparison of 25 words. Approximately 50 operations per clause.

With $500 \times 10 = 5{,}000$ total clauses:
$5{,}000 \times 50 = 250{,}000$ bitwise operations.

\paragraph{Vote summation.}
$5{,}000$ additions plus 10 argmax comparisons. Trivial.

\begin{engineernote}
The total inference cost is dominated by clause evaluation: about
250,000 bitwise operations. Compare this to a conventional two-layer
network with 784 inputs, 300 hidden neurons, and 10 outputs:
$(784 \times 300 + 300 \times 10) = 238{,}200$ floating-point
multiply-accumulates. The operation counts are similar, but each
Tsetlin operation is a bitwise AND (0.1~pJ) versus a floating-point
MAC (roughly 4~pJ)~\cite{horowitz2014energy}, giving a roughly 40$\times$ energy advantage per
operation.
\end{engineernote}


% ----------------------------------------------------------------------------
\section{What Do Learned Clauses Look Like?}
\label{sec:learned-clauses}
% ----------------------------------------------------------------------------

After training, each clause is a conjunction of selected literals.
We can visualise a clause by highlighting which pixel positions
are included:

\begin{figure}[htbp]
\centering
\begin{tikzpicture}[scale=0.14]
  % Grid for clause visualization
  \fill[bitgray!5] (0,0) rectangle (28,28);
  \draw[bitgray!15, very thin, step=1] (0,0) grid (28,28);

  % Positive literals (pixel must be bright) - shown in blue
  % Representing a typical clause for digit "7"
  \foreach \x/\y in {
    8/24, 9/24, 10/24, 11/24, 12/24, 13/24, 14/24, 15/24,
    16/24, 17/24, 18/24, 19/24,
    18/23, 19/23,
    17/22,
    16/21,
    15/20,
    14/19, 15/19,
    13/17, 14/17,
    13/16
  }{
    \fill[bitblue!50] (\x,\y) rectangle ++(1,1);
  }

  % Negative literals (pixel must be dark) - shown in orange
  \foreach \x/\y in {
    5/24, 6/24, 7/24,
    20/24, 21/24, 22/24,
    5/20, 6/20, 7/20,
    5/16, 6/16, 7/16,
    5/12, 6/12, 7/12,
    20/12, 21/12, 22/12
  }{
    \fill[bitorange!30] (\x,\y) rectangle ++(1,1);
  }

  % Legend
  \node[font=\small\sffamily] at (14, -2)
    {\textcolor{bitblue}{Blue}: $x_k$ included (pixel must be bright)};
  \node[font=\small\sffamily] at (14, -3.5)
    {\textcolor{bitorange}{Orange}: $\bar{x}_k$ included (pixel must be dark)};
\end{tikzpicture}
\caption{Visualisation of a learned clause for digit ``7.'' Blue pixels
  must be bright (the clause includes $x_k$); orange pixels must be
  dark (the clause includes $\bar{x}_k$). Unshaded pixels are excluded
  and can be anything. The clause detects a horizontal top bar with a
  diagonal downstroke, while requiring the left and lower-right regions
  to be empty.}
\label{fig:clause-vis}
\end{figure}

Key observations about learned clauses:
\begin{itemize}
  \item Clauses typically include 20--80 literals (out of 1568
    available). They are specific enough to detect a pattern but
    general enough to match variations.
  \item Positive literals ($x_k$) mark pixels that must be bright
    (ink). They trace the expected stroke positions.
  \item Negative literals ($\bar{x}_k$) mark pixels that must be
    dark (background). They ensure the surrounding area is empty,
    distinguishing similar digits.
  \item The excluded literals (no constraint) handle the natural
    variation in handwriting---the clause does not care about pixels
    that vary across examples of the same digit.
\end{itemize}

\begin{keyinsight}
Tsetlin Machine clauses are inherently interpretable. Each clause
is a human-readable rule: ``if these pixels are bright AND those
pixels are dark, then this might be digit~$k$.'' No post-hoc
explanation method is needed---the model \emph{is} the explanation.
This is a significant advantage for debugging, auditing, and
understanding what the model has learned.
\end{keyinsight}


% ----------------------------------------------------------------------------
\section{Training Dynamics}
\label{sec:training-dynamics}
% ----------------------------------------------------------------------------

Training a Tsetlin Machine on MNIST typically proceeds as follows:

\paragraph{Epoch 1--2.} Clauses are mostly random. Many include
contradictory literal pairs ($x_k$ and $\bar{x}_k$), so they always
output~0. Test accuracy is around 60--70\%.

\paragraph{Epoch 3--5.} Type~I feedback begins to resolve
contradictions. Clauses start specialising on common patterns (e.g.,
the straight vertical stroke of ``1,'' the closed loop of ``0'').
Test accuracy reaches 85--90\%.

\paragraph{Epoch 5--15.} Clauses refine. Some become very specific
(many literals) and match rare handwriting styles. Others remain
general and match common patterns. Test accuracy reaches 93--96\%.

\paragraph{Epoch 15--50.} Marginal gains. The system learns to
distinguish confusable pairs like (3, 8), (4, 9), and (5, 6).
Test accuracy plateaus at 96--97\% (with simple thresholding)~\cite{abeyrathna2021integer}.

\begin{figure}[htbp]
\centering
\begin{tikzpicture}[>=Stealth]
  \begin{axis}[
    width=10cm, height=6cm,
    xlabel={Epoch}, ylabel={Test accuracy (\%)},
    xmin=0, xmax=50, ymin=50, ymax=100,
    grid=major,
    grid style={bitgray!20},
    every axis label/.style={font=\small\sffamily},
    tick label style={font=\small},
    legend pos=south east,
    legend style={font=\small\sffamily}
  ]
    % Typical learning curve
    \addplot[thick, bitblue, smooth] coordinates {
      (1, 65) (2, 78) (3, 85) (5, 90) (8, 93)
      (10, 94.5) (15, 95.5) (20, 96.2) (30, 96.7)
      (40, 96.9) (50, 97.0)
    };
    \addlegendentry{$m=500$, $s=5.0$}

    % More clauses
    \addplot[thick, bitgreen, smooth, dashed] coordinates {
      (1, 68) (2, 82) (3, 88) (5, 92) (8, 94.5)
      (10, 95.5) (15, 96.5) (20, 97.0) (30, 97.3)
      (40, 97.5) (50, 97.6)
    };
    \addlegendentry{$m=2000$, $s=5.0$}
  \end{axis}
\end{tikzpicture}
\caption{Typical training curves for a Tsetlin Machine on MNIST with
  simple thresholding ($\theta = 127$). More clauses give higher
  accuracy but slower training. The curve shows rapid initial
  improvement followed by gradual refinement.}
\label{fig:training-curve}
\end{figure}


% ----------------------------------------------------------------------------
\section{Accuracy Results and Comparison}
\label{sec:tm-accuracy}
% ----------------------------------------------------------------------------

Published results for Tsetlin Machines on
MNIST~\cite{granmo2018tsetlin, abeyrathna2021integer}:

\begin{table}[htbp]
\centering
\caption{Tsetlin Machine accuracy on MNIST under different
  configurations.}
\label{tab:tm-accuracy}
\begin{tabular}{@{}llcr@{}}
  \toprule
  \textbf{Configuration} & \textbf{Encoding} & \textbf{Clauses/class}
    & \textbf{Test accuracy} \\
  \midrule
  Baseline TM & Threshold ($\theta=127$) & 500 & 96.5\% \\
  Large TM & Threshold ($\theta=127$) & 2000 & 97.2\% \\
  Thermometer TM & Thermometer ($L=8$) & 2000 & 98.2\% \\
  Weighted TM~\cite{phoulady2020weighted}
    & Thermometer ($L=8$) & 2000 & 98.6\% \\
  Convolutional TM~\cite{granmo2019convtm}
    & Patches & 8000 & 99.3\% \\
  \midrule
  2-layer FP network & Continuous & --- & 98.0\% \\
  \bottomrule
\end{tabular}
\end{table}

The baseline Tsetlin Machine with simple thresholding achieves
96--97\%, which is 1--2\% below a standard floating-point network.
With thermometer encoding and more clauses, it matches or exceeds
the FP baseline. The convolutional variant~\cite{granmo2019convtm} reaches 99.3\%
(mean; 99.4\% peak), which is competitive with CNNs.

\begin{engineernote}
For our hybrid architecture, we target the baseline TM configuration
(simple thresholding, 500 clauses/class). This gives 96--97\%
accuracy from the Tsetlin stage alone. The Logic Gate Network and
LUT stages will then refine this further, aiming for $>$97\% total.
\end{engineernote}


% ----------------------------------------------------------------------------
\section{Memory Footprint}
\label{sec:tm-memory}
% ----------------------------------------------------------------------------

Each clause stores a bitmask of 1568~bits (which literals are
included). With 500 clauses per class and 10 classes:

\begin{equation}
  \text{Model size} = 10 \times 500 \times 1568 \text{ bits}
    = 7{,}840{,}000 \text{ bits} = 980 \text{ kB}
  \label{eq:tm-model-size}
\end{equation}

During training, each automaton also stores its state (typically 8
bits for $N \leq 128$), adding another 980~kB. At inference time,
only the bitmasks are needed.

\begin{engineernote}
The 980~kB inference model fits comfortably in the L2 cache of any
modern CPU (typically 256~kB to 8~MB)~\cite{wheeldon2020pervasive}. This means clause evaluation
runs entirely from cache, with zero main memory accesses during
inference. Compare this to a floating-point network where 784 $\times$
300 $\times$ 4~bytes = 940~kB just for the first layer's weights---similar
size but requiring FP arithmetic instead of bitwise operations.
\end{engineernote}


% ----------------------------------------------------------------------------
\section{The Clause Output as Features}
\label{sec:clause-features}
% ----------------------------------------------------------------------------

For the hybrid architecture, the Tsetlin Machine's role is not just
classification but \concept{feature extraction}. The clause outputs
themselves are valuable binary features that capture learned patterns
in the data.

With 500 clauses per class and 10 classes, we have 5000 clause
outputs---each a binary value indicating whether a specific pattern
is present in the current image. These 5000 bits form a rich,
learned feature representation that can be fed into the Logic Gate
Network (\cref{ch:diff-gates}).

\begin{keyinsight}
The Tsetlin Machine transforms a 1568-dimensional literal vector
(raw pixel features) into a 5000-dimensional clause output vector
(learned pattern features). Each dimension answers a specific
yes/no question: ``Does this image contain pattern~$j$?'' This
binary feature representation is the interface between the Tsetlin
stage and the Logic Gate Network stage of the hybrid architecture.
\end{keyinsight}

The clause outputs have several desirable properties as features:
\begin{itemize}
  \item \textbf{Binary:} Perfect for feeding into logic gates.
  \item \textbf{Interpretable:} Each feature corresponds to a
    readable clause.
  \item \textbf{Sparse:} For any given input, most clauses output~0.
    Only the 10--50 clauses that match fire. This sparsity can be
    exploited for efficiency.
  \item \textbf{Diverse:} The stochastic training ensures different
    clauses detect different patterns, giving a rich feature set.
\end{itemize}


% ----------------------------------------------------------------------------
\section{Limitations and What Comes Next}
\label{sec:tm-limitations}
% ----------------------------------------------------------------------------

The standalone Tsetlin Machine has several limitations that motivate
the hybrid approach:

\begin{enumerate}
  \item \textbf{Linear voting:} The clause vote is a simple sum. It
    cannot learn nonlinear combinations of clause outputs. If
    recognising a digit requires ``clause~17 fires AND clause~42 does
    NOT fire,'' the linear vote cannot express this.

  \item \textbf{Independent clauses:} Each clause is trained
    independently (modulo the shared vote threshold). There is no
    mechanism for clauses to explicitly cooperate or specialise.

  \item \textbf{Slow convergence:} The stochastic, one-state-at-a-time
    updates are robust but slow compared to gradient descent. Training
    a large Tsetlin Machine can take 50--200 epochs.

  \item \textbf{No composition:} Clauses operate directly on the
    input features. There is no hierarchy---no ``clause of clauses''
    that builds higher-level concepts from lower-level patterns.
\end{enumerate}

The Logic Gate Network~\cite{petersen2022difflogic} addresses limitations 1 and 4: it learns
nonlinear combinations of clause outputs arranged in a layered
hierarchy. The LUT compilation~\cite{liu2024kan} addresses efficiency: it precomputes
the scoring functions into lookup tables for minimal-cost inference.

In Part~III, we will see how the Logic Gate Network takes the Tsetlin
Machine's clause outputs and learns to combine them through logic
gates, creating a deeper and more expressive classifier while
maintaining the all-Boolean character of the architecture.


% Part III: Logic Gate Networks
\part{Logic Gate Networks --- Learning the Circuit}
\label{part:lgn}

% ============================================================================
\chapter{Differentiable Gate Selection}
\label{ch:diff-gates}
% ============================================================================

This chapter introduces the Logic Gate Network (LGN)~\cite{petersen2022difflogic}, a neural network
where each node is a 2-input logic gate. The key innovation is a
training method that uses gradient descent to \emph{select} which of
the 16 possible gate types each node should implement. After training,
the continuous selection mechanism is removed, leaving a pure Boolean
circuit.


% ----------------------------------------------------------------------------
\section{The 16 Two-Input Boolean Functions}
\label{sec:sixteen-gates}
% ----------------------------------------------------------------------------

As we saw in \cref{sec:truth-tables}, two binary inputs $(a, b)$
have $2^2 = 4$ input combinations, and each combination maps to 0
or~1. This gives $2^4 = 16$ possible functions. We list all 16:

\begin{table}[htbp]
\centering
\caption{All 16 two-input Boolean functions, ordered by their 4-bit
  truth table code. The code lists outputs for inputs $(0,0),
  (0,1), (1,0), (1,1)$.}
\label{tab:all-16}
\begin{tabular}{@{}clcccc@{}}
  \toprule
  \textbf{Index} & \textbf{Name} & $(0,0)$ & $(0,1)$ & $(1,0)$ & $(1,1)$ \\
  \midrule
  0  & FALSE (constant 0) & 0 & 0 & 0 & 0 \\
  1  & AND              & 0 & 0 & 0 & 1 \\
  2  & $A \wedge \bar{B}$ (inhibit B) & 0 & 0 & 1 & 0 \\
  3  & $A$ (project A)    & 0 & 0 & 1 & 1 \\
  4  & $\bar{A} \wedge B$ (inhibit A) & 0 & 1 & 0 & 0 \\
  5  & $B$ (project B)    & 0 & 1 & 0 & 1 \\
  6  & XOR              & 0 & 1 & 1 & 0 \\
  7  & OR               & 0 & 1 & 1 & 1 \\
  8  & NOR              & 1 & 0 & 0 & 0 \\
  9  & XNOR             & 1 & 0 & 0 & 1 \\
  10 & $\bar{B}$ (NOT B)  & 1 & 0 & 1 & 0 \\
  11 & $A \vee \bar{B}$ (imply A) & 1 & 0 & 1 & 1 \\
  12 & $\bar{A}$ (NOT A)  & 1 & 1 & 0 & 0 \\
  13 & $\bar{A} \vee B$ (imply B) & 1 & 1 & 0 & 1 \\
  14 & NAND             & 1 & 1 & 1 & 0 \\
  15 & TRUE (constant 1)  & 1 & 1 & 1 & 1 \\
  \bottomrule
\end{tabular}
\end{table}

Each function is uniquely identified by its 4-bit truth table code.
We denote the truth table of gate type $g$ as $\TruthTable(g) \in
\{0,1\}^4$, indexed by the input combination.


% ----------------------------------------------------------------------------
\section{The Gate Selection Problem}
\label{sec:gate-selection}
% ----------------------------------------------------------------------------

The Logic Gate Network~\cite{petersen2022difflogic,petersen2024difflogic} consists of layers of nodes. Each node takes
two binary inputs and produces one binary output. The question is:
which of the 16 gate types should each node implement?

If we fix the gate types, evaluation is trivial---just apply the logic
operations. The challenge is \emph{learning} the gate types from data.

\begin{keyinsight}
A logic gate is a discrete object: it is AND or it is OR; there is no
``halfway between AND and OR.'' Gradient descent cannot directly
optimise a discrete choice. The LGN solves this by replacing the
discrete choice with a continuous probability distribution over gate
types, optimising this distribution with gradients, and then
discretising after training.
\end{keyinsight}


% ----------------------------------------------------------------------------
\section{The Continuous Relaxation}
\label{sec:relaxation}
% ----------------------------------------------------------------------------

For each node $i$ in the network, we maintain a vector of 16
\concept{logits}:
\begin{equation}
  \bvec{z}_i = (z_i^{(0)}, z_i^{(1)}, \ldots, z_i^{(15)})
    \in \mathbb{R}^{16}
  \label{eq:logits}
\end{equation}

These logits are converted to a probability distribution using the
\concept{softmax} function:
\begin{equation}
  p_i^{(g)} = \frac{\exp(z_i^{(g)} / \tau)}
    {\sum_{g'=0}^{15} \exp(z_i^{(g')} / \tau)}
  \label{eq:softmax-gate}
\end{equation}

where $\tau > 0$ is the \concept{temperature} parameter.

The node's output is then a \emph{weighted average} of all 16 gate
outputs:
\begin{equation}
  \hat{y}_i = \sum_{g=0}^{15} p_i^{(g)} \cdot \TruthTable(g)[a_i, b_i]
  \label{eq:relaxed-output}
\end{equation}

where $a_i, b_i \in \{0, 1\}$ are the node's two inputs, and
$\TruthTable(g)[a_i, b_i]$ is the output of gate type $g$ on those
inputs.

\begin{example}
Suppose a node has inputs $a = 1, b = 0$. The truth table outputs for
this input combination are: AND gives 0, OR gives 1, XOR gives 1,
NAND gives 1, etc. If the probabilities are
$p^{(\text{AND})} = 0.6$ and $p^{(\text{OR})} = 0.3$ and
$p^{(\text{others})} = 0.1$, then:
\[
  \hat{y} = 0.6 \times 0 + 0.3 \times 1 + 0.1 \times (\text{others})
  = 0.3 + \cdots \approx 0.4
\]
The output is continuous (between 0 and 1), not binary. This is the
price of differentiability.
\end{example}


% ----------------------------------------------------------------------------
\section{Temperature Annealing}
\label{sec:annealing}
% ----------------------------------------------------------------------------

The temperature $\tau$ controls how ``soft'' the gate selection is:

\begin{itemize}
  \item $\tau \to \infty$: All gate types have equal probability
    ($p^{(g)} = 1/16$). The output is a uniform average. Maximum
    exploration, minimum commitment.
  \item $\tau \to 0$: The probability concentrates on a single gate
    type (the one with the largest logit). The output approaches a
    hard Boolean value. Maximum commitment, zero exploration.
\end{itemize}

During training, we anneal $\tau$ from a high value to a low value:

\begin{equation}
  \tau(t) = \tau_{\text{start}} \cdot
    \left(\frac{\tau_{\text{end}}}{\tau_{\text{start}}}\right)^{t / t_{\max}}
  \label{eq:anneal}
\end{equation}

where $t$ is the training step and $t_{\max}$ is the total number of
steps.

\begin{figure}[htbp]
\centering
\begin{tikzpicture}[>=Stealth]
  \begin{axis}[
    width=10cm, height=5cm,
    xlabel={Training progress (\%)}, ylabel={Temperature $\tau$},
    xmin=0, xmax=100, ymin=0, ymax=5,
    grid=major, grid style={bitgray!20},
    every axis label/.style={font=\small\sffamily},
    tick label style={font=\small}
  ]
    \addplot[thick, bitblue, smooth, domain=0:100, samples=50]
      {5 * (0.01/5)^(x/100)};

    \node[font=\scriptsize\sffamily, bitblue] at (axis cs:15, 4)
      {explore};
    \node[font=\scriptsize\sffamily, bitblue] at (axis cs:85, 0.8)
      {commit};
  \end{axis}
\end{tikzpicture}
\caption{Temperature annealing schedule. The temperature starts high
  (soft selection, exploration) and decays to near zero (hard
  selection, commitment). The exponential decay ensures a gradual
  transition.}
\label{fig:annealing}
\end{figure}

\begin{analogy}
Temperature annealing is like a hiring decision. Early in the process,
you keep many candidates under consideration (high temperature). As
interviews proceed, you narrow the field. Eventually, you commit to
one candidate (low temperature). The gradual narrowing avoids premature
decisions based on limited information.
\end{analogy}


% ----------------------------------------------------------------------------
\section{Why This Works: The Gradient Flow}
\label{sec:gradient-flow}
% ----------------------------------------------------------------------------

The relaxed output (\cref{eq:relaxed-output}) is differentiable with
respect to the logits $\bvec{z}_i$. The gradient of the loss with
respect to logit $z_i^{(g)}$ tells us: ``if we increase the
probability of gate type $g$ at node $i$, does the loss go up or
down?''

Specifically, if $\mathcal{L}$ is the loss:
\begin{equation}
  \frac{\partial \mathcal{L}}{\partial z_i^{(g)}} =
    \frac{\partial \mathcal{L}}{\partial \hat{y}_i} \cdot
    \frac{\partial \hat{y}_i}{\partial z_i^{(g)}}
  \label{eq:gate-grad}
\end{equation}

The second factor involves the softmax derivative and the truth table
values. Standard backpropagation through the network computes the
first factor.

\begin{engineernote}
The implementation is straightforward in any automatic differentiation
framework. For each node, the forward pass computes:
\begin{enumerate}
  \item Look up the 4 truth table outputs for all 16 gate types (a
    constant $16 \times 4$ matrix).
  \item Index into this matrix using the binary inputs $a_i, b_i$ to
    get a 16-element vector of gate outputs.
  \item Compute softmax of logits to get probabilities.
  \item Dot product of probabilities and gate outputs to get
    $\hat{y}_i$.
\end{enumerate}
The backward pass is handled automatically by the framework.
\end{engineernote}


% ----------------------------------------------------------------------------
\section{Discretisation: From Soft to Hard}
\label{sec:discretisation}
% ----------------------------------------------------------------------------

After training with fully annealed temperature, each node's
probability distribution is sharply peaked at one gate type. We
``harden'' the network by selecting the most probable gate:

\begin{equation}
  g_i^* = \arg\max_g \; z_i^{(g)}
  \label{eq:hardening}
\end{equation}

The logits and softmax machinery are then discarded. What remains is
a fixed Boolean circuit: each node is a specific logic gate with two
inputs and one output.

\begin{warning}
Discretisation introduces a small accuracy drop because the
continuous weighted average is replaced by a single gate. The drop
is typically 0.1--0.5\% on MNIST~\cite{petersen2024difflogic}. If the annealing schedule is too
aggressive (temperature drops too fast), the drop can be larger
because nodes commit before finding good gate types.
\end{warning}


% ----------------------------------------------------------------------------
\section{Network Topology}
\label{sec:lgn-topology}
% ----------------------------------------------------------------------------

The LGN arranges nodes in layers. Each node in layer $l$ takes two
inputs from layer $l-1$ (or from the input). The connectivity can be:

\begin{description}[leftmargin=1em]
  \item[Random sparse:] Each node's two inputs are randomly chosen
    from the previous layer. This is the approach
    in~\cite{petersen2022difflogic}. Despite the randomness, the
    learned gate types compensate, and the network achieves
    surprisingly good accuracy.

  \item[Fully connected (impractical):] Every pair of nodes in the
    previous layer could be an input. This has quadratic cost and is
    not used.

  \item[Structured:] Inputs are chosen based on some fixed pattern
    (e.g., adjacent clause outputs). This can encode domain knowledge
    but limits flexibility.
\end{description}

For our hybrid architecture, we use random sparse connectivity.
Each layer has a fixed number of nodes, and each node's two inputs
are independently sampled (with replacement) from the previous layer's
outputs.

\begin{engineernote}
A typical LGN for MNIST~\cite{petersen2024difflogic}: 4--6 layers, each with 2000--5000 nodes,
random connectivity. The final layer's outputs are aggregated
(e.g., by popcount or a small voting circuit) to produce the
classification decision. Total gate count: 10,000--30,000.
At inference, each gate is a single table lookup (4 entries) or a
bitwise operation.
\end{engineernote}


% ----------------------------------------------------------------------------
\section{The Straight-Through Estimator Alternative}
\label{sec:ste}
% ----------------------------------------------------------------------------

An alternative to the softmax relaxation is the \concept{straight-through
estimator} (STE)~\cite{bengio2013ste}. During the forward pass, the node
commits to a hard gate type (the argmax). During the backward pass, the
gradient is computed \emph{as if} the soft relaxation were used.

This avoids the accuracy gap between soft and hard evaluation, but
introduces a bias in the gradient (since the backward pass does not
match the forward pass). In practice, both approaches work well for
LGNs; the softmax relaxation is more principled, while STE can be
faster.


% ----------------------------------------------------------------------------
\section{Summary}
\label{sec:ch8-summary}
% ----------------------------------------------------------------------------

The Logic Gate Network makes Boolean circuits learnable through three
ideas:
\begin{enumerate}
  \item \textbf{Enumerate all gates:} There are only 16 two-input
    Boolean functions. This finite set enables exhaustive
    consideration.
  \item \textbf{Softmax relaxation:} A probability distribution over
    gate types makes the output continuous and differentiable.
    Gradient descent optimises the logits.
  \item \textbf{Temperature annealing:} Gradually reducing the
    temperature transitions from exploration (all gates equally
    likely) to commitment (one gate selected).
\end{enumerate}

After training and discretisation, the LGN is a pure Boolean circuit
with zero continuous-valued operations. The next chapter describes
how to train this network on Tsetlin Machine clause outputs.

% ============================================================================
\chapter{Training the Logic Gate Network on Clause Outputs}
\label{ch:lgn-training}
% ============================================================================

This chapter describes how to train a Logic Gate Network (LGN)~\cite{petersen2022difflogic} using
the clause outputs from a pre-trained Tsetlin Machine~\cite{granmo2018tsetlin} as input. This
is Phase~2 of the hybrid architecture: the Tsetlin Machine is frozen,
and the LGN learns to combine its clause outputs into a better
classifier.


% ----------------------------------------------------------------------------
\section{The Training Setup}
\label{sec:lgn-setup}
% ----------------------------------------------------------------------------

After Phase~1 (Tsetlin Machine training), we have:
\begin{itemize}
  \item A trained Tsetlin Machine with $m$ clauses per class and
    $K = 10$ classes, giving $m \times K$ total clauses.
  \item For each training image, a binary vector of clause outputs:
    $\bvec{c} \in \{0, 1\}^{m \times K}$.
\end{itemize}

The LGN takes $\bvec{c}$ as input and produces a classification
decision. The goal is to learn a Boolean circuit that achieves higher
accuracy than the Tsetlin Machine's~\cite{granmo2018tsetlin} simple vote-and-argmax.

\begin{figure}[htbp]
\centering
\begin{tikzpicture}[
  box/.style={draw, rounded corners, minimum width=2.4cm,
    minimum height=1cm, font=\small\sffamily, align=center,
    line width=0.6pt},
  arr/.style={-{Stealth[length=4pt]}, thick},
]
  \node[box, fill=bitgreen!10] (tm) at (0, 0)
    {Tsetlin Machine\\(frozen)};
  \node[box, fill=bitorange!10] (lgn) at (5, 0)
    {Logic Gate Network\\(training)};
  \node[box, fill=bitpurple!10] (out) at (10, 0)
    {Classification\\output};

  \draw[arr] (-3, 0) -- (tm)
    node[midway, above, font=\scriptsize\sffamily] {image};
  \draw[arr] (tm) -- (lgn)
    node[midway, above, font=\scriptsize\sffamily] {clause bits};
  \draw[arr] (lgn) -- (out)
    node[midway, above, font=\scriptsize\sffamily] {gate outputs};
\end{tikzpicture}
\caption{Phase~2 training: the Tsetlin Machine is frozen and provides
  binary clause outputs. The LGN trains on these outputs using
  gradient descent with the softmax relaxation.}
\label{fig:lgn-setup}
\end{figure}


% ----------------------------------------------------------------------------
\section{Precomputing Clause Outputs}
\label{sec:precompute}
% ----------------------------------------------------------------------------

Since the Tsetlin Machine is frozen during Phase~2, we can precompute
all clause outputs for the entire training set. This is a one-time
cost that dramatically speeds up LGN training.

\begin{engineernote}
For 60,000 training images with 5,000 clauses (500 per class $\times$
10 classes), the precomputed dataset is:
\[
  60{,}000 \times 5{,}000 \text{ bits}
  = 300{,}000{,}000 \text{ bits}
  = 37.5 \text{ MB}
\]
This fits easily in memory. Each training epoch of the LGN then
reads from this cached binary matrix instead of re-evaluating the
Tsetlin Machine.
\end{engineernote}


% ----------------------------------------------------------------------------
\section{Network Architecture for MNIST}
\label{sec:lgn-arch}
% ----------------------------------------------------------------------------

A concrete LGN architecture for MNIST:

\begin{description}[leftmargin=1em]
  \item[Input layer:] 5,000 bits (clause outputs from the TM).

  \item[Hidden layers:] 4 layers, each with 4,000 nodes. Each node
    has two randomly selected inputs from the previous layer. Each
    node maintains 16 logits for gate selection.

  \item[Output layer:] 10 groups of nodes, one per digit class. Each
    group has some number of output gates. The popcount of each
    group's outputs gives the confidence score for that class.

  \item[Classification:] $\hat{y} = \arg\max_k \; \text{popcount}
    (\text{group}_k)$.
\end{description}

During training, the continuous relaxation produces real-valued
outputs in $[0, 1]$, and we use cross-entropy loss:
\begin{equation}
  \mathcal{L} = -\log \frac
    {\exp(s_y)}
    {\sum_{k=0}^{9} \exp(s_k)}
  \label{eq:lgn-loss}
\end{equation}
where $s_k$ is the (continuous) score for class $k$ and $y$ is the
true label.


% ----------------------------------------------------------------------------
\section{Training Procedure}
\label{sec:lgn-procedure}
% ----------------------------------------------------------------------------

\begin{algorithm}[htbp]
\caption{Training the Logic Gate Network (Phase~2).}
\label{alg:lgn-train}
\begin{algorithmic}[1]
  \State \textbf{Input:} Precomputed clause outputs $\{(\bvec{c}_i, y_i)\}_{i=1}^{60000}$
  \State Initialise logits $\bvec{z}$ randomly (small values near 0)
  \State Set temperature $\tau \leftarrow \tau_{\text{start}}$ (e.g., 5.0)
  \For{epoch $= 1, 2, \ldots, E$}
    \For{each mini-batch $B$ of clause output vectors}
      \State \textbf{Forward pass:}
      \For{each node $i$ in topological order}
        \State $p_i^{(g)} \leftarrow \text{softmax}(\bvec{z}_i / \tau)$
          for $g = 0, \ldots, 15$
        \State $\hat{y}_i \leftarrow \sum_g p_i^{(g)} \cdot
          \TruthTable(g)[a_i, b_i]$
      \EndFor
      \State Compute loss $\mathcal{L}$ using \cref{eq:lgn-loss}
      \Statex
      \State \textbf{Backward pass:}
      \State Compute $\frac{\partial \mathcal{L}}{\partial \bvec{z}}$
        via backpropagation
      \State Update logits: $\bvec{z} \leftarrow \bvec{z} -
        \eta \frac{\partial \mathcal{L}}{\partial \bvec{z}}$
    \EndFor
    \State Anneal temperature: $\tau \leftarrow \tau \cdot \gamma$
      \Comment{$\gamma < 1$, e.g., 0.95}
  \EndFor
  \Statex
  \State \textbf{Discretise:}
  \State For each node $i$: $g_i^* \leftarrow \arg\max_g \; z_i^{(g)}$
  \State Discard logits; fix each node to gate type $g_i^*$
\end{algorithmic}
\end{algorithm}


% ----------------------------------------------------------------------------
\section{What the LGN Learns}
\label{sec:lgn-learns}
% ----------------------------------------------------------------------------

The LGN learns nonlinear combinations of clause outputs that the
Tsetlin Machine's linear voting cannot express. Examples:

\paragraph{Clause conjunction.}
``If clause~17 (top-bar pattern) fires AND clause~42 (bottom-curve
pattern) fires, this is likely a `5'.'' The LGN can implement this
with a single AND gate connecting these two clause outputs.

\paragraph{Clause inhibition.}
``If clause~23 (vertical-stroke pattern) fires AND clause~89
(horizontal-cross pattern) does NOT fire, this is a `1' not a `4'.''
The LGN implements this with a gate of type 2 ($A \wedge \bar{B}$).

\paragraph{Feature hierarchy.}
Through multiple layers, the LGN builds a hierarchy: layer~1 combines
pairs of clause outputs, layer~2 combines those combinations, and so
on. By layer~4, each gate's output depends on up to $2^4 = 16$
original clause outputs, enabling complex decision boundaries.

\begin{keyinsight}
The LGN~\cite{petersen2022difflogic} adds \emph{depth} to the Tsetlin Machine's flat structure.
While the TM evaluates clauses independently and sums the results,
the LGN creates a deep Boolean circuit that can express arbitrary
logical relationships between clause outputs. This compositional
structure is what allows the hybrid to exceed the TM's accuracy.
\end{keyinsight}


% ----------------------------------------------------------------------------
\section{Accuracy Improvement}
\label{sec:lgn-accuracy}
% ----------------------------------------------------------------------------

Typical accuracy progression for the hybrid architecture on MNIST~\cite{lecun1998mnist}:

\begin{table}[htbp]
\centering
\caption{Accuracy at each stage of the hybrid architecture.}
\label{tab:stage-accuracy}
\begin{tabular}{@{}lcc@{}}
  \toprule
  \textbf{Stage} & \textbf{Test accuracy} & \textbf{Improvement} \\
  \midrule
  TM alone (vote + argmax) & 96.5\% & --- \\
  TM + LGN (after Phase 2) & 97.5--98.0\% & +1.0--1.5\% \\
  TM + LGN + LUT (after Phase 3) & 97.3--97.8\% & $-$0.2\% (discretisation) \\
  \bottomrule
\end{tabular}
\end{table}

The LGN consistently adds 1--1.5\% accuracy over the TM's linear
voting. The LUT compilation stage introduces a small accuracy loss
(0.1--0.3\%) from quantising continuous scoring functions, which is
addressed in Part~IV.


% ----------------------------------------------------------------------------
\section{Training Considerations}
\label{sec:lgn-considerations}
% ----------------------------------------------------------------------------

\paragraph{Learning rate.}
The logits are the only trainable parameters. A learning rate of
$\eta = 0.01$ with the Adam optimiser works well.

\paragraph{Number of epochs.}
30--100 epochs are typically sufficient. The annealing schedule should
be calibrated so that $\tau$ reaches near zero by the final epoch.

\paragraph{Random connectivity.}
The random wiring of inputs to nodes is fixed before training and
does not change. Different random seeds produce slightly different
final accuracies (variance of about 0.2--0.3\%). The LGN is
surprisingly robust to the specific connectivity pattern.

\paragraph{Gate type distribution.}
After training, the distribution of gate types is informative:
\begin{itemize}
  \item AND and OR gates dominate (together 40--50\% of nodes).
  \item XOR and XNOR are rare but important (5--10\%).
  \item Constant gates (FALSE, TRUE) appear at 10--15\%, effectively
    pruning unused connections.
  \item Projection gates ($A$ or $B$, ignoring the other input)
    appear at 10--15\%, indicating that one of the node's inputs
    is irrelevant.
\end{itemize}

\begin{engineernote}
The presence of projection and constant gates suggests that the
network is over-parameterised. Post-training pruning (removing
constant gates and collapsing projections) can reduce the gate
count by 20--30\% with no accuracy loss.
\end{engineernote}


% ----------------------------------------------------------------------------
\section{The Discrete Inference Circuit}
\label{sec:inference-circuit}
% ----------------------------------------------------------------------------

After discretisation, the LGN is a fixed Boolean circuit. Evaluating
it on a new input requires:

\begin{enumerate}
  \item Compute clause outputs from the Tsetlin Machine (bitwise AND
    operations).
  \item For each gate in topological order, look up the output from
    the gate's truth table using the two inputs.
  \item Aggregate the final layer's outputs to produce a
    classification.
\end{enumerate}

Each gate evaluation is a single operation: form a 2-bit index from
the inputs, look up the result in a 4-entry table (which fits in a
single byte). With 16,000 gates across 4 layers, the total is 16,000
table lookups---trivial for any processor.

\begin{keyinsight}
The discretised LGN has \emph{zero} floating-point operations. Every
gate is a 1-byte table lookup. Combined with the Tsetlin Machine's
bitwise AND operations, the entire pipeline from pixels to
classification uses only Boolean operations, integer addition, and
small table lookups. This is the all-Boolean inference promise of
the hybrid architecture.
\end{keyinsight}


% ----------------------------------------------------------------------------
\section{Summary}
\label{sec:ch9-summary}
% ----------------------------------------------------------------------------

The Logic Gate Network training (Phase~2) transforms the Tsetlin
Machine's linear classifier into a deep Boolean circuit:

\begin{enumerate}
  \item Freeze the Tsetlin Machine and precompute clause outputs.
  \item Train a multi-layer network of 2-input gates using softmax
    relaxation and temperature annealing.
  \item Discretise: select the most probable gate type for each node.
  \item The result is a pure Boolean circuit that improves accuracy by
    1--1.5\% over the TM alone.
\end{enumerate}

Part~IV addresses the remaining piece: how to handle any scoring
functions (like the final vote aggregation) that still involve
arithmetic beyond Boolean operations, by compiling them into lookup
tables.


% Part IV: LUT Compilation
\part{LUT Compilation --- From Functions to Tables}
\label{part:lut}

% ============================================================================
\chapter{B-Splines and Learnable Activation Functions}
\label{ch:splines}
% ============================================================================

This chapter introduces the mathematical foundation for the third
component of the hybrid architecture: learnable functions that are
smooth during training but can be discretised into lookup tables for
deployment. The key tool is the B-spline, a flexible curve defined by
a small set of control points.


% ----------------------------------------------------------------------------
\section{Why Learnable Functions?}
\label{sec:why-learnable}
% ----------------------------------------------------------------------------

In a conventional neural network, activation functions (ReLU, sigmoid,
etc.) are fixed by the architect. The \concept{Kolmogorov--Arnold
Network} (KAN)~\cite{liu2024kan} proposes a radical alternative: let
the network \emph{learn} its activation functions during training.

This is motivated by the Kolmogorov--Arnold representation theorem~\cite{kolmogorov1957representation}:

\begin{theorem}[Kolmogorov--Arnold, 1957~\cite{kolmogorov1957representation}]
\label{thm:ka}
Any continuous function $f : [0,1]^n \to \mathbb{R}$ can be written as:
\begin{equation}
  f(x_1, \ldots, x_n) = \sum_{q=0}^{2n} \Phi_q\!\left(
    \sum_{p=1}^{n} \phi_{q,p}(x_p)\right)
  \label{eq:ka-theorem}
\end{equation}
where $\phi_{q,p} : [0,1] \to \mathbb{R}$ and
$\Phi_q : \mathbb{R} \to \mathbb{R}$ are continuous functions.
\end{theorem}

The theorem says that multivariate functions can be decomposed into
compositions and sums of \emph{univariate} functions ($\phi$ and
$\Phi$). If we can learn these univariate functions well, we can
represent any continuous mapping.

\begin{analogy}
Traditional neural networks are like a factory with fixed machines
(ReLU, sigmoid) connected by adjustable conveyor belts (weights). The
KAN approach is like a factory where the machines themselves are
reprogrammable---each machine learns the optimal transformation for
its position in the assembly line.
\end{analogy}


% ----------------------------------------------------------------------------
\section{B-Splines: A Primer}
\label{sec:bsplines}
% ----------------------------------------------------------------------------

A \concept{B-spline} (basis spline) is a smooth, piecewise polynomial
curve defined by \emph{control points} and
\emph{knots}~\cite{deboor1978splines}.

\begin{definition}[B-spline curve]
Given $n+1$ control points $c_0, c_1, \ldots, c_n$ and a knot vector
$\bvec{t} = (t_0, t_1, \ldots, t_{n+k})$, the B-spline of degree
$k{-}1$ is:
\begin{equation}
  S(x) = \sum_{i=0}^{n} c_i \cdot B_{i,k}(x)
  \label{eq:bspline}
\end{equation}
where $B_{i,k}(x)$ are the B-spline basis functions, defined
recursively by the Cox--de~Boor formula.
\end{definition}

For our purposes, the important properties are:

\begin{enumerate}
  \item \textbf{Smooth:} B-splines are continuously differentiable
    (up to degree $k{-}2$). Gradients flow through them.
  \item \textbf{Local control:} Moving one control point affects
    only a nearby portion of the curve. This makes learning stable.
  \item \textbf{Compact representation:} A curve is defined by a
    small number of control points (typically 8--16), regardless of
    how many points you evaluate it at.
  \item \textbf{Evaluable:} Given an input $x$, the output $S(x)$
    can be computed in $O(k)$ operations.
\end{enumerate}

\begin{figure}[htbp]
\centering
\begin{tikzpicture}[>=Stealth]
  \begin{axis}[
    width=10cm, height=5cm,
    xlabel={Input $x$}, ylabel={Output $S(x)$},
    xmin=0, xmax=1, ymin=-0.5, ymax=1.5,
    grid=major, grid style={bitgray!20},
    every axis label/.style={font=\small\sffamily},
    tick label style={font=\small},
  ]
    % A smooth learned function
    \addplot[very thick, bitblue, smooth, domain=0:1, samples=100]
      {0.2 + 0.8*sin(deg(3.14*x)) - 0.3*sin(deg(6.28*x))};

    % Control points
    \addplot[only marks, mark=*, mark size=3pt, bitorange]
      coordinates {
        (0, 0.2) (0.125, 0.5) (0.25, 0.9)
        (0.375, 1.1) (0.5, 0.8) (0.625, 0.4)
        (0.75, 0.1) (0.875, 0.3) (1.0, 0.5)
      };

    \node[font=\scriptsize\sffamily, bitorange] at (axis cs:0.85, 1.2)
      {control points};
    \node[font=\scriptsize\sffamily, bitblue] at (axis cs:0.15, 1.2)
      {learned curve};
  \end{axis}
\end{tikzpicture}
\caption{A B-spline curve defined by 9 control points (orange dots).
  The smooth curve (blue) passes near the control points. During
  training, gradient descent adjusts the control points to minimise
  the loss.}
\label{fig:bspline}
\end{figure}


% ----------------------------------------------------------------------------
\section{KAN: Putting Splines on Edges}
\label{sec:kan}
% ----------------------------------------------------------------------------

In a traditional neural network, computation is:
\begin{equation}
  y = \sigma(w \cdot x + b)
  \label{eq:trad-neuron-recap}
\end{equation}
The \emph{weight} $w$ on the edge is a single number (linear
scaling), and the \emph{activation} $\sigma$ on the node is a fixed
function.

In a KAN, this is reversed:
\begin{equation}
  y = \sum_{i} \phi_i(x_i)
  \label{eq:kan-neuron}
\end{equation}
The \emph{edge} carries a learnable function $\phi_i$ (a B-spline),
and the \emph{node} is a simple summation.

\begin{figure}[htbp]
\centering
\begin{tikzpicture}[
  input/.style={circle, draw, minimum size=0.6cm, font=\small},
  neuron/.style={circle, draw, minimum size=0.9cm, font=\small,
    fill=bitblue!10},
  spline/.style={draw, rounded corners, fill=bitorange!10,
    minimum width=1.2cm, minimum height=0.5cm, font=\tiny\sffamily},
  >=Stealth
]
  % Traditional
  \node[font=\small\sffamily\bfseries] at (0, 3) {Traditional};
  \node[input] (t1) at (-1, 1.5) {$x_1$};
  \node[input] (t2) at (-1, 0) {$x_2$};
  \node[neuron] (tn) at (1.5, 0.75) {$\sigma(\Sigma)$};
  \draw[->, thick] (t1) -- (tn) node[midway, above, font=\tiny, sloped] {$w_1$};
  \draw[->, thick] (t2) -- (tn) node[midway, below, font=\tiny, sloped] {$w_2$};

  % KAN
  \node[font=\small\sffamily\bfseries] at (6, 3) {KAN};
  \node[input] (k1) at (5, 1.5) {$x_1$};
  \node[input] (k2) at (5, 0) {$x_2$};
  \node[neuron] (kn) at (8.5, 0.75) {$\Sigma$};
  \node[spline] (s1) at (6.7, 1.5) {$\phi_1$};
  \node[spline] (s2) at (6.7, 0) {$\phi_2$};
  \draw[->, thick] (k1) -- (s1);
  \draw[->, thick] (s1) -- (kn);
  \draw[->, thick] (k2) -- (s2);
  \draw[->, thick] (s2) -- (kn);
\end{tikzpicture}
\caption{Traditional neuron vs.\ KAN neuron. In a traditional neuron,
  edges carry scalar weights and the node applies a fixed activation.
  In a KAN, edges carry learnable functions (B-splines) and the node
  simply sums.}
\label{fig:kan-vs-trad}
\end{figure}

Each learnable function $\phi_i$ is a B-spline~\cite{deboor1978splines} with (say) 8 control
points. The control points are the trainable parameters. During
training, backpropagation adjusts these control points to minimise
the loss.


% ----------------------------------------------------------------------------
\section{Relevance to Our Architecture}
\label{sec:relevance}
% ----------------------------------------------------------------------------

Our hybrid architecture does not use a full KAN. Instead, it borrows
the key idea: \emph{learnable univariate functions that can be
compiled to lookup tables.}

The specific application is in the scoring layer. After the LGN
produces binary outputs, we need to aggregate them into class scores.
A simple approach is popcount (count the 1s for each class group),
but a learned scoring function can be more effective:

\begin{equation}
  s_k = \phi_k\!\left(\sum_{j \in \text{group}_k} g_j\right)
  \label{eq:scoring}
\end{equation}

Here, $\sum g_j$ is the popcount of the gate outputs for class $k$,
and $\phi_k$ is a learned function that maps this count to a score.
During training, $\phi_k$ is a B-spline. After training, it is
compiled into a lookup table.

\begin{keyinsight}
The computation-memory trade-off at work: instead of computing a
complex nonlinear function at inference time, we precompute it for
all possible input values and store the results in a table. The
function's complexity (how many control points, how many spline
segments) affects only the \emph{training} cost. The inference cost
is always a single table lookup, regardless of the function's
complexity.
\end{keyinsight}


% ----------------------------------------------------------------------------
\section{Training the Scoring Functions}
\label{sec:scoring-training}
% ----------------------------------------------------------------------------

The scoring functions can be trained jointly with the LGN (during
Phase~2) or in a separate Phase~2b after the LGN is discretised.
The latter approach is simpler:

\begin{enumerate}
  \item Discretise the LGN (fix all gate types).
  \item For each training image, compute the popcount vector
    $(n_0, n_1, \ldots, n_9)$ where $n_k$ is the number of output
    gates for class $k$ that fire.
  \item Train the scoring functions $\phi_0, \ldots, \phi_9$ on
    these popcount vectors using cross-entropy loss and gradient
    descent.
  \item The only trainable parameters are the B-spline control points
    (8--16 per function, 80--160 total for 10 classes).
\end{enumerate}

This is a tiny optimisation problem that converges in a few epochs.


% ----------------------------------------------------------------------------
\section{Why B-Splines and Not Polynomials?}
\label{sec:why-splines}
% ----------------------------------------------------------------------------

One might ask: why not use a simple polynomial $\phi(x) = a_0 +
a_1 x + a_2 x^2 + \cdots$?

\begin{enumerate}
  \item \textbf{Numerical stability:} High-degree polynomials
    oscillate wildly (Runge's phenomenon). B-splines are locally
    defined and stable.
  \item \textbf{Local control:} Adjusting one control point of a
    B-spline affects only a nearby region. Adjusting one polynomial
    coefficient affects the entire curve.
  \item \textbf{Discretisation quality:} A B-spline with 8 control
    points can be sampled at, say, 64 points to create a lookup table
    with very small approximation error. A polynomial of degree 8
    might require more samples to capture its oscillations.
\end{enumerate}

\begin{engineernote}
For the scoring functions in our architecture, the input is a
popcount value in the range $[0, n_{\text{gates}}]$. With $n_{\text{gates}} = 200$
per class, the lookup table has 201 entries. Each entry is an 8-bit
or 16-bit integer. The total table size is $10 \times 201 \times 2 =
4{,}020$ bytes---negligible.
\end{engineernote}


% ----------------------------------------------------------------------------
\section{Summary}
\label{sec:ch10-summary}
% ----------------------------------------------------------------------------

This chapter introduced B-splines as learnable functions that bridge
the gap between continuous training and discrete deployment:

\begin{enumerate}
  \item The Kolmogorov--Arnold theorem~\cite{kolmogorov1957representation} guarantees that compositions
    of univariate functions can represent any continuous mapping.
  \item B-splines~\cite{deboor1978splines} provide smooth, locally controllable, compactly
    parameterised univariate functions.
  \item During training, B-splines support standard backpropagation
    through their control points.
  \item After training, each B-spline can be sampled and stored as a
    lookup table---converting a continuous function into a discrete
    memory access.
\end{enumerate}

The next chapter shows how to compile these functions into lookup
tables and analyses the accuracy-efficiency trade-off of
discretisation.

% ============================================================================
\chapter{Compiling to Lookup Tables}
\label{ch:lut}
% ============================================================================

This chapter describes Phase~3 of the hybrid architecture: compiling
any remaining continuous functions into discrete lookup tables (LUTs).
After this step, the entire inference pipeline operates on Boolean
logic, integer arithmetic, and memory reads---with zero floating-point
operations.


% ----------------------------------------------------------------------------
\section{The LUT Compilation Idea}
\label{sec:lut-idea}
% ----------------------------------------------------------------------------

A lookup table replaces computation with memory: instead of evaluating
$f(x)$ at runtime, we precompute $f(x)$ for every possible input $x$
and store the results. At inference time, we simply read $f(x)$ from
memory.

\begin{equation}
  \text{LUT}[i] = f(i) \quad \text{for } i = 0, 1, \ldots, M{-}1
  \label{eq:lut-def}
\end{equation}

This is only practical when the number of possible inputs $M$ is
small. For a B-spline~\cite{deboor1978splines} scoring function whose input is a popcount
value in $[0, 200]$, $M = 201$---trivially small.

\begin{figure}[htbp]
\centering
\begin{tikzpicture}[
  box/.style={draw, rounded corners, minimum width=2cm,
    minimum height=0.8cm, font=\small\sffamily, align=center},
  arr/.style={-{Stealth[length=4pt]}, thick},
]
  % Training time
  \node[font=\sffamily\bfseries] at (3, 3) {Training time};
  \node[box, fill=bitorange!10] (input1) at (0, 2) {popcount\\value};
  \node[box, fill=bitblue!10] (spline) at (3.5, 2) {B-spline\\$\phi_k(x)$};
  \node[box, fill=bitgreen!10] (score1) at (7, 2) {class\\score};
  \draw[arr] (input1) -- (spline);
  \draw[arr] (spline) -- (score1);

  % Compilation arrow
  \draw[->, very thick, bitpurple, dashed] (3.5, 1.2) -- (3.5, 0.3)
    node[midway, right, font=\scriptsize\sffamily] {compile};

  % Inference time
  \node[font=\sffamily\bfseries] at (3, -0.5) {Inference time};
  \node[box, fill=bitorange!10] (input2) at (0, -1.5) {popcount\\value};
  \node[box, fill=bitpurple!10] (lut) at (3.5, -1.5) {LUT\\read};
  \node[box, fill=bitgreen!10] (score2) at (7, -1.5) {class\\score};
  \draw[arr] (input2) -- (lut);
  \draw[arr] (lut) -- (score2);
\end{tikzpicture}
\caption{LUT compilation replaces a B-spline evaluation (training time)
  with a simple table read (inference time). The function is computed
  once and stored; inference is a single memory access.}
\label{fig:lut-compile}
\end{figure}


% ----------------------------------------------------------------------------
\section{The Compilation Process}
\label{sec:compilation}
% ----------------------------------------------------------------------------

For each scoring function $\phi_k$:

\begin{enumerate}
  \item \textbf{Determine the input range.} The input is a popcount
    of binary gate outputs, so $x \in \{0, 1, \ldots, n_k\}$ where
    $n_k$ is the number of output gates for class $k$.

  \item \textbf{Evaluate the trained B-spline} at every integer
    input: compute $\phi_k(0), \phi_k(1), \ldots, \phi_k(n_k)$ in
    floating point.

  \item \textbf{Quantise to integers.} Convert each floating-point
    score to a fixed-point integer representation:
    \begin{equation}
      \text{LUT}_k[i] = \text{round}\!\left(
        \frac{\phi_k(i) - \phi_{\min}}{\phi_{\max} - \phi_{\min}}
        \times (2^b - 1)\right)
      \label{eq:quantise}
    \end{equation}
    where $b$ is the bit width (e.g., $b = 8$ for 8-bit integers)
    and $\phi_{\min}, \phi_{\max}$ are the range of the function.

  \item \textbf{Store the table.} The result is an array of $n_k + 1$
    integer values.
\end{enumerate}

\begin{engineernote}
With 8-bit quantisation and 201 entries per class:
\[
  \text{Total LUT size} = 10 \times 201 \times 1 \text{ byte}
  = 2{,}010 \text{ bytes} \approx 2 \text{ kB}
\]
This is negligible. Even with 16-bit entries, the total is 4~kB.
\end{engineernote}


% ----------------------------------------------------------------------------
\section{Quantisation Error}
\label{sec:quant-error}
% ----------------------------------------------------------------------------

Quantising a continuous function to $b$-bit integers introduces
\concept{quantisation error}. The maximum error per entry is:

\begin{equation}
  \epsilon_{\max} = \frac{\phi_{\max} - \phi_{\min}}{2^{b+1}}
  \label{eq:quant-error}
\end{equation}

For $b = 8$: $\epsilon_{\max} = \frac{\phi_{\max} - \phi_{\min}}{512}$.
Since the scoring function's range is typically small (e.g., $[-5, 5]$
for log-probability scores), the quantisation error is about 0.02 per
entry. This rarely changes the argmax classification.

\begin{keyinsight}
The classification decision depends on the \emph{relative ordering}
of class scores, not their absolute values. Quantisation preserves
relative ordering as long as the quantisation error is smaller than
the gap between the top two class scores. For well-trained models,
this gap is usually large (the correct class is confident), so 8-bit
quantisation introduces negligible accuracy loss.
\end{keyinsight}


% ----------------------------------------------------------------------------
\section{ULEEN: An Extreme LUT Architecture}
\label{sec:uleen}
% ----------------------------------------------------------------------------

The ULEEN architecture~\cite{susskind2024uleen} pushes the LUT idea
to its logical extreme: \emph{the entire network is lookup tables.}

In ULEEN:
\begin{enumerate}
  \item The input bits are grouped into small subsets (e.g., 8 bits
    each).
  \item Each subset indexes into a LUT with $2^8 = 256$ entries.
  \item The LUT outputs are combined (summed or concatenated) and
    fed into the next layer of LUTs.
  \item This continues for 2--3 layers.
\end{enumerate}

ULEEN achieves 96.2\% accuracy on MNIST~\cite{lecun1998mnist} with a model size of only
16.9~kB. Inference requires only memory reads and integer
additions---no multiplications, no logic gate evaluations, not even
comparisons (beyond the final argmax).

\begin{table}[htbp]
\centering
\caption{Comparison of LUT-based approaches.}
\label{tab:lut-comparison}
\begin{tabular}{@{}lccc@{}}
  \toprule
  \textbf{Approach} & \textbf{Accuracy} & \textbf{Model size}
    & \textbf{Operations} \\
  \midrule
  ULEEN~\cite{susskind2024uleen} & 96.2\% & 16.9 kB & Table reads + adds \\
  Our hybrid (projected) & 97--98\% & $\sim$1 MB & AND + gate LUT + score LUT \\
  2-layer FP network & 98.0\% & $\sim$940 kB & FP multiply-add \\
  \bottomrule
\end{tabular}
\end{table}

\begin{engineernote}
ULEEN shows that pure-LUT architectures are viable. Our hybrid
approach trades a larger model size for higher accuracy by using the
Tsetlin Machine and LGN to build a richer feature representation
before the final LUT scoring.
\end{engineernote}


% ----------------------------------------------------------------------------
\section{The Complete Inference Cost}
\label{sec:inference-cost}
% ----------------------------------------------------------------------------

After LUT compilation, the total inference cost for one MNIST image:

\begin{table}[htbp]
\centering
\caption{Operation count for the complete hybrid inference pipeline.}
\label{tab:inference-cost}
\begin{tabular}{@{}llr@{}}
  \toprule
  \textbf{Stage} & \textbf{Operation type} & \textbf{Count} \\
  \midrule
  Booleanise + expand & Comparisons, NOT & $\sim$1,600 \\
  TM clause evaluation & Bitwise AND, compare & $\sim$250,000 \\
  LGN gate evaluation & 4-entry LUT reads & $\sim$16,000 \\
  Popcount & Integer addition & $\sim$5,000 \\
  Score lookup & 201-entry LUT reads & 10 \\
  Argmax & Integer comparisons & 9 \\
  \midrule
  \textbf{Total} & & $\sim$273,000 \\
  \bottomrule
\end{tabular}
\end{table}

All operations are: bitwise logic, integer comparison, integer
addition, and memory reads from small tables. The energy cost per
operation is 0.1--5~pJ (logic and SRAM), compared to 3.7~pJ for each
floating-point multiply in a conventional network~\cite{horowitz2014energy}.

\begin{keyinsight}
The hybrid architecture achieves its energy efficiency not by having
fewer operations, but by using \emph{cheaper} operations. Each
operation costs 10--40$\times$ less energy than a floating-point
multiply-accumulate, and the model fits entirely in L1/L2 cache,
avoiding the 6,400$\times$ cost of DRAM accesses~\cite{horowitz2014energy}.
\end{keyinsight}


% ----------------------------------------------------------------------------
\section{Generalisation: LUTs for Any Bounded-Input Function}
\label{sec:generalisation}
% ----------------------------------------------------------------------------

The LUT compilation technique applies whenever:
\begin{enumerate}
  \item The input to the function has a small, finite number of
    possible values.
  \item The function is known (trained) before deployment.
  \item There is enough memory to store the table.
\end{enumerate}

In our architecture, the scoring functions have at most 201 possible
inputs (popcount values 0--200). But the same technique could be
applied to:

\begin{itemize}
  \item \textbf{Non-linear combination of two integer inputs:}
    If two popcounts $a \in [0, 50]$ and $b \in [0, 50]$ need to be
    combined, a 2D LUT of $51 \times 51 = 2601$ entries stores the
    result.
  \item \textbf{Activation functions in quantised networks~\cite{hubara2017quantized}:} A
    quantised activation with 8-bit input has $256$ possible values,
    all storable in a 256-entry table.
  \item \textbf{Distance computations:} If both inputs are bounded
    integers, Hamming distance or other metrics can be precomputed.
\end{itemize}

The cost of a LUT scales as $\prod_i |D_i|$ where $|D_i|$ is the
number of possible values for input $i$. For one input with 256
values: 256 entries. For two inputs with 256 values each:
$256^2 = 65{,}536$ entries. For three: $256^3 \approx 16.7$ million.
The exponential growth limits LUTs to functions of few inputs, which
is why we use them only for the final scoring layer.


% ----------------------------------------------------------------------------
\section{Summary}
\label{sec:ch11-summary}
% ----------------------------------------------------------------------------

LUT compilation is the final step that eliminates all floating-point
arithmetic from the inference pipeline:

\begin{enumerate}
  \item Trained B-spline scoring functions are evaluated at every
    possible integer input.
  \item The results are quantised to fixed-precision integers and
    stored in small tables.
  \item At inference time, a table read replaces the function
    evaluation.
  \item The quantisation error is negligible for well-trained models
    with 8-bit precision.
\end{enumerate}

With this step complete, we have all three components of the hybrid
architecture. Part~V assembles them into a single system and analyses
the result.


% Part V: The Complete System
\part{The Complete Hybrid Architecture}
\label{part:hybrid}

% ============================================================================
\chapter{Assembling the Hybrid Architecture}
\label{ch:assembly}
% ============================================================================

This chapter presents the complete hybrid architecture end-to-end. We
walk through the full pipeline---from a raw MNIST~\cite{lecun1998mnist} image to a digit
classification---showing exactly how data flows through each stage and
how the three training phases produce the final model.


% ----------------------------------------------------------------------------
\section{The Complete Architecture}
\label{sec:complete-arch}
% ----------------------------------------------------------------------------

\begin{figure}[htbp]
\centering
\begin{tikzpicture}[
  stage/.style={draw, rounded corners=4pt, minimum width=2.8cm,
    minimum height=1.4cm, align=center, font=\small\sffamily,
    line width=0.7pt},
  data/.style={font=\scriptsize\sffamily, bitgray},
  arr/.style={-{Stealth[length=5pt]}, thick},
]
  % Stages
  \node[stage, fill=bitgray!10] (input) at (0, 0)
    {\textbf{Input}\\$28 \times 28$ pixels};
  \node[stage, fill=bitorange!10, draw=bitorange] (bool) at (3.8, 0)
    {\textbf{Booleanise}\\threshold + expand};
  \node[stage, fill=bitgreen!10, draw=bitgreen] (tm) at (7.6, 0)
    {\textbf{Tsetlin Machine}\\5000 clauses};
  \node[stage, fill=bitorange!10, draw=bitorange] (lgn) at (11.4, 0)
    {\textbf{Logic Gate Net}\\16000 gates};
  \node[stage, fill=bitpurple!10, draw=bitpurple] (lut) at (15.2, 0)
    {\textbf{LUT Score}\\10 table reads};

  % Arrows with data descriptions
  \draw[arr] (input) -- (bool);
  \node[data, above] at (1.9, 0.9) {784 uint8};

  \draw[arr] (bool) -- (tm);
  \node[data, above] at (5.7, 0.9) {1568 bits};

  \draw[arr] (tm) -- (lgn);
  \node[data, above] at (9.5, 0.9) {5000 bits};

  \draw[arr] (lgn) -- (lut);
  \node[data, above] at (13.3, 0.9) {10 popcounts};

  % Output
  \draw[arr] (lut) -- ++(2, 0) node[right, font=\small\sffamily\bfseries]
    {$\hat{y} \in \{0, \ldots, 9\}$};
  \node[data] at (16.8, 0.9) {argmax};

  % Data types below
  \node[data] at (0, -1.2) {uint8};
  \node[data] at (3.8, -1.2) {bit};
  \node[data] at (7.6, -1.2) {bit};
  \node[data] at (11.4, -1.2) {bit $\to$ int};
  \node[data] at (15.2, -1.2) {int8/16};
\end{tikzpicture}
\caption{The complete hybrid inference pipeline. Data types shrink or
  remain compact at every stage. No floating-point values appear
  anywhere.}
\label{fig:complete-pipeline}
\end{figure}


% ----------------------------------------------------------------------------
\section{Walkthrough: Classifying a Single Image}
\label{sec:walkthrough}
% ----------------------------------------------------------------------------

Let us trace the classification of a specific MNIST image---the
digit~``7''---through the entire pipeline.

\subsection{Step 1: Booleanisation}

The raw image has pixel values in $[0, 255]$. Thresholding at
$\theta = 127$:

\begin{itemize}
  \item The horizontal top stroke (pixels at row~5, columns 8--19)
    has values $> 200$. These become 1.
  \item The diagonal downstroke has values $150$--$255$. These
    become 1.
  \item Background pixels have values $0$--$30$. These become 0.
\end{itemize}

After threshold: 784 bits. After literal expansion: 1568 bits (each
pixel contributes $x_k$ and $\bar{x}_k$).

\subsection{Step 2: Tsetlin Machine Clause Evaluation}

The 1568-bit literal vector is matched against 5000 clause bitmasks
(500 clauses per class, 10 classes).

For this ``7'' image:
\begin{itemize}
  \item Class ``7'' has $\sim$30--50 positive clauses that fire
    (matching the top bar, the diagonal, etc.) and $\sim$5--10
    negative clauses that fire.
  \item Class ``1'' might have $\sim$10--15 clauses that fire
    (the diagonal alone is somewhat ``1''-like).
  \item Classes ``0'', ``3'', ``8'' might have a few firing clauses.
  \item Other classes have very few firing clauses.
\end{itemize}

The output is a 5000-bit binary vector: 1 where a clause fired, 0
elsewhere. Typically $\sim$100--200 bits are 1 (sparse).

\subsection{Step 3: Logic Gate Network}

The 5000-bit clause vector passes through 4 layers of logic gates
(16,000 gates total). Each gate reads two bits and produces one bit.

The LGN combines clause outputs nonlinearly. For example:
\begin{itemize}
  \item Gate at layer 1: AND of clause~17 (``top horizontal bar
    present'') and clause~204 (``no closed loop at top'')
    $\Rightarrow$ 1 (both conditions met for ``7'').
  \item Gate at layer 2: OR of two layer-1 outputs, combining
    evidence from different stroke detectors.
  \item And so on, building increasingly abstract features.
\end{itemize}

The output layer has 10 groups. The group for class ``7'' has the
most 1-bits.

\subsection{Step 4: Popcount and LUT Scoring}

For each class $k$, count the 1-bits in its output group:
$n_k = \text{popcount}(\text{group}_k)$.

For our ``7'' image: $n_7 \approx 150$ (many gates fired),
$n_1 \approx 40$, other classes lower.

Each popcount indexes into the class's LUT:
$s_k = \text{LUT}_k[n_k]$.

The LUT maps the raw popcount to a calibrated score. The trained
scoring function might boost confident predictions (high popcount)
and suppress ambiguous ones.

\subsection{Step 5: Argmax}

$\hat{y} = \arg\max_k \; s_k = 7$. Correct.


% ----------------------------------------------------------------------------
\section{The Three Training Phases}
\label{sec:three-phases}
% ----------------------------------------------------------------------------

\begin{table}[htbp]
\centering
\caption{Summary of the three training phases.}
\label{tab:three-phases}
\begin{tabular}{@{}clll@{}}
  \toprule
  \textbf{Phase} & \textbf{Component} & \textbf{Method}
    & \textbf{Arithmetic} \\
  \midrule
  1 & Tsetlin Machine & Automata feedback & Integer only \\
  2 & Logic Gate Network & Gradient descent & Floating-point \\
  2b & Scoring functions & Gradient descent & Floating-point \\
  3 & LUT compilation & Evaluate + quantise & One-time FP \\
  \bottomrule
\end{tabular}
\end{table}

\paragraph{Phase 1: Tsetlin Machine~\cite{granmo2018tsetlin} (Chapters 5--7).}
Train clause automata on booleanised MNIST images. The output is
a set of clause bitmasks. Training uses only integer arithmetic
(increment/decrement automaton states, random number generation).
Typical duration: 20--50 epochs, a few minutes on a single CPU core.

\paragraph{Phase 2: Logic Gate Network~\cite{petersen2022difflogic} (Chapters 8--9).}
Freeze the TM. Precompute clause outputs for all training images.
Train the LGN using softmax relaxation and temperature annealing.
This phase \emph{does} use floating-point arithmetic (gradient
computation, softmax, etc.). Typical duration: 30--100 epochs,
10--30 minutes with a GPU or a few hours on CPU.

\paragraph{Phase 2b: Scoring functions (Chapter 10).}
Optionally train B-spline~\cite{deboor1978splines} scoring functions on the LGN's popcount
outputs. Tiny optimisation problem: 80--160 parameters, converges
in seconds.

\paragraph{Phase 3: LUT compilation (Chapter 11).}
Evaluate each scoring function at every integer input. Quantise to
fixed-point integers. Store as lookup tables. This is a one-time
offline computation taking milliseconds.

After Phase~3, the model contains:
\begin{itemize}
  \item 5000 clause bitmasks (each 1568 bits) $= 980$ kB
  \item 16,000 gate types (each 4 bits) $= 8$ kB
  \item 16,000 wiring connections (each $\sim$13 bits) $= 26$ kB
  \item 10 scoring LUTs ($\sim$201 entries $\times$ 2 bytes each)
    $= 4$ kB
  \item \textbf{Total: $\sim$1,020 kB}
\end{itemize}


% ----------------------------------------------------------------------------
\section{Data Flow Summary}
\label{sec:data-flow}
% ----------------------------------------------------------------------------

\begin{table}[htbp]
\centering
\caption{Data dimensions at each stage.}
\label{tab:data-flow}
\begin{tabular}{@{}llrl@{}}
  \toprule
  \textbf{Stage} & \textbf{Data type} & \textbf{Size}
    & \textbf{Representation} \\
  \midrule
  Raw input & uint8 & 784 B & Pixel intensities \\
  Booleanised & bits & 98 B & Bright/dark \\
  Expanded literals & bits & 196 B & Feature $+$ negation \\
  Clause outputs & bits & 625 B & Pattern detections \\
  Gate outputs & bits & 2000 B & Composed features \\
  Popcount scores & uint16 & 20 B & Per-class counts \\
  LUT scores & int16 & 20 B & Calibrated scores \\
  Prediction & uint8 & 1 B & Digit label \\
  \bottomrule
\end{tabular}
\end{table}

Note how the data remains compact throughout. The widest
representation is the gate output layer at 2000 bytes (16,000 bits),
which fits in a single L1 cache line on many architectures.


% ----------------------------------------------------------------------------
\section{End-to-End Pseudocode}
\label{sec:pseudocode}
% ----------------------------------------------------------------------------

\begin{algorithm}[htbp]
\caption{Complete hybrid inference for one MNIST image.}
\label{alg:full-inference}
\begin{algorithmic}[1]
  \Require Pixel array $\bvec{p} \in [0,255]^{784}$
  \Require Clause masks $\bvec{M} \in \{0,1\}^{5000 \times 1568}$
  \Require Gate types $\bvec{g} \in \{0,\ldots,15\}^{16000}$
  \Require Gate wiring $(\text{in}_1, \text{in}_2)$ for each gate
  \Require Scoring LUTs $\text{LUT}_k$ for $k = 0, \ldots, 9$
  \Ensure Predicted digit $\hat{y}$
  \Statex
  \State \textbf{// Step 1: Booleanise}
  \For{$k = 1$ to $784$}
    \State $x_k \leftarrow \mathbf{1}[p_k > 127]$
    \State $\bar{x}_k \leftarrow 1 - x_k$
  \EndFor
  \State Pack into bit vector $\bvec{b} \in \{0,1\}^{1568}$
  \Statex
  \State \textbf{// Step 2: Evaluate Tsetlin Machine clauses}
  \For{$j = 1$ to $5000$}
    \State $c_j \leftarrow \mathbf{1}[(\bvec{M}_j \wedge \bvec{b}) = \bvec{M}_j]$
      \Comment{Bitwise AND + compare}
  \EndFor
  \Statex
  \State \textbf{// Step 3: Evaluate Logic Gate Network}
  \State $\bvec{h}^{(0)} \leftarrow (c_1, c_2, \ldots, c_{5000})$
  \For{each layer $l = 1, 2, \ldots, L$}
    \For{each gate $i$ in layer $l$}
      \State $a \leftarrow h^{(l-1)}_{\text{in}_1(i)}$;
        $\;\; b \leftarrow h^{(l-1)}_{\text{in}_2(i)}$
      \State $h^{(l)}_i \leftarrow \TruthTable(g_i)[2a + b]$
        \Comment{4-entry LUT}
    \EndFor
  \EndFor
  \Statex
  \State \textbf{// Step 4: Popcount + LUT scoring}
  \For{$k = 0$ to $9$}
    \State $n_k \leftarrow \text{popcount}(\text{output group}_k)$
    \State $s_k \leftarrow \text{LUT}_k[n_k]$
  \EndFor
  \Statex
  \State \textbf{// Step 5: Classify}
  \State $\hat{y} \leftarrow \arg\max_k \; s_k$
\end{algorithmic}
\end{algorithm}

Every operation in this pseudocode is one of: integer comparison,
bitwise AND/OR/NOT, table lookup, integer addition, or integer
comparison. There is not a single floating-point operation.


% ----------------------------------------------------------------------------
\section{Summary}
\label{sec:ch12-summary}
% ----------------------------------------------------------------------------

The assembled hybrid architecture achieves the thesis stated in
\cref{ch:problem}: a neural network that classifies MNIST digits
using zero floating-point operations at inference time. The model
size is approximately 1~MB, fits in L2 cache, and processes each
image with roughly 273,000 cheap operations.

The next chapter analyses the architecture's properties in more
detail: accuracy, speed, energy, interpretability, and limitations.

% ============================================================================
\chapter{Analysis and Properties}
\label{ch:analysis}
% ============================================================================

This chapter analyses the hybrid architecture from multiple
angles: accuracy, computational cost, energy efficiency,
interpretability, and theoretical properties.


% ----------------------------------------------------------------------------
\section{Accuracy Analysis}
\label{sec:accuracy-analysis}
% ----------------------------------------------------------------------------

\begin{table}[htbp]
\centering
\caption{Expected accuracy at each stage and comparison with
  baselines.}
\label{tab:accuracy-analysis}
\begin{tabular}{@{}lcl@{}}
  \toprule
  \textbf{Model} & \textbf{Test accuracy} & \textbf{Notes} \\
  \midrule
  TM alone (simple threshold) & 96.5\% & 500 clauses/class \\
  TM + LGN (continuous) & 97.5--98.0\% & Before discretisation \\
  TM + LGN + LUT (final) & 97.3--97.8\% & After quantisation \\
  \midrule
  2-layer FP network & 98.0\% & 784-300-10, ReLU \\
  Logistic regression & 92.0\% & Linear baseline \\
  Random forest & 97.0\% & Non-neural baseline \\
  LeNet-5 CNN & 99.0\% & Convolutional FP \\
  \bottomrule
\end{tabular}
\end{table}

The hybrid achieves accuracy comparable to a standard 2-layer FP
network (within 0.5\%) while using fundamentally different
operations. The remaining gap comes from two sources:

\begin{enumerate}
  \item \textbf{Booleanisation loss:} Thresholding pixels discards
    grayscale information. Thermometer encoding ($L = 4$) would
    recover $\sim$0.5\% accuracy at the cost of 4$\times$ larger input.

  \item \textbf{Discretisation loss:} The LGN's softmax relaxation
    is more expressive than the hardened circuit. Temperature
    annealing minimises but does not eliminate this gap.
\end{enumerate}


% ----------------------------------------------------------------------------
\section{Computational Cost Comparison}
\label{sec:cost-comparison}
% ----------------------------------------------------------------------------

\begin{table}[htbp]
\centering
\caption{Operation counts per image for different architectures.}
\label{tab:op-counts}
\begin{tabular}{@{}lrrr@{}}
  \toprule
  \textbf{Architecture} & \textbf{FP MACs} & \textbf{Bitwise ops}
    & \textbf{LUT reads} \\
  \midrule
  2-layer FP (784-300-10) & 238,200 & 0 & 0 \\
  Our hybrid & 0 & $\sim$252,000 & $\sim$16,010 \\
  ULEEN~\cite{susskind2024uleen} & 0 & 0 & $\sim$15,000 \\
  \bottomrule
\end{tabular}
\end{table}

The hybrid has a comparable number of operations to the FP baseline,
but each operation is $10$--$40\times$ cheaper in energy.


% ----------------------------------------------------------------------------
\section{Energy Efficiency}
\label{sec:energy}
% ----------------------------------------------------------------------------

Using the energy figures from \cref{tab:energy-cost}
(Horowitz~\cite{horowitz2014energy}, 45\,nm technology):

\paragraph{Floating-point baseline.}
$238{,}200$ FP MACs $\times$ $(3.7 + 0.9)$ pJ
$\approx 1{,}096{,}000$ pJ $\approx 1.1 \; \mu$J per image.

\paragraph{Hybrid architecture.}
\begin{itemize}
  \item Bitwise ops: $252{,}000 \times 0.1$ pJ $= 25{,}200$ pJ
  \item Small LUT reads (4-entry, register): $16{,}000 \times 0.1$
    pJ $\approx 1{,}600$ pJ
  \item Score LUT reads (L1 cache): $10 \times 5.0$ pJ $= 50$ pJ
  \item Popcount + adds: $5{,}000 \times 0.9$ pJ $= 4{,}500$ pJ
\end{itemize}
Total: $\sim$31,000 pJ $= 0.031 \; \mu$J per image.

\begin{keyinsight}
The hybrid architecture is approximately 35$\times$ more energy
efficient than the floating-point baseline for the same task. This
factor comes from two sources: cheaper operations (logic vs.\ FP
multiply) and smaller model (fits in cache vs.\ requires DRAM).
\end{keyinsight}

These are theoretical estimates based on operation energy costs~\cite{horowitz2014energy,wheeldon2020pervasive}.
Actual energy savings depend on the hardware platform, memory
hierarchy, and implementation efficiency. Real-world speedups of
$10$--$30\times$ are more realistic due to overhead.


% ----------------------------------------------------------------------------
\section{Memory and Cache Behaviour}
\label{sec:cache}
% ----------------------------------------------------------------------------

\begin{table}[htbp]
\centering
\caption{Model size breakdown and cache residency.}
\label{tab:model-size}
\begin{tabular}{@{}lrl@{}}
  \toprule
  \textbf{Component} & \textbf{Size} & \textbf{Fits in} \\
  \midrule
  Clause bitmasks & 980 kB & L2 cache \\
  Gate types & 8 kB & L1 cache \\
  Gate wiring & 26 kB & L1 cache \\
  Scoring LUTs & 4 kB & L1 cache \\
  \midrule
  \textbf{Total} & \textbf{1,018 kB} & L2 cache \\
  \bottomrule
\end{tabular}
\end{table}

The clause bitmasks are the dominant component (96\% of model size).
They fit in L2 cache (typically 256 kB--8 MB on modern CPUs) but not
in L1 (typically 32--64 kB). The LGN and LUT components easily fit
in L1.

\begin{engineernote}
For embedded deployment on microcontrollers with limited memory, the
clause bitmasks can be compressed. Sparse clauses (those with few
included literals) can be stored as lists of included literal indices
rather than full bitmasks, potentially reducing the TM component to
100--200 kB. This requires a different evaluation strategy (iterate
over included literals instead of bitwise AND) but saves memory.
\end{engineernote}


% ----------------------------------------------------------------------------
\section{Interpretability}
\label{sec:interpretability}
% ----------------------------------------------------------------------------

Each component of the hybrid architecture has different
interpretability properties:

\paragraph{Tsetlin Machine~\cite{granmo2018tsetlin} (high interpretability).}
Each clause is a conjunction of pixel conditions. A human can inspect
a clause and understand exactly what pattern it detects: ``pixel (14, 10)
must be bright AND pixel (5, 20) must be dark AND \ldots''. Clause
visualisation (\cref{fig:clause-vis}) makes this concrete.

\paragraph{Logic Gate Network~\cite{petersen2022difflogic} (medium interpretability).}
Each gate is a named Boolean function (AND, OR, XOR, etc.), and its
two inputs are identifiable. In principle, one could trace the logic
from output to input. In practice, with 16,000 gates in 4 layers,
this is difficult for humans but straightforward for automated
analysis tools.

\paragraph{LUT scoring (high interpretability).}
Each scoring table is a simple function mapping popcount to score.
It can be plotted and inspected directly. The shape reveals how
confidence relates to the number of firing gates.

\begin{keyinsight}
The hybrid architecture is substantially more interpretable than a
conventional neural network. A standard network's weights are opaque
floating-point numbers; the hybrid's parameters are readable Boolean
conditions, named logic gates, and plottable scoring curves. This is
valuable for debugging, auditing, and building trust in the model.
\end{keyinsight}


% ----------------------------------------------------------------------------
\section{Theoretical Properties}
\label{sec:theory}
% ----------------------------------------------------------------------------

\paragraph{Expressiveness.}
The LGN~\cite{petersen2024difflogic} with all 16 gate types is functionally complete
(\cref{sec:completeness}): it can represent any Boolean function of
its inputs given sufficient depth and width. The Tsetlin Machine~\cite{granmo2018tsetlin}
produces clauses that can approximate any Boolean function of the
pixels (since any function can be expressed as a disjunction of
conjunctions, i.e., a sum of clauses). Together, the architecture
is a universal Boolean function approximator.

\paragraph{Convergence.}
The Tsetlin Machine has formal convergence
guarantees~\cite{zhang2021convergence}: with sufficient states $N$,
each automaton converges to the optimal action. The LGN training
inherits the convergence properties of gradient descent on
differentiable functions.

\paragraph{Robustness.}
The Boolean nature of the architecture provides some inherent
robustness: small perturbations to pixel values (below the threshold)
do not change the booleanised input. However, adversarial examples
that flip pixels across the threshold can still fool the model.


% ----------------------------------------------------------------------------
\section{Limitations}
\label{sec:limitations}
% ----------------------------------------------------------------------------

\begin{enumerate}
  \item \textbf{Accuracy ceiling.} The $\sim$97.5\% accuracy on MNIST~\cite{lecun1998mnist} is
    below what CNNs achieve (99\%+). The gap comes from the fixed
    booleanisation and the lack of spatial structure (no convolution).

  \item \textbf{Scalability.} The Tsetlin Machine's clause count
    grows with task complexity. For larger images (e.g., CIFAR-10),
    many more clauses and literals are needed, and training time
    increases quadratically.

  \item \textbf{Two-phase training.} The sequential training
    (TM first, then LGN) means the TM cannot benefit from the LGN's
    feedback. Joint training would be ideal but is difficult because
    the TM and LGN use different learning paradigms.

  \item \textbf{Random connectivity.} The LGN's random wiring is
    simple but suboptimal. Learned or structured connectivity could
    improve accuracy but complicates the architecture.

  \item \textbf{No spatial awareness.} Unlike CNNs, the flat
    architecture treats every pixel independently. It cannot exploit
    the fact that nearby pixels are correlated. The Convolutional
    Tsetlin Machine~\cite{granmo2019convtm} addresses this but adds
    complexity.
\end{enumerate}


% ----------------------------------------------------------------------------
\section{Summary}
\label{sec:ch13-summary}
% ----------------------------------------------------------------------------

The hybrid architecture achieves its design goals:
\begin{itemize}
  \item \textbf{Accuracy:} 97--98\% on MNIST, competitive with
    FP baselines.
  \item \textbf{Efficiency:} 35$\times$ theoretical energy reduction
    vs.\ FP networks.
  \item \textbf{Simplicity:} Zero FP operations at inference. All
    operations are bitwise logic, integer math, and table lookups.
  \item \textbf{Interpretability:} Clauses, gates, and scoring
    tables are all human-inspectable.
  \item \textbf{Model size:} $\sim$1 MB, fits in L2 cache.
\end{itemize}

These properties come with trade-offs: lower accuracy than deep CNNs,
limited scalability to complex tasks, and sequential training that
prevents joint optimisation. The next chapter discusses how to
implement this architecture in Rust.

% ============================================================================
\chapter{Implementation Considerations}
\label{ch:implementation}
% ============================================================================

This chapter discusses how to implement the hybrid architecture in
Rust, covering data structures, performance optimisations, and the
key design decisions for a proof-of-concept implementation.


% ----------------------------------------------------------------------------
\section{Why Rust?}
\label{sec:why-rust}
% ----------------------------------------------------------------------------

The choice of Rust is motivated by three factors:

\begin{enumerate}
  \item \textbf{Performance:} Rust compiles to native code with
    zero-cost abstractions. Bitwise operations map directly to
    hardware instructions with no interpreter or runtime overhead.

  \item \textbf{Safety:} Rust's ownership system prevents data races,
    buffer overflows, and memory leaks at compile time. For a
    research implementation with complex bit manipulation, this
    catches bugs early.

  \item \textbf{Portability:} Rust runs on every platform: x86,
    ARM, RISC-V, WebAssembly. The same code can target a server,
    a laptop, a Raspberry Pi, or a browser.
\end{enumerate}

\begin{engineernote}
Rust's ecosystem includes excellent support for SIMD intrinsics
(\texttt{std::arch}), which are essential for the bitwise parallel
operations in clause evaluation. The \texttt{packed\_simd} or
\texttt{std::simd} APIs allow writing architecture-aware bitwise
code that exploits 256-bit or 512-bit registers.
\end{engineernote}


% ----------------------------------------------------------------------------
\section{Data Representation}
\label{sec:data-repr}
% ----------------------------------------------------------------------------

\subsection{Bit Vectors}

The fundamental data structure is the \concept{packed bit vector}:
an array of 64-bit unsigned integers where each bit represents one
Boolean value.

For 1568 literals: $\lceil 1568 / 64 \rceil = 25$ \texttt{u64} words.

Key operations:
\begin{itemize}
  \item \textbf{AND:} Bitwise AND of two packed vectors (25
    instructions).
  \item \textbf{Equality check:} Compare two packed vectors word by
    word (25 comparisons with early exit).
  \item \textbf{Popcount:} Sum \texttt{count\_ones()} across all
    words.
  \item \textbf{NOT:} Bitwise NOT of each word (25 instructions).
\end{itemize}

\subsection{Clause Representation}

Each clause is stored as a single packed bit vector (the include
mask). Clause evaluation:

\begin{verbatim}
fn clause_matches(mask: &[u64; 25], input: &[u64; 25]) -> bool {
    for i in 0..25 {
        if (mask[i] & input[i]) != mask[i] {
            return false;
        }
    }
    true
}
\end{verbatim}

The early exit on the first non-matching word makes this very fast
for clauses that do not match (which is the common case due to
sparsity).

\subsection{Automaton States}

During training, each automaton stores its state as a \texttt{u8}
(for $N \leq 128$, states 1--256). The automata for one clause are
stored in a contiguous array of 1568 \texttt{u8} values.

\subsection{Gate Network}

Each gate stores:
\begin{itemize}
  \item Gate type: 4 bits (one of 16 types), packed as \texttt{u8}.
  \item Input indices: two \texttt{u16} values pointing to the
    previous layer.
\end{itemize}

Total per gate: 5 bytes. For 16,000 gates: 80 kB.


% ----------------------------------------------------------------------------
\section{Performance-Critical Paths}
\label{sec:critical-paths}
% ----------------------------------------------------------------------------

\subsection{Clause Evaluation (Inference Bottleneck)}

Clause evaluation dominates inference time. Optimisations:

\begin{enumerate}
  \item \textbf{SIMD:} Use 256-bit AVX2 or 512-bit AVX-512
    instructions to process 4 or 8 words per instruction instead
    of 1. This reduces the inner loop from 25 iterations to 4--7.

  \item \textbf{Early exit:} Most clauses do not match any given
    input. Checking the first few words (which often contain the
    most discriminative literals) allows early termination for
    non-matching clauses.

  \item \textbf{Cache layout:} Store clause masks contiguously in
    memory so that iterating over clauses has good spatial locality.
\end{enumerate}

\subsection{Automaton Updates (Training Bottleneck)}

Training is dominated by automaton state updates. Each update
requires:
\begin{enumerate}
  \item A random number (for stochastic feedback).
  \item A state increment or decrement.
  \item Clamping to the valid range.
\end{enumerate}

Optimisations:
\begin{itemize}
  \item Use a fast PRNG (e.g., xoshiro256 or PCG) that produces
    64 random bits per call, providing 64 random decisions at once.
  \item Vectorise the state updates using SIMD on the \texttt{u8}
    state arrays.
\end{itemize}


% ----------------------------------------------------------------------------
\section{Module Structure}
\label{sec:modules}
% ----------------------------------------------------------------------------

A suggested Rust project structure:

\begin{verbatim}
bitwise-digits/
  src/
    lib.rs          -- Public API
    bitvec.rs       -- Packed bit vector operations
    booleanise.rs   -- Thresholding + literal expansion
    tsetlin/
      mod.rs        -- TM public interface
      automaton.rs  -- Single automaton state machine
      clause.rs     -- Clause evaluation and storage
      feedback.rs   -- Type I and Type II feedback logic
      machine.rs    -- Full TM: clauses + voting
    lgn/
      mod.rs        -- LGN public interface
      gate.rs       -- Gate types and truth tables
      network.rs    -- Layered gate network
      train.rs      -- Softmax relaxation + annealing
    lut/
      mod.rs        -- LUT public interface
      spline.rs     -- B-spline evaluation
      compile.rs    -- Spline-to-LUT compilation
      score.rs      -- Scoring with compiled LUTs
    pipeline.rs     -- End-to-end inference
    mnist.rs        -- MNIST data loading
  benches/
    inference.rs    -- Inference benchmarks
    training.rs     -- Training benchmarks
  tests/
    integration.rs  -- End-to-end accuracy tests
\end{verbatim}


% ----------------------------------------------------------------------------
\section{Training Pipeline}
\label{sec:training-pipeline}
% ----------------------------------------------------------------------------

\begin{algorithm}[htbp]
\caption{Complete training pipeline.}
\label{alg:full-training}
\begin{algorithmic}[1]
  \State \textbf{// Phase 1: Train Tsetlin Machine}
  \State Load MNIST training images
  \State Booleanise all images (threshold + expand)
  \State Initialise TM~\cite{granmo2018tsetlin}: 500 clauses/class, $s = 5.0$, $T = 50$
  \For{epoch $= 1$ to $50$}
    \For{each image $(\bvec{x}, y)$ in random order}
      \State Evaluate clauses, compute votes
      \State Apply Type~I/II feedback
    \EndFor
    \State Evaluate test accuracy
  \EndFor
  \State Freeze TM clause masks
  \Statex
  \State \textbf{// Phase 2: Train Logic Gate Network}
  \State Precompute clause outputs for all training images
  \State Initialise LGN~\cite{petersen2022difflogic}: 4 layers, 4000 gates/layer, random wiring
  \For{epoch $= 1$ to $100$}
    \For{each mini-batch of clause output vectors}
      \State Forward pass with softmax relaxation
      \State Compute cross-entropy loss
      \State Backward pass, update logits
    \EndFor
    \State Anneal temperature
  \EndFor
  \State Discretise: fix each gate to its argmax type
  \Statex
  \State \textbf{// Phase 3: Compile LUTs}
  \State Compute popcount vectors for training set
  \State Train B-spline scoring functions (optional)
  \State Evaluate splines at all integer inputs
  \State Quantise to int8/int16 lookup tables
  \State Save complete model
\end{algorithmic}
\end{algorithm}


% ----------------------------------------------------------------------------
\section{Testing and Validation}
\label{sec:testing}
% ----------------------------------------------------------------------------

\paragraph{Unit tests.}
\begin{itemize}
  \item Bit vector operations: AND, OR, NOT, popcount on known inputs.
  \item Automaton transitions: verify reward/penalty move states
    correctly.
  \item Clause evaluation: manually constructed clauses on known
    images.
  \item Gate evaluation: all 16 gate types on all 4 input
    combinations.
\end{itemize}

\paragraph{Integration tests.}
\begin{itemize}
  \item Train a TM on a tiny subset (100 images) and verify it
    learns non-trivial clauses.
  \item Train an LGN on synthetic clause outputs and verify gate
    type convergence.
  \item End-to-end: train on full MNIST and verify test accuracy
    exceeds 95\%.
\end{itemize}

\paragraph{Benchmarks.}
\begin{itemize}
  \item Single-image inference latency.
  \item Full test-set evaluation throughput (images/second).
  \item Training time per epoch for each phase.
  \item Comparison with a reference FP implementation (e.g., a
    simple PyTorch 2-layer network).
\end{itemize}


% ----------------------------------------------------------------------------
\section{Benchmarking Against the Floating-Point Baseline}
\label{sec:benchmarking}
% ----------------------------------------------------------------------------

The FP baseline for comparison:

\begin{itemize}
  \item \textbf{Architecture:} 2-layer MLP, 784-300-10, ReLU
    activation.
  \item \textbf{Training:} SGD with learning rate 0.01, 50 epochs.
  \item \textbf{Expected accuracy:} $\sim$98\% on MNIST~\cite{lecun1998mnist}.
  \item \textbf{Implementation:} Either in Rust (using
    \texttt{ndarray}) or Python/PyTorch for reference.
\end{itemize}

The benchmark should measure:
\begin{enumerate}
  \item \textbf{Accuracy:} Test set accuracy for both models.
  \item \textbf{Inference speed:} Time per image, averaged over the
    10,000 test images.
  \item \textbf{Model size:} Total bytes of the saved model.
  \item \textbf{Training time:} Wall-clock time for complete training.
\end{enumerate}

\begin{engineernote}
For a fair speed comparison, both implementations should be optimised:
the FP baseline should use BLAS-accelerated matrix operations, and
the hybrid should use SIMD bitwise operations. Comparing unoptimised
implementations would be misleading in either direction.
\end{engineernote}


% ----------------------------------------------------------------------------
\section{Summary}
\label{sec:ch14-summary}
% ----------------------------------------------------------------------------

The implementation of the hybrid architecture in Rust is
straightforward because the core operations (bitwise AND, integer
increment, table lookup) map directly to hardware instructions.
The main engineering challenges are:

\begin{enumerate}
  \item Efficient packed bit vector operations with SIMD.
  \item Fast random number generation for stochastic feedback.
  \item Clean separation of training phases (integer TM, FP LGN,
    one-time LUT compilation).
  \item Thorough testing at every level to catch bit-manipulation bugs.
\end{enumerate}

With these components in place, the proof of concept can validate the
hybrid architecture's accuracy, speed, and energy promises on real
hardware.


% --- Back matter ---
\backmatter

% Appendices
\appendix
\part*{Appendices}
\addcontentsline{toc}{part}{Appendices}

% ============================================================================
\chapter{Notation Reference}
\label{app:notation}
% ============================================================================

\begin{table}[htbp]
\centering
\caption{General notation.}
\begin{tabular}{@{}ll@{}}
  \toprule
  \textbf{Symbol} & \textbf{Meaning} \\
  \midrule
  $x, y, z$ & Scalars (lowercase italic) \\
  $\bvec{x}, \bvec{w}$ & Vectors (bold lowercase) \\
  $\bmat{W}$ & Matrices (bold uppercase) \\
  $\{0, 1\}$ & Binary / Boolean values \\
  $\{-1, +1\}$ & Bipolar binary values \\
  $\bar{x}$ & Boolean complement: $\bar{x} = 1 - x$ \\
  $\wedge, \vee$ & Logical AND, OR \\
  $\oplus$ & Logical XOR \\
  $|S|$ & Cardinality of set $S$ \\
  $[n]$ & The set $\{1, 2, \ldots, n\}$ \\
  $\mathbf{1}[\cdot]$ & Indicator function: 1 if condition is true, 0 otherwise \\
  \bottomrule
\end{tabular}
\end{table}

\begin{table}[htbp]
\centering
\caption{MNIST and data notation.}
\begin{tabular}{@{}ll@{}}
  \toprule
  \textbf{Symbol} & \textbf{Meaning} \\
  \midrule
  $p_k$ & Raw pixel value, $p_k \in \{0, \ldots, 255\}$ \\
  $\theta$ & Booleanisation threshold (default 127) \\
  $x_k$ & Boolean pixel feature: $x_k = \mathbf{1}[p_k > \theta]$ \\
  $\bar{x}_k$ & Negated feature: $\bar{x}_k = 1 - x_k$ \\
  $\ell_k$ & A literal (either $x_k$ or $\bar{x}_k$) \\
  $n$ & Number of input features (784 pixels) \\
  $K$ & Number of classes (10 for MNIST) \\
  $y$ & True class label \\
  $\hat{y}$ & Predicted class label \\
  \bottomrule
\end{tabular}
\end{table}

\begin{table}[htbp]
\centering
\caption{Tsetlin Machine notation.}
\begin{tabular}{@{}ll@{}}
  \toprule
  \textbf{Symbol} & \textbf{Meaning} \\
  \midrule
  $N$ & Number of states per action in each automaton \\
  $s$ & Automaton state, $s \in \{1, \ldots, 2N\}$ \\
  $C_j^+, C_j^-$ & Positive/negative clause $j$ \\
  $v_k$ & Clause vote for class $k$ \\
  $T$ & Vote clipping threshold \\
  $s$ (parameter) & Specificity parameter \\
  $m$ & Number of clauses per class \\
  $\bvec{m}$ & Clause include bitmask \\
  \bottomrule
\end{tabular}
\end{table}

\begin{table}[htbp]
\centering
\caption{Logic Gate Network notation.}
\begin{tabular}{@{}ll@{}}
  \toprule
  \textbf{Symbol} & \textbf{Meaning} \\
  \midrule
  $g$ & Gate type index, $g \in \{0, \ldots, 15\}$ \\
  $\TruthTable(g)$ & Truth table of gate type $g$ (4-bit vector) \\
  $\bvec{z}_i$ & Logit vector for node $i$ ($\in \mathbb{R}^{16}$) \\
  $p_i^{(g)}$ & Probability of gate type $g$ at node $i$ \\
  $\tau$ & Softmax temperature \\
  $\hat{y}_i$ & Relaxed (continuous) output of node $i$ \\
  $g_i^*$ & Selected gate type after discretisation \\
  \bottomrule
\end{tabular}
\end{table}

\begin{table}[htbp]
\centering
\caption{LUT and scoring notation.}
\begin{tabular}{@{}ll@{}}
  \toprule
  \textbf{Symbol} & \textbf{Meaning} \\
  \midrule
  $\phi_k$ & Learned scoring function for class $k$ \\
  $c_i$ & B-spline control point $i$ \\
  $B_{i,k}(x)$ & B-spline basis function \\
  $\text{LUT}_k$ & Compiled lookup table for class $k$ \\
  $b$ & Quantisation bit width \\
  $s_k$ & Score for class $k$ \\
  $\popcount(\cdot)$ & Number of 1-bits in a binary word \\
  \bottomrule
\end{tabular}
\end{table}

% ============================================================================
\chapter{The 16 Two-Input Boolean Functions}
\label{app:gates}
% ============================================================================

This appendix provides the complete reference for all 16 two-input
Boolean functions. Each function is identified by its 4-bit truth
table code (outputs for inputs $(0,0), (0,1), (1,0), (1,1)$).

\begin{table}[htbp]
\centering
\caption{Complete table of all 16 two-input Boolean functions.}
\label{tab:all-16-full}
\begin{tabular}{@{}ccccccll@{}}
  \toprule
  \textbf{Index} & $(0,0)$ & $(0,1)$ & $(1,0)$ & $(1,1)$
    & \textbf{Code} & \textbf{Name} & \textbf{Expression} \\
  \midrule
  0  & 0 & 0 & 0 & 0 & 0000 & FALSE     & $0$ \\
  1  & 0 & 0 & 0 & 1 & 0001 & AND       & $A \wedge B$ \\
  2  & 0 & 0 & 1 & 0 & 0010 & Inhibit B & $A \wedge \bar{B}$ \\
  3  & 0 & 0 & 1 & 1 & 0011 & Project A & $A$ \\
  4  & 0 & 1 & 0 & 0 & 0100 & Inhibit A & $\bar{A} \wedge B$ \\
  5  & 0 & 1 & 0 & 1 & 0101 & Project B & $B$ \\
  6  & 0 & 1 & 1 & 0 & 0110 & XOR       & $A \oplus B$ \\
  7  & 0 & 1 & 1 & 1 & 0111 & OR        & $A \vee B$ \\
  8  & 1 & 0 & 0 & 0 & 1000 & NOR       & $\overline{A \vee B}$ \\
  9  & 1 & 0 & 0 & 1 & 1001 & XNOR      & $\overline{A \oplus B}$ \\
  10 & 1 & 0 & 1 & 0 & 1010 & NOT B     & $\bar{B}$ \\
  11 & 1 & 0 & 1 & 1 & 1011 & Imply A   & $A \vee \bar{B}$ \\
  12 & 1 & 1 & 0 & 0 & 1100 & NOT A     & $\bar{A}$ \\
  13 & 1 & 1 & 0 & 1 & 1101 & Imply B   & $\bar{A} \vee B$ \\
  14 & 1 & 1 & 1 & 0 & 1110 & NAND      & $\overline{A \wedge B}$ \\
  15 & 1 & 1 & 1 & 1 & 1111 & TRUE      & $1$ \\
  \bottomrule
\end{tabular}
\end{table}


% ----------------------------------------------------------------------------
\section{Properties and Relationships}
\label{sec:gate-properties}
% ----------------------------------------------------------------------------

\paragraph{Complementary pairs.}
Each function $g$ has a complement $\bar{g} = 15 - g$ (bitwise NOT
of the truth table):
\begin{center}
\begin{tabular}{@{}ll@{}}
  AND (1) $\leftrightarrow$ NAND (14) &
  OR (7) $\leftrightarrow$ NOR (8) \\
  XOR (6) $\leftrightarrow$ XNOR (9) &
  FALSE (0) $\leftrightarrow$ TRUE (15) \\
  Project A (3) $\leftrightarrow$ NOT A (12) &
  Project B (5) $\leftrightarrow$ NOT B (10) \\
  Inhibit B (2) $\leftrightarrow$ Imply B (13) &
  Inhibit A (4) $\leftrightarrow$ Imply A (11) \\
\end{tabular}
\end{center}

\paragraph{Symmetric functions.}
Functions where swapping $A$ and $B$ does not change the output:
AND (1), OR (7), XOR (6), XNOR (9), NAND (14), NOR (8), FALSE (0),
TRUE (15). The remaining 8 functions are asymmetric.

\paragraph{Functionally complete subsets.}
\begin{itemize}
  \item $\{\text{NAND}\}$ alone is functionally complete.
  \item $\{\text{NOR}\}$ alone is functionally complete.
  \item $\{\text{AND}, \text{NOT}\}$ is functionally complete.
  \item $\{\text{AND}, \text{OR}, \text{NOT}\}$ is the classical
    complete set.
\end{itemize}

The LGN~\cite{petersen2022difflogic} uses all 16 functions, which is more than sufficient for
completeness. This redundancy is beneficial: having all gate types
available means the network can represent any function with fewer
layers than a NAND-only circuit.


% ----------------------------------------------------------------------------
\section{Implementation as Lookup Tables}
\label{sec:gate-lut}
% ----------------------------------------------------------------------------

Each gate type is implemented as a 4-entry lookup table stored in
the low 4 bits of a single byte:

\begin{verbatim}
const GATE_LUT: [u8; 16] = [
    0b0000, // 0:  FALSE
    0b0001, // 1:  AND
    0b0010, // 2:  A AND NOT B
    0b0011, // 3:  A
    0b0100, // 4:  NOT A AND B
    0b0101, // 5:  B
    0b0110, // 6:  XOR
    0b0111, // 7:  OR
    0b1000, // 8:  NOR
    0b1001, // 9:  XNOR
    0b1010, // 10: NOT B
    0b1011, // 11: A OR NOT B
    0b1100, // 12: NOT A
    0b1101, // 13: NOT A OR B
    0b1110, // 14: NAND
    0b1111, // 15: TRUE
];

fn eval_gate(gate_type: u8, a: bool, b: bool) -> bool {
    let index = ((a as u8) << 1) | (b as u8);
    (GATE_LUT[gate_type as usize] >> index) & 1 == 1
}
\end{verbatim}

This evaluates any of the 16 gate types in constant time with two
shifts, one OR, one AND, and one comparison---no branching required.


% Bibliography
\printbibliography[heading=bibintoc]

\end{document}
